\documentclass[10pt,a4paper]{article}

% ----------------- Paquetes -------------------%
\usepackage[T1]{fontenc}
\usepackage[utf8]{inputenc}
\usepackage[a4paper,margin=2.5cm]{geometry}
\usepackage{mathptmx}             
\usepackage{changepage} 
\usepackage{setspace}
\usepackage[spanish]{babel}
\usepackage[labelfont=bf,labelsep=space]{caption}
\usepackage{amssymb}
\usepackage{amsmath}
\usepackage{ebgaramond}
\usepackage{placeins}
\usepackage{etoolbox}
\usepackage{setspace}
\usepackage{parskip} 
\usepackage{tcolorbox}
\usepackage{needspace}
\usepackage{xparse}
\usepackage{ocgx2}
\usepackage{hyperref}
\usepackage{hypcap}


%%%%%%%%%%%%%%%%%%%%%%%%%%%%%%%%%%%%%%%%%%%%%%%%%%%%%%%%%%%%%%%%
% --- Fija primero tus valores globales "buenos" ---
\setstretch{1.2}                 % Interlineado global
\setlength{\parskip}{0.7\baselineskip}
\setlength{\parindent}{0pt}
\newlength{\GlobalParskip}
\setlength{\GlobalParskip}{\parskip}

\newcounter{eqnum}[subsection]
\renewcommand{\theeqnum}{\arabic{eqnum}}
\newcounter{alphaitem}
\newcounter{numitem}

%% ------------------- Recuadros----------------------%%
\newcommand{\puntosclave}[1]{%
  \begin{tcolorbox}[
    colframe=black,
    title={Puntos clave},
    before upper={\setlength{\parindent}{0.5cm}\setlength{\parskip}{0.5\baselineskip}}
  ]
  #1
  \end{tcolorbox}
}

\newcommand{\modificaciones}[1]{
  \begin{tcolorbox}[
    colback=white,
    colframe=red,
    title={Modificaciones a realizar},
    before upper={\setlength{\parindent}{0.5cm}\setlength{\parskip}{0.5\baselineskip}}
  ]
  #1
  \end{tcolorbox}
}

%% ------------------- Listas----------------------%%
    % --- Lista alfabética ---
    \newenvironment{la}
      {\begin{list}{\alph{alphaitem})}{%
          \usecounter{alphaitem}%
          \setlength{\leftmargin}{1cm}%
          \setlength{\parskip}{3pt}% 
          \setlength{\itemsep}{0pt}% ← espacio entre ítems
          \setlength{\parsep}{0pt}%  ← párrafos dentro de un ítem
      }}
      {\end{list}}
      
    % --- Lista numérica ---
    \newenvironment{lnum}
      {\begin{list}{\arabic{numitem}.}{%
          \usecounter{numitem}%
          \setlength{\leftmargin}{1cm}%
          \setlength{\parskip}{3pt}% 
          \setlength{\itemsep}{0pt}% ← espacio entre ítems
          \setlength{\parsep}{0pt}%  ← párrafos dentro de un ítem
      }}
      {\end{list}}
      
% ------------------- Imágenes----------------------%
    \graphicspath{{Img/}}% Rutas por defecto 
    % --- Imagen adaptativa ---
    % Registros de dimensión definidos una sola vez
    \newdimen\imgcap
    \newdimen\imgmin
    \newdimen\imgmax
    \newdimen\imgremain
    \newdimen\imgh
    \newdimen\imgneed
    
    \newcommand{\imagenfit}[3]{%
      \addvspace{10pt}% separación antes
      \par\begingroup
        % Parámetros de composición (asignaciones, no \newdimen)
        \imgcap=1\baselineskip            % alto estimado de título+fuente
        \imgmin=0.15\textheight
        \imgmax=0.15\textheight
        \imgremain=\pagegoal
        \ifdim\pagegoal=\maxdimen \imgremain=0pt\fi
        \advance\imgremain by -\pagetotal
        \advance\imgremain by -\imgcap
        % Decisión de altura destino
        \ifdim\imgremain<\imgmin
          \newpage
          \imgh=0.20\textheight
        \else
          \imgh=\imgremain
          \ifdim\imgh>\imgmax \imgh=\imgmax \fi
        \fi
        % Reservar espacio para evitar cortes del bloque completo
        \imgneed=\imgh \advance\imgneed by \imgcap
        \Needspace{\imgneed}
        % Cuerpo
        \noindent\begin{minipage}{\textwidth}
          \centering
          \captionof{figure}{\textemdash\ #2}\par\vspace{0.3em}% título arriba
          \includegraphics[width=\linewidth,height=\imgh,keepaspectratio]{\detokenize{#1}}\par
          \vspace{0.3em}\small\textit{Fuente: }#3% fuente abajo
        \end{minipage}%
      \endgroup
      \par\addvspace{10pt}%
    }

%------------ Bloque de RECUADRO ESPECIAL -------------%
    \newenvironment{specialblock}[1]{%
      \begin{quote} % crea sangría a izquierda y derecha
      \small % texto más pequeño
      \noindent\textbf{#1}\\[-0.3em] % título en negrita
      \rule{\linewidth}{0.4pt}\\[0.6em]}{
      \end{quote}}
      

%-------------------------- Author Font -----------------------%
    \newcommand{\authorfont}[1]{{\fontfamily{EBGaramond-TLF}\selectfont #1}}

%-------------------------- Anexos -----------------------%
    \makeatletter
    \newcounter{anexo}
    \renewcommand{\theanexo}{\Roman{anexo}}
    \newcommand{\anexo}[1]{%
      \clearpage
      \phantomsection
      \refstepcounter{anexo}%
      \noindent\textbf{Anexo \theanexo.- }#1\par\medskip
      \addcontentsline{stc}{section}{Anexo \theanexo.- #1}%
    }
    \makeatother

    %-------------------------- Lista de figuras (personal) -----------------------%
\makeatletter
\newcommand{\MyListOfFigures}{%
  \clearpage
  \phantomsection
  {\centering\bfseries\Large Índice de figuras\par}%
  \medskip
  % Entrada en nuestro índice de contenido (stc)
  \addcontentsline{stc}{section}{Índice de figuras}%
  % Recuperación del listado (sin cabecera propia)
  \begingroup
    \parindent=0pt\parskip=2pt
    \@starttoc{lof}%
  \endgroup
}
\makeatother

%-------------------------- Bibliografía -----------------------%
\makeatletter
% Contador para las referencias
\newcounter{myref}
% Cabecera de la bibliografía, entra en el índice 'stc'
\newcommand{\MyBibliography}{%
  \clearpage
  \phantomsection
  {\centering\bfseries\Large Bibliografía y referencias\par}%
  \medskip
  \addcontentsline{stc}{section}{Bibliografía y referencias}%
  \setcounter{myref}{0}%
}
\newcommand{\MyBibItem}[4]{%
  \refstepcounter{myref}%
  \noindent[\themyref]\ #1.\ \emph{#2}. #3. #4\par
}
\makeatother

%--------- Establecer cuatro niveles de índice --------%
    \setcounter{tocdepth}{4}   
    \setcounter{secnumdepth}{4}% numera hasta \paragraph

%%%%%%%%%%%%%%%%%%%%%%%%%%%%%%%%

\makeatletter

    %--------- Interlineado Global --------%
    \edef\GlobalBaseLineStretch{\baselinestretch}
    
    %--------- Espaciado de seccinones --------%
    \renewcommand\section{\@startsection{section}
      {1}{\z@}
      {2.5ex \@plus 1ex \@minus .2ex} % espacio antes
      {1.5ex \@plus .2ex}             % espacio después
      {\normalfont\Large\bfseries}    % formato del título
    }
    \renewcommand\subsection{\@startsection{subsection}{2}{\z@}
      {3ex \@plus 0.8ex \@minus .2ex}
      {1.5ex \@plus .2ex}
      {\normalfont\large\bfseries}
    }
    \renewcommand\subsubsection{\@startsection{subsubsection}{3}{\z@}
      {3ex \@plus 0.5ex \@minus .2ex}
      {1ex \@plus .2ex}
      {\normalfont\normalsize\bfseries}
    }
    \renewcommand\paragraph{\@startsection{paragraph}{4}{\z@}%
    {-2ex \@plus -0.5ex \@minus -.2ex}% espacio antes
    {0.8ex \@plus .2ex}% espacio después
    {\normalfont\normalsize\itshape}}% ← cursiva en lugar de negrita

    %--------- Ecuaciones --------%   
    % Variables globales para almacenar títulos y notas
    \gdef\leq@current@section{}%
    \gdef\leq@current@subsection{}%
    \gdef\leq@last@printed@section{}%
    \gdef\leq@last@printed@subsection{}%
    
    % Comando para añadir nota (invisible en el cuerpo)
    \newcommand{\eqnote}[1]{%
      \addcontentsline{leq}{leqnote}{#1}%
    }
    
    % Comando para ecuaciones
    \newcommand{\eqblock}[2]{%
      % Verificar si necesitamos imprimir el título de sección
      \ifx\leq@current@section\leq@last@printed@section\else
        \addcontentsline{leq}{leqsection}{\leq@current@section}%
        \global\let\leq@last@printed@section\leq@current@section
        \gdef\leq@last@printed@subsection{}% Reset subsección al cambiar sección
      \fi
      % Verificar si necesitamos imprimir el título de subsección
      \ifx\leq@current@subsection\leq@last@printed@subsection\else
        \ifx\leq@current@subsection\@empty\else
          \addcontentsline{leq}{leqsubsection}{\leq@current@subsection}%
          \global\let\leq@last@printed@subsection\leq@current@subsection
        \fi
      \fi
      % Ecuación
      \refstepcounter{eqnum}%
      \begin{equation}
        #1\tag{E.\theeqnum}
      \end{equation}%
      \addcontentsline{leq}{leqt}{[E.\theeqnum] -- #2}%
      \addcontentsline{leq}{leqm}{#1}%
    }
    
    %%% ----- Índice de ecuaciones ----------%%%
    % Tamaño para []
    \newcommand{\leqIndexFont}{\normalsize}
    \newcommand{\leqTitleFont}{\tiny}
    \newcommand{\leqNoteFont}{\tiny}
    % Cabecera de sección 
    \newcommand\l@leqsection[2]{%
      \addvspace{25pt}% Mayor separación antes de nueva sección
      \noindent\leqIndexFont
      \makebox[\linewidth]{\rule{.38\linewidth}{0.4pt}\ \ \textbf{  \MakeUppercase{#1}}\ \ \rule{.38\linewidth}{0.4pt}}%
      \par\vspace{-10pt}
    }
    % Cabecera de subsección 
    \newcommand\l@leqsubsection[2]{%
      \par\addvspace{15pt}%
      \noindent\leqIndexFont #1
      \par\vspace{-7pt}
    }
    % Título de ecuación 
    \newcommand\l@leqt[2]{%
    \par\addvspace{10pt}%
      \par\noindent\leqTitleFont #1\par
    }
    % Miniatura de ecuación
    \newcommand\l@leqm[2]{%
      \vspace{0.3\baselineskip}%
      \par\scriptsize\noindent\hspace{1.5em}\(#1\)\par%
    }
    % Nota de ecuación 
    \newcommand\l@leqnote[2]{%
      \vspace{0.2\baselineskip}%
      \par\noindent\leqNoteFont\hspace{1.5em}\textbf{Nota.--} #1\par\vspace{0.5\baselineskip}%
    }
    % Generación del índice
    \newcommand{\mylistofequations}{%
      \clearpage
      \phantomsection
      % Título visual de la lista (ajústalo a tu gusto):
      {\centering\bfseries\Large Índice de ecuaciones\par}
      % Entrada en nuestro índice 'stc':
      \addcontentsline{stc}{section}{Índice de ecuaciones}%
      % Llamada a tu lista real (usa el nombre exacto de tu macro):
      \@starttoc{leq}%
    }
    
    %--------- Captura de títulos de sección --------%
    \let\leq@original@section\section
    \let\leq@original@subsection\subsection
    % Flag para saber si estamos en una sección asterisco
    \newif\ifLeqStarSection
    \LeqStarSectionfalse
    
    % Redefinir \section CON CONTROL DE ASTERISCO
    \renewcommand{\section}{%
      \@ifstar{\leq@section@star}{\@dblarg\leq@section@nostar}%
    }
    \def\leq@section@star#1{%
      \global\LeqStarSectiontrue
      \gdef\leq@current@section{#1}%
      \gdef\leq@current@subsection{}%
      \leq@original@section{#1}%
      \global\LeqStarSectionfalse
    }
    \def\leq@section@nostar[#1]#2{%
      \global\LeqStarSectionfalse
      \gdef\leq@current@section{#2}%
      \gdef\leq@current@subsection{}%
      \leq@original@section[#1]{#2}%
    }
    % Redefinir \subsection
    \renewcommand{\subsection}{%
      \@ifstar{\leq@subsection@star}{\@dblarg\leq@subsection@nostar}%
    }
    \def\leq@subsection@star#1{%
      \global\LeqStarSectiontrue
      \gdef\leq@current@subsection{#1}%
      \leq@original@subsection{#1}%
      \global\LeqStarSectionfalse
    }
    \def\leq@subsection@nostar[#1]#2{%
      \global\LeqStarSectionfalse
      \gdef\leq@current@subsection{#2}%
      \leq@original@subsection[#1]{#2}%
    }

    % ----- Restauración de valores Globales de estilo ----------%
    % Flags
    \newif\ifInFootnote
    \InFootnotefalse
    \pretocmd{\@footnotetext}{\InFootnotetrue}{}{}
    \apptocmd{\@footnotetext}{\InFootnotefalse}{}{}
    % Profundidad de listas (soporta anidadas)
    \newcount\MyListDepth \MyListDepth=0
    \AtBeginEnvironment{la}{\advance\MyListDepth by 1}
    \AfterEndEnvironment{la}{\advance\MyListDepth by -1}
    \AtBeginEnvironment{lnum}{\advance\MyListDepth by 1}
    \AfterEndEnvironment{lnum}{\advance\MyListDepth by -1}
    \AtBeginEnvironment{itemize}{\advance\MyListDepth by 1}
    \AfterEndEnvironment{itemize}{\advance\MyListDepth by -1}
    \AtBeginEnvironment{enumerate}{\advance\MyListDepth by 1}
    \AfterEndEnvironment{enumerate}{\advance\MyListDepth by -1}
    % Restauración diferida: hazlo al INICIO del próximo párrafo real
    \newif\if@pendingRestore
    \@pendingRestorefalse
    \newcommand{\RequestRestore}{\global\@pendingRestoretrue}
    \everypar{%
      \if@pendingRestore
        \global\@pendingRestorefalse
        \ifInFootnote\else
          \ifnum\MyListDepth=0
            \setlength{\parskip}{\GlobalParskip}%
            \setstretch{\GlobalBaseLineStretch}%
          \fi
        \fi
      \fi
    }
    % Al final de bloques:
    \AfterEndEnvironment{equation}{\RequestRestore}
    \AfterEndEnvironment{equation}{\RequestRestore}
    \AfterEndEnvironment{la}{\RequestRestore}
    \AfterEndEnvironment{lnum}{\RequestRestore}
    % Al final de macros:
    \apptocmd{\eqblock}{\RequestRestore}{}{}
    \apptocmd{\imagenfit}{\RequestRestore}{}{}
    
\makeatother

% ===================== ToC propio con 4 niveles (único bloque) =====================
\makeatletter

% --- Escritores: añadimos una línea al .stc en cada cabecera numerada ---
% Guardamos las versiones actuales para llamar al comportamiento normal
\let\MyTOC@orig@section\section
\let\MyTOC@orig@subsection\subsection
\let\MyTOC@orig@subsubsection\subsubsection
\let\MyTOC@orig@paragraph\paragraph

% SECTION
\renewcommand{\section}{%
  \@ifstar{\MyTOC@section@star}{\@dblarg\MyTOC@section@nostar}%
}
\def\MyTOC@section@star#1{%
  \MyTOC@orig@section{#1}% sin entrada al .stc
}
\def\MyTOC@section@nostar[#1]#2{%
  \MyTOC@orig@section[{#1}]{#2}% comportamiento normal (escribe al .toc)
  % Escribimos también nuestra línea al .stc (texto con \numberline)
  \addcontentsline{stc}{section}{\protect\numberline{\thesection}#2}%
}

% SUBSECTION
\renewcommand{\subsection}{%
  \@ifstar{\MyTOC@subsection@star}{\@dblarg\MyTOC@subsection@nostar}%
}
\def\MyTOC@subsection@star#1{%
  \MyTOC@orig@subsection{#1}%
}
\def\MyTOC@subsection@nostar[#1]#2{%
  \MyTOC@orig@subsection[{#1}]{#2}%
  \addcontentsline{stc}{subsection}{\protect\numberline{\thesubsection}#2}%
}

% SUBSUBSECTION
\renewcommand{\subsubsection}{%
  \@ifstar{\MyTOC@subsubsection@star}{\@dblarg\MyTOC@subsubsection@nostar}%
}
\def\MyTOC@subsubsection@star#1{%
  \MyTOC@orig@subsubsection{#1}%
}
\def\MyTOC@subsubsection@nostar[#1]#2{%
  \MyTOC@orig@subsubsection[{#1}]{#2}%
  \addcontentsline{stc}{subsubsection}{\protect\numberline{\thesubsubsection}#2}%
}

% PARAGRAPH
\renewcommand{\paragraph}{%
  \@ifstar{\MyTOC@paragraph@star}{\@dblarg\MyTOC@paragraph@nostar}%
}
\def\MyTOC@paragraph@star#1{%
  \MyTOC@orig@paragraph{#1}%
}
\def\MyTOC@paragraph@nostar[#1]#2{%
  \MyTOC@orig@paragraph[{#1}]{#2}%
  % Si no numeras los \paragraph, cambia \theparagraph por texto plano sin \numberline
  \addcontentsline{stc}{paragraph}{\protect\numberline{\theparagraph}#2}%
}

% --- Impresor del índice propio (lee el .stc) ---
\newcommand{\MyTOC}{%
  \par\bigskip
  {\centering\bfseries\Large Índice de contenido\par}%
  \medskip
  \begingroup
    \parindent=0pt\parskip=2pt
    \@starttoc{stc}%
  \endgroup
}

\makeatother
% ===================== fin ToC propio =====================


%%%%%%%%%%%%%%


% --- Datos del tema ---
\title{Teoría impositiva (VI): Imposición directa\\
\large Teoría y comparaciones internacionales}
\author{Marcelo Ortega}
\date{Agosto 2026}

\usepackage{soul}\sethlcolor{yellow}% resaltado amarillo: \hl{texto} (Panel TCEE, ⌃H)
\begin{document}
\maketitle
\newpage

% --- Página de resumen y modificaciones ---
\puntosclave{ %% Escribir puntos clave Apuntes CECO
     
    }
\vspace{1cm}

\modificaciones{
    
    }

\newpage
\MyTOC
\newpage

\mylistofequations
\clearpage

\MyListOfFigures
\clearpage


\pagenumbering{arabic}
\setcounter{page}{1}


\section{Introducción}



\subsection{Contextualización}


\subsubsection{Relevancia}




\subsubsection{Problemática}



\subsection{Estructura de la exposición}


\clearpage
\section{Imposición directa: ¿por qué gravar la generación de riqueza?}
\puntosclave{
    
}

El Estado debe basar su principal instrumento de imposición directa en la renta, y no en el patrimonio, en el consumo, en el número de puertas y ventanas de una vivienda, o en un impuesto de capitación uniforme para cada ciudadano. Esa decisión no es axiomática ni fue nunca evidente por sí misma. Por tanto, en lo que ahora nos centraremos es en dos preguntas entrelazadas: una normativa, ¿qué principio de justicia tributaria justifica elegir la renta como base, frente a la alternativa de vincular el impuesto al beneficio que cada uno recibe del Estado?; y otra positiva e histórica, ¿por qué, de hecho, los Estados modernos acabaron convergiendo en la renta como base preferente, y en qué momento y por qué causas concretas ocurrió esa convergencia?

Empezaremos por la segunda porque es la que mejor revela algo que la pregunta normativa oculta: la renta no ganó el debate teórico antes de imponerse en la práctica. Se impuso en la práctica primero, empujada por la necesidad fiscal de la guerra, y solo después la doctrina construyó el aparato de justificación que hoy damos por sentado. 

\subsection{Desarrollo histórico: de los signos externos de riqueza a la renta como base universal}
La renta se consolidó como base de la imposición directa porque fue, en cada momento histórico decisivo, el mejor compromiso disponible entre dos exigencias que rara vez coinciden: medir la capacidad económica real del contribuyente, y hacerlo con los medios administrativos y de información de los que el Estado disponía en cada época.

\subsubsection{Historia pre-industrial: Limitaciones administrativas}
Antes de que existiera la posibilidad técnica de observar y verificar la renta de un individuo, los Estados premodernos gravaban lo único que podían observar sin necesidad de contabilidad ni declaración voluntaria: signos externos y visibles de riqueza. 

En la Inglaterra de los siglos XVII y XVIII, el \textit{Land Tax} (impuesto uniforme sobre la tierra introducido a finales del XVII) fue la principal fuente de ingresos de la Corona durante más de un siglo, y convivió con figuras hoy pintorescas pero reveladoras de esta misma lógica: el impuesto sobre las ventanas, que gravaba las viviendas según el número de ventanas que tuvieran y el impuesto sobre las casas, que golpeaba sobre todo las residencias urbanas de Londres, mientras el de ventanas afectaba más a las grandes mansiones rurales. [Probable] La lógica común a todas estas figuras es puramente administrativa: gravan lo que el recaudador puede contar sin cooperación del contribuyente y sin necesidad de conocer su renta real, precisamente porque el Estado premoderno carecía de la capacidad de verificación —sin sistema bancario desarrollado, sin contabilidad mercantil generalizada, sin registro fiable de rentas del trabajo o del capital— que exigiría gravar directamente la renta.

El salto a la renta como base directa de imposición llega, y esto es el dato que más conviene retener de todo este subapartado, como instrumento de financiación bélica de emergencia, no como reforma de justicia tributaria en tiempos de paz. El primer impuesto sobre la renta de la historia moderna lo introdujo William Pitt el Joven, entonces Primer Ministro y Canciller del Exchequer británico, en su presupuesto de diciembre de 1798, entrando en vigor el 9 de enero de 1799, con la finalidad explícita de financiar "la prosecución de la guerra" contra la Francia napoleónica: un tipo del 1\% sobre rentas superiores a 60 libras, y del 10\% sobre las superiores a 200, organizado ya en varias cédulas según el origen de la renta (tierra, valores públicos, actividad comercial, rentas del trabajo, y más tarde dividendo.\footnote{%
    Se trata del mismo esquema cedular que luego reaparecerá en toda Europa continental y que ya identificamos, en el Anexo I, como el resto conceptual de la vieja teoría de la fuente—.
} El impuesto se concibió, expresamente, como medida temporal, y su trayectoria posterior lo confirma con una literalidad casi cómica —fue derogado en 1802 al firmarse la Paz de Amiens con Napoleón; reintroducido en 1803 al reanudarse las hostilidades; y derogado de nuevo en 1816, un año después de Waterloo, entre lo que las crónicas parlamentarias de la época describen como "una atronadora salva de aplausos", ordenando el Parlamento británico que todos los documentos relativos al impuesto fueran recogidos, troceados y triturados para borrar cualquier rastro de su existencia—.

\subsubsection{Historia industrial: }
El impuesto permaneció así, formalmente extinguido, durante más de un cuarto de siglo, hasta que Robert Peel lo reintrodujo en 1842. No ya para financiar una guerra, sino para cubrir un déficit presupuestario estructural que las viejas figuras (el impuesto sobre la malta, sobre las casas, sobre las ventanas) ya no podían sostener. Este momento, por tanto, marca el verdadero punto de inflexión:, pues es la primera vez que un impuesto sobre la renta sobrevive a la paz, y desde entonces no ha vuelto a desaparecer del sistema fiscal británico.

El patrón se repite, con variaciones nacionales, al otro lado del Atlántico. El Congreso estadounidense introdujo un impuesto sobre la renta durante la Guerra de Secesión (1861), derogado en 1872 una vez terminado el conflicto temporal que ya vimos en Gran Bretaña. No obstante, cuando el Congreso intentó restablecerlo en tiempos de paz, mediante la Revenue Act de 1894 (un tipo del $2\%$ sobre rentas superiores a $4.000\$$), el Tribunal Supremo lo declaró inconstitucional (Pollock v. Farmers Loan \& Trust Co., 1895): la Constitución exige que cualquier impuesto directo se reparta entre los estados en proporción a su población censada (art. I, secc. 2 y 9), y el Tribunal calificó el impuesto sobre la renta procedente de la propiedad como un impuesto directo en ese sentido constitucional, haciendo inviable en la práctica cualquier impuesto federal sobre la renta (repartir la recaudación entre estados en proporción a la población, y no a su riqueza real, habría sido absurdo y regresivo). Hicieron falta casi dos décadas y una enmienda constitucional (16ª) para neutralizar expresamente esa doctrina, otorgando al Congreso la potestad de establecer y recaudar impuestos sobre las rentas, cualquiera que sea su procedencia, sin necesidad de reparto entre los diversos Estados. Solo entonces, con el obstáculo constitucional despejado, se aprueba la Revenue Act de 1913 y nace el impuesto federal sobre la renta estadounidense en su forma moderna y permanente.

El propio patrón (guerra que crea la necesidad fiscal, capacidad administrativa que la hace verificable, y solo entonces consolidación como institución permanente) tiene, además de estos dos ejemplos concretos, un marco teórico propio dentro de la ciencia política y la historia económica. TILLY (1990) formuló la tesis belicista de la formación estatal: los Estados modernos construyeron su capacidad administrativa y fiscal fundamentalmente para financiar la guerra, y esa misma infraestructura, creada con fines bélicos, quedó después disponible para usos fiscales en tiempos de paz: \textit{the state made war, and war made the state}. BESLEY y PERSSON (2009) formalizaron después esta intuición histórica en un modelo económico riguroso: la capacidad fiscal de un Estado no es un dato dado desde el principio, sino el resultado de una inversión deliberada y costosa en infraestructura legal y administrativa, inversión que los Estados solo emprenden cuando la amenaza hace que el coste de no disponer de esa capacidad recaudatoria supere al coste de construirla. 

Es precisamente esta misma lógica la que explica por qué la masificación del impuesto sobre la renta en el siglo XX, ya como figura central, ocurre exactamente donde y cuando la teoría de TILLY y de BESLEY-PERSSON predice. Durante las dos guerras mundiales, cuando los Estados combatientes necesitaron movilizar recursos a una escala que ninguna figura anterior podía proporcionar, obligando a extender el impuesto (por primera vez) a las rentas medias y bajas del trabajo. Se estableció entonces la retención en origen, que hizo posible recaudar de millones de asalariados sin depender de su declaración voluntaria, y que en el Reino Unido llegó a coexistir con tipos marginales del $99,25\%$ durante la Segunda Guerra Mundial, manteniéndose en torno al $95\%$ hasta bien entrada la década de 1970.

En definitiva, la renta no se impuso como base de la imposición directa porque un consenso doctrinal previo determinara que era la medida más justa de capacidad económica, sino porque fue, sucesivamente en cada episodio bélico decisivo desde 1799, el instrumento que la capacidad administrativa del momento permitía verificar con mayor amplitud posible de base. 

\subsection{Argumentos a favor: ¿por qué establecer figuras de imposición directa?}
Esta historia, de contingencia y necesidad fiscal, es precisamente lo que hace tan importante preguntarse si la justificación normativa que hoy se ofrece para la renta resiste el escrutinio con independencia de cómo llegamos hasta aquí.

La doctrina clásica ofrece, desde el siglo XVIII, dos respuestas radicalmente distintas a la pregunta de qué determina cuánto debe pagar cada contribuyente.

\subsubsection{Eficiencia: Principio del beneficio (WICKSELL, 1896; LINDAHL, 1919)}
El principio del beneficio, desarrollado con rigor analítico por los economistas suecos WICKSELL (1896) y, sobre todo, LINDAHL (1919), sostiene una lógica de intercambio: el impuesto es, en esencia, el precio que cada ciudadano paga por los servicios públicos que recibe, exactamente como pagaría por cualquier bien en el mercado.

\subsubsection{Equidad: Capacidad de pago (SMITH, 1776)}
SMITH (1776) formuló sus célebres cuatro cánones de la imposición, y el primero de ellos es la formulación fundacional del principio de capacidad de pago: se paga en función de lo que se tiene, no de lo que se recibe (\textit{los súbditos de todo Estado deben contribuir al sostenimiento del gobierno, tan cerca como sea posible, en proporción a sus respectivas capacidades}). 

La razón por la que la imposición directa general se ha construido casi universalmente sobre capacidad de pago y no sobre beneficio es, en esencia, por limitaciones prácticas. El principio del beneficio exige poder medir, individuo a individuo, cuánto se beneficia cada uno de servicios como la defensa nacional, la administración de justicia o la seguridad pública, mientras que la capacidad de pago solo exige observar la renta del contribuyente, un dato mucho más tratable. 

Por eso, aunque el principio del beneficio no ha desaparecido de los sistemas fiscales comparados, ha quedado confinado a las figuras donde sí es posible esa medición individualizada: las tasas por servicios concretos y los tributos locales vinculados a servicios de proximidad claramente territorializables (educación, policía, parques). La imposición directa general, en cambio, se ha construido casi enteramente sobre capacidad de pago, precisamente porque es la única de las dos doctrinas capaz de sostener un impuesto de alcance nacional y de base amplia.

\paragraph{Eficiencia: ¿por qué debe ser progresiva? (MILL, 1848)}
Aceptado que se paga según capacidad, queda por resolver cuánta más capacidad debe traducirse en cuánto más impuesto. Es decir, el fundamento de la progresividad. 

MILL (1848), formuló el criterio del sacrificio igual: la justicia tributaria exige que el impuesto imponga a cada contribuyente el mismo sacrificio de utilidad, no la misma cantidad de dinero. De este principio derivan tres variantes técnicas, cada una con una implicación distinta sobre la tarifa que debería aplicarse: (i) el sacrificio absoluto igual exige que la pérdida total de utilidad sea idéntica para todos los contribuyentes, lo que apunta hacia cierta progresividad, pero de forma moderada; (ii) el sacrificio proporcional igual exige que la pérdida de utilidad sea la misma proporción de la utilidad total de cada uno, resultado más exigente pero compatible con una tarifa solo levemente progresiva o incluso proporcional, bajo ciertos supuestos; y, (iii) el sacrificio marginal igual, que persigue minimizar el sacrificio agregado de toda la sociedad, y conduce que el óptimo social sea la igualación de rentas después de impuestos hasta donde la recaudación necesaria lo permitiera, el fundamento teórico más extremo en favor de la progresividad.\footnote{Este resultado fue formalizado por EDGEWORTH (1897), como veremos más adelante.
}

\paragraph{Limitaciones: ¿se pueden realizar comparaciones interpersonales de utilidad? (L. ROBBINS, 1932)}
Esa crítica llega, sobre todo, de la mano de ROBBINS (1932), que dio forma definitiva a una objeción que hoy es prácticamente unánime entre los economistas: toda la teoría del sacrificio presupone que la utilidad de distintas personas puede compararse y sumarse entre sí, como si existiera una vara común capaz de medir cuánto sufre cada uno al perder un euro. ROBBINS mostró que esta comparación interpersonal de utilidad no tiene fundamento científico: es un juicio de valor y, por tanto, constituye un problema epistémico fundamental. La consecuencia práctica de esta crítica es la que explica por qué la literatura de imposición óptima sustituyó, décadas después, el aparato del sacrificio igual por un marco distinto (FBS paramétricas, con grado de aversión a la desigualdad explícito). No porque el problema que MILL y EDGEWORTH planteaban hubiera dejado de ser relevante, sino porque su solución técnica resultó, con el tiempo, insostenible en sus propios términos.

\subsection{Argumentos en contra: Estado Leviatán (BRENNAN y BUCHANAN, 1980)}
Cerramos este primer bloque conceptual con la objeción más radical de todas a las figuras de tirbutación directa. BRENNAN y BUCHANAN (1980), invierten por completo el supuesto sobre el que descansa toda la teoría de las finanzas públicas que hemos presentado hasta ahora. 

La doctrina ortodoxa asume implícitamente que el Estado es un planificador social benevolente, que elige la base y la tarifa del impuesto para maximizar el bienestar social sujeto a una restricción de recaudación necesaria. BRENNAN y BUCHANAN proponen modelar, en cambio, al Estado como un Leviatán maximizador de recaudación: una entidad que, una vez que la Constitución le concede una potestad tributaria, la explota hasta el límite de lo posible, no hasta el óptimo social.

Bajo este supuesto, se produce una inversión completa de la prescripción normativa.\footnote{
Evidentemente, los propios autores no la declaran como una descripción realista, sino más bien como un ejercicio de diseño constitucional prudente, en la tradición de MONTESQUIEU: \textit{es una experiencia eterna que todo hombre que tiene poder se ve tentado a abusar de él}.
} La teoría convencional de la imposición óptima recomienda bases amplias, comprehensivas y de difícil elusión, precisamente porque una base así distribuye la carga sobre la mayor cantidad posible de actividad económica y minimiza la distorsión marginal por cada euro recaudado. Pero si el Estado no es un planificador benevolente sino un Leviatán que va a extraer toda la recaudación que la base le permita, entonces esa misma amplitud y comprehensividad, lejos de ser una virtud, es lo que maximiza el poder recaudatorio explotable del Leviatán: una base amplia y difícil de eludir no es solo eficiente, es también la que menos límite efectivo impone al apetito fiscal del Estado. 

La recomendación constitucional que se sigue de este razonamiento es que los ciudadanos, situados en el momento hipotético de diseñar su propia constitución fiscal, podrían preferir deliberadamente restringir al Estado a bases estrechas, más fácilmente evitables y de menor rendimiento potencial, precisamente para limitar de antemano cuánto puede extraer, aunque esa elección sacrifique la eficiencia asignativa que la teoría ortodoxa persigue.

Evidentemente, la tesis Leviatán no es la ganadora del debate, pero sí conviene tenerla en cuenta como la advertencia metodológica que debe acompañar la exposición. Cada vez que evaluemos si una base imponible debería ser más amplia o más comprehensiva, conviene saber que existe esta vertiente doctrinal que invierte la recomendación clásica sobre la emplitud de la base y depende de un solo supuesto: si confiamos, o no, en que el Estado que va a administrar esa base amplia se comportará como un planificador benevolente.


\clearpage
\section{El objeto de gravamen: ¿cómo determinar la base imponible?}
\puntosclave{
    
}

Hemos visto por qué la renta se impuso como base casi por accidente bélico antes que por convicción teórica y por qué, normativamente, ni siquiera hoy existe un consenso cerrado sobre si la capacidad de pago es el criterio correcto ni sobre si la comprehensividad de la base es, sin matices, deseable. Esto nos deja en un punto incómodo. Sabemos que la renta se impuso como base de la imposición directa por necesidad bélica y capacidad administrativa, más que por convicción teórica previa, y sabemos que ni siquiera su justificación normativa está libre de una objeción de fondo. 

Lo que corresponde ahora es la pregunta que no podíamos responder sin presuponerla ya resuelta: una vez decidido que gravamos la generación de riqueza, ¿qué contamos exactamente como renta? Es decir, ¿cómo definimos la base imponible? 

La respuesta, como siempre, divide a la doctrina en dos escuelas, con razones propias y defendibles. Como veremos, la escuela ganadora del debate arrastra un problema interno tan serio que generó su propia alternativa de base completa, y que incluso resuelto todo esto queda una última pregunta, la de si "renta" y "patrimonio" son, en el fondo, la misma magnitud fragmentada en dos objetos por razones puramente administrativas. Las tres piezas —teorías de la renta, la respuesta del gasto al problema temporal de SHS, y la fragmentación renta/patrimonio— no son tres preguntas distintas: son la misma pregunta —qué contamos como el "algo" que gravamos— vista desde tres ángulos sucesivos.

\subsection{Base imponible: ¿cómo determinarla desde la óptica del ingreso?}

[Seguro] La doctrina hacendística clásica, de raíz alemana decimonónica, ofreció dos respuestas incompatibles entre sí a la pregunta de qué es renta, y conviene presentarlas con el mismo peso argumental, porque la que finalmente se impuso no lo hizo por ser la única defendible, sino por ganar un debate en el que la otra tenía —y sigue teniendo— razones propias de mérito.

\subsubsection{Teoría de la fuente: WAGNER y HERMANN ()}
La teoría de la fuente o renta-producto, asociada a WAGNER y HERMANN, sostiene que solo es renta el flujo periódico y regular que procede de una fuente productiva estable y duradera. 

Bajo esta óptica, una plusvalía esporádica por la venta de un bien, una herencia recibida, o un premio de lotería no son renta, precisamente por no constituir el fruto periódico de una fuente que perdura en el tiempo. No es una posición arbitraria. Tiene, al menos, tres razones de peso propio que conviene analizar. Primera, una razón de sostenibilidad fiscal intergeneracional. Al gravar solo el fruto y nunca el árbol, la teoría de la fuente protege deliberadamente la capacidad de esa fuente para seguir generando renta en el futuro, evitando que el impuesto erosione el propio capital productivo del que depende la recaudación de mañana (\textit{no comerse el grano de siembra}). Segunda, una razón de verificabilidad administrativa. Un flujo periódico y recurrente es mucho más fácil de observar, declarar y fiscalizar que un incremento patrimonial disperso y ocasional, sobre todo en contextos de capacidad administrativa limitada. Tercera, una razón de exigencia de liquidez. Al gravar solo lo que efectivamente fluye de forma recurrente, nunca exige al contribuyente encontrar dinero en efectivo para pagar un impuesto sobre algo que no ha convertido en flujo disponible.

\subsubsection{Teoría del incremento patrimonial: SCHANZ (1896), HAIG (1921) y SIMONS (1938)}
La teoría del incremento patrimonial neto o renta-entrada (accretion theory) invierte el criterio. 

Par esta vertiente, es renta cualquier incremento del poder económico del sujeto, con independencia de su periodicidad y de su origen. Su culminación técnica es el concepto SCHANZ-HAIG-SIMONS:

\eqblock{
\begin{aligned}
    Y=C+\Delta W
\end{aligned}
}

Se considera renta, tanto el consumo como la variación neta de patrimonio. Sus argumentos de superioridad son, sobre todo, de equidad. En primer lugar, mejora la equidad horizontal, porque dos personas con idéntico incremento de capacidad económica tributan igual con independencia de la forma en que ese incremento llegue a ellas. Y, en segundo lugar, mejora la equidad vertical, porque una base más comprehensiva capta mejor la capacidad económica real de las rentas altas, que con frecuencia obtienen una proporción mayor de su renta por vías que la teoría de la fuente dejaría sistemáticamente fuera. 

Pero la definición SHS no gana el debate sin coste. En particular, con su implementación material surgen dos problemas que hay que tener en cuenta:\footnote{%
    Podríamos mencionar un tercero, aunque no es tan relevante. Sería el referido al Estado como Leviatán. Una base tan comprehensiva como SHS en su forma pura, es, precisamente por su amplitud, la que se podría explotar con mayor eficacia. De esta forma, la misma propiedad que hace a SHS superior en términos de equidad bajo el supuesto de planificador benevolente es la que lo hace más peligroso bajo el supuesto de Leviatán, y esa tensión no se resuelve dentro de la propia teoría de la renta, depende, otra vez, de qué del comportamiento presunto.
}

\paragraph{Principio de realización: ¿cuál es el periodo relevante?}
El primer problema que enfrenta, prácticamente irresoluble, es el de la determinación del periodo impositivo. 

La razón es más sencilla de lo que parece. Imaginemos que dos agentes económicos (A y B) generan anualmente $10.000€$. De esta cantidad, consumen el $80\%$ y ahorran el $20\%$. Imaginemos que el tipo efectivo que enfrentan es de un $10\%$ y, por tanto, cada año ingresan a la hacienda Pública $1.000€$ por el impuesto directo sobre esos $10.000€$ generados durante el periodo. Pero, ¿qué sucede con los rendimientos obtenidos del ahorro? 

Supongamos que ambos invierten sus ahorros en cuentas indexadas, distintas, pero que generan la misma rentabilidad promedio del $10\%$. Al cabo de cinco años, la cuenta del agente A se ha mantenido estable, proporcionándole unos rendimientos efectivos del $10\%$. Sin embargo, la cuenta del agente B, más volátil, ha proporcionado unos rendimientos irregulares de: $(+35\%, −15\%, +10\%, +25\%, −5\%)$. Como han percibido los mismos rendimientos promedio (un $10\%$), cabría esperar que ambos agentes acabaran pagando la misma cantidad total de impuesto, puesto que su capacidad económica generada por el ahorro es idéntica. Pero no es así, y aquí está el problema.

Si Hacienda grava cada año el rendimiento positivo que genera la cuenta al tipo del $10\%$, pero no permite compensar de ninguna manera las pérdidas de los años en que el rendimiento es negativo (simplemente no se paga impuesto) el resultado es que, partiendo de los $2.000€$ ahorrados: el agente A terminaría los cinco años con $2.000€ × 1,10^5 = 3.221€$ y tributando $122,1€$ en total (sobre la ganancia total, $1.221€$); pero el agente B habría pagado un total de $156,05€$, ¡sobre una ganancia total de $997,8€$!.\footnote{
Para probarlo, simplemente hay que calcular:
\begin{lnum}
    \item Año 1: $2.000€ × 1,35 = 2.700€$. Ganancia de $700€$, impuesto: $70€$;
    \item Año 2: $2.700€ × 0,85 = 2.295€$. Pérdida de $405€$. No hay impuesto, pero tampoco hay compensación: esos $405€$ de pérdida, a efectos fiscales, simplemente desaparecen;
    \item Año 3: $2.295€ × 1,10 = 2.524,5€$. Ganancia de $229,5€$, impuesto: $22,95€$;
    \item Año 4: $2.524,5€ × 1,25 = 3.155,6€$. Ganancia de $631,1€$, impuesto: $63,1€$; y,
    \item Año 5: $3.155,6€ × 0,95 = 2.997,8€$. Pérdida de $157,8€$. Sin impuesto, sin compensación.
\end{lnum}
El segundo problema obvio y todavía más incómodo, el relativo a la ganancia, queda escondido en las cifras. Es la llamada \textbf{erosión por volatilidad}, un efecto puramente matemático de la capitalización compuesta: multiplicar por $1,35$ y después por $0,85$ no te devuelve al punto de partida multiplicado por $1,10^2$, sino a algo menor y completamente ajeno a cualquier decisión fiscal. Evidentemente, el sistema tributario no crea este primer problema; pero sí lo agrava ya que penaliza dos veces al agente cuya suerte ha sido más irregular, primero con una riqueza final menor por pura matemática financiera, y después con una factura fiscal mayor por el diseño del periodo impositivo.
}

Es decir, dos agentes con exactamente la misma rentabilidad media anual, pagan cantidades de impuesto sustancialmente distintas (casi un $28\%$ más el segundo, solo porque su rentabilidad fue irregular en lugar de estable). Y, encima, recae una carga fiscal mayor sobre el agente cuya ganancia ha sido, en conjunto, menor.\footnote{%
    Como sabemos, esto puede resolverse parcialmente permitiendo la compensación de pérdidas y ganancias, pero es solo un parche que no resuelve verdaderamente el problema.
}

Por otro lado, al problema de la irregularidad temporal se suma un segundo obstáculo. ¿Qué pasa si el activo en cuestión carece de mercado líquido y transparente? Imaginemos que, en lugar de invertir sus $2.000€$, uno de nuestros agentes hubiera destinado su ahorro a comprar un inmueble de $70.000€$ (con financiación adicional). Si quisiéramos gravar, año a año, la variación de valor de ese inmueble necesitaríamos conocer su valor al final de cada periodo impositivo, lo cual resulta absolutamente imposible.\footnote{%
    Como es evidente, un inmueble no se transacciona todos los días, de modo que cualquier cifra que le asignemos en los años intermedios es una estimación. Y las estimaciones disponibles, además, pueden no coincidir entre sí, provocando inequidad horizontal.
} El único momento en que el valor de ese inmueble deja de ser una estimación disputable es, por tanto, el de la venta. Y eso, por otro lado, genera un nuevo problema relativo a la fuente: dos sujetos que destinan su ahorro a una fuente de rendimientos diferente, \textit{pueden} acabar soportando cargas fiscales distintas.\footnote{%
    \textbf{NOTA.-} Para concluir que los dos agentes soportan una carga fiscal diferente es imprescindible que la arquitectura del impuesto sea de tarifa progresiva por tramos. Si la arquitectura es proporcional, soportarán la misma carga fiscal. 

La única diferencia en este caso es que uno \textit{diferirá} su carga a diferentes años y esto sí que le proporcionará una ganancia económica. Solo debemos tener en cuenta el tipo de descuento (siendo este homogeneo para ambos agentes) y veremos como el agente que tributa anualmente (que no inmoviliza su activo), obtendrá una ganancia real mayor que el otro.
} 

En definitiva, ambos ejemplos son dos caras de la misma moneda y pretenden ilustrar la condición necesaria para que la definición SHS sea verdaderamente equitativa en términos horizontales: la base imponible sujeta a gravamen debe ser la renta vitalicia que obtiene una persona, no su renta anual. Por tanto, el periodo impositivo ideal sería el de toda su vida, pues solo al final de una trayectoria desaparecen, por definición, tanto el problema de la irregularidad interanual como el de la valoración de activos no vendidos, porque todo habría terminado por realizarse en algún momento. 

Ahora bien, ¿qué Administración podría permitirse una figura impositiva de estas características? ¿qué Administración podría permitirse diferir el pago de impuestos a un único momento en la vida del contribuyente? Exigírselo es, simplemente, inviable. En consecuencia, la solución de compromiso que todos los sistemas fiscales comparados adoptan, no resuelve el problema; simplemente lo administra, combinando un periodo impositivo deliberadamente corto, con correctores parciales y desagregación de figuras de imposición directa.\footnote{%
    Entre muchas otras: la compensación de pérdidas entre ejercicios, que atenúa la penalización a la irregularidad; el principio de realización, que renuncia a gravar lo que no puede valorarse con certeza, aceptando a cambio diferir esa tributación hasta que la propia transacción revele la ganancia real; las exenciones sobre base ante ganancias extraordinarias: pero también la distinción entre impuestos sobre el patrimonio, impuesto sobre sociedades, sobre sucesiones y donaciones, etc.
}\textsuperscript{,}\footnote{%
    En realidad, sí existe uno que podría permitirse este tipo de gestión: aquel que presente una distribución perfectamente igualitaria de la riqueza; o, quizás un caso menos rígido, aquel cuyo pico de la distribución de riqueza se encuentre perfectamente en el medio.
}

\paragraph{Distorsiones intertemporales: ¿por qué penaliza el ahorro?}
Pero hay una segunda objeción propuesta por KALDOR (1955), mucho más profunda, que evidencia la contradiccón existente entre el diseño SHS y sus propios objetivos. Si introducimos la dimensión temporal en las decisiones de los agentes llegaremos a la conclusión de que, independientemente de cómo se diseñe el impuesto, si optamos por la definición SHS se distorsionan las decisiones intermporales de estos. 

Veamoslo con un ejemplo. Imaginemos dos agente que ganan $100€$ de salario, sujetos a un tipo del $30\%$ y un rendimiento de mercado del $10\%$.

El Agente A decide consumir siempre lo que le queda después de impuestos, $70€$. Por tanto, el tipo efectivo sobre su renta generada es exactamente $30\%$, sin sorpresas. Sin embargo, el Agente B decide ahorrar e invertir esos $70€$. Bajo un SHS comprehensivo el rendimiento de ese capital tributará al $30\%$, por tanto dejará $4,9€$ netos. De esta forma, el agente B dispondrá de $74,90€$ para consumir en el periodo siguiente. Hasta aquí, nada raro, el rendimiento es renta nueva, tributa como cualquier otra renta, sin doble cómputo de ningún euro. Ahora, preguntémonos una cosa: ¿cuál es el tipo efectivo total que estos agente han soportado, comparando lo que finalmente pueden consumir frente a lo que hubieran podido consumir si nada hubiera tributado nunca? 

Sin ningún impuesto: el Agente A hubiera consumido $100€$ en el periodo, mientras que el Agente B hubiera dispuesto de $110€$. Por tanto, ¿cuál es realmente el tipo efectivo que soporta el Agente B? Con el sistema anterior dispone de $74,9€$ frente a $110€$, por tanto, su tipo efectivo total es de $(110 − 74,9) / 110 = 31,9\%$. Si lo comparamos con el consumidor impaciente, que consume de inmediato $70€$ o $100€$, vemos que su tipo efectivo es $(100-70)/100=30\%$, exactamente el nominal, sin desviación alguna.

Es decir, el ahorrador soporta un tipo efectivo del $31,9\%$ sobre exactamente el mismo origen de riqueza que al consumidor le cuesta un $30\%$. Esa diferencia, que crece sin límite cuanto más tiempo permanezca el capital invertido, por el efecto añadido (compuesto) de gravar cada ronda sucesiva de rendimientos, aparece porque el sistema fiscal altera el precio relativo entre consumir hoy y consumir mañana. 

En ausencia de impuestos, el mercado le ofrece al ahorrador la posibilidad de cambiar 100€ de hoy por 110€ de dentro de un año. Sin embargo, con SHS comprehensivo, ese mismo intercambio le sale más caro: paga más, en términos relativos, por posponer el consumo de lo que el tipo de interés de mercado justificaría, precisamente porque el fisco se queda con una parte del propio rendimiento que compensaba la espera. Esta ha sido, con diferencia, la crítica clásica contra la definición SHS, pues resulta irresoluble incluso introduciendo divergencias respecto a su concepción más pura (como lo hace en la práctica totalidad de casos). No es un defecto corregible, sino una propiedad estructural de cualquier base que defina la renta incluyendo la variación del patrimonio. 

En definitiva, el resultado de este debate no es una victoria limpia de una teoría sobre la otra, sino una victoria parcial. La práctica totalidad de los sistemas fiscales comparados adoptó, a lo largo del siglo XX, la lógica de renta-entrada como principio general (de ahí que hablemos de \textit{impuestos sintéticos}), pero ninguno la llevó hasta sus últimas consecuencias, conservando en su diseño técnico restos identificables de la vieja teoría de la fuente y ajustes parciales con el fin de resolver las complicaciones que esta tiene.\footnote{%
    El ejemplo más elocuente es la propia distinción entre rendimientos y ganancias patrimoniales que sobrevive en casi cualquier IRPF comparado. Una distinción que, desde la lógica de SHS no debería existir, pero que persiste precisamente porque la vieja intuición de la teoría de la fuente nunca terminó de disolverse del todo.
}

\subsection{Base Imponible: ¿cómo determinarla desde la óptica del gasto?}
Ahora bien, el mismo KALDOR (1955), propuso sustituir por completo la base imponible con el fin de remediar esta distorsión de elección. En su tesis propone gravar no la renta, sino el gasto personal: 

\eqblock{
\begin{aligned}
    Y − \Delta S = C
\end{aligned}
}{}

Es decir el consumo efectivamente realizado, renta menos ahorro del periodo. Bajo esta base, el ahorro desaparece por completo del cómputo en el momento en que se genera, y solo se somete a gravamen cuando, eventualmente, se convierte en consumo.

\subsubsection{Dualización de bases: ¿cómo podría implementarse en la práctica?}
El Informe MEADE (1978), uno de los trabajos de mayor autoridad en el ámbito de la Teoría Impositiva, llevó esta idea hasta el máximo nivel de detalle jurídico-administrativo, distinguiendo dos variantes de implementación. ¿Por qué? Porque, evidentemente, existe una dificultad clara: ¿cómo implementarlo si nadie declara directamente cuánto consume? 

Antes esto, los autores del Informa proponen diferenciar dos bases:
\begin{la}
    \item \textbf{Transacciones reales}. La primera base, $R$ (cash-flow tax), grava únicamente transacciones reales. Es decir, ingresos menos toda inversión en activos reales. Bajo esta fórmula, el consumo se obtiene, por diferencia, restando de la renta todo lo que se ha invertido;\footnote{
    Esta base fue propuesta de forma independiente también por ANDREWS (1974)
    } y,
    \item \textbf{Flujos financieros}. La segunda base, $R+F$, añade además los flujos financieros. Es decir, los préstamos recibidos se suman a la base (financian consumo) y las devoluciones de préstamo se restan (reducen el consumo), lo que resulta necesario si se quiere capturar también el consumo financiado con deuda y no solo el financiado con renta corriente.
\end{la} 

En definitiva, lo que KALDOR resuelve es la neutralidad intertemporal perfecta. La decisión de consumir hoy o consumir dentro de veinte años deja de estar distorsionada por el sistema fiscal, porque el ahorro nunca tributa mientras permanece como tal. 

Sin embargo, la solución genera tres problemas nuevos. El primero es un problema de \textbf{transición}. Quien ya construyó su patrimonio bajo un sistema de renta se vería obligado a tributar una segunda vez, ahora como consumo, al gastar ese mismo patrimonio bajo el nuevo sistema. Evidentemente, podrían diseñarse reglas transitorias para evitar ese solapamiento, pero serían de considerable complejidad técnica. El segundo es el de las \textbf{grandes compras discretas}. \textcolor{red}{\textbf{(OJO -> Verificar esto con la tesis de EECKHOUT: compra de vivienda no es inversión)} Por ejemplo, la adquisición de una vivienda es, en algunos casos, una inversión y no un acto de consumo del periodo en que se paga. Esto obliga a idear mecanismos de imputación del valor de uso a lo largo de los años, el mismo problema de imputación de rentas de la vivienda habitual que se experimenta en otros contextos.} El tercero es de \textbf{movilidad internacional}. Un impuesto sobre el gasto resulta fácil de eludir colocando el ahorro en jurisdicciones que gravan la renta pero no el consumo diferido. Es decir, habría un incentivo a la competencia fiscal por bases móviles, como también se identifica en  otras figuras (ZODROW y MIESZKOWSKI).\footnote{%
    La India y Sri Lanka ensayaron, siguiendo directamente la propuesta de KALDOR (que llegó a asesorar personalmente al gobierno indio en la reforma), sendos impuestos sobre el gasto personal a comienzos de los años sesenta, y ambos lo abandonaron en pocos años por las dificultades administrativas que acabamos de describir, sobre todo la de verificar con precisión los flujos financieros bajo la base $R+F$ en economías con sistemas bancarios entonces poco desarrollados. Es el único intento histórico de implementación real de la propuesta, y su fracaso (más administrativo que teórico) es la razón por la que el gasto sigue siendo, hoy, el paradigma teóricamente más elegante para resolver el problema temporal de SHS, y a la vez el menos probado en la práctica.
}

\subsection{Fragmentación: ¿debemos fragmentar la base en diferentes figuras o no?}
\textcolor{red}{\textbf{OJO -> Hay que determinar si es conveniente o no separar este apartado. ¿Tiene lógica preguntarnos esto? O, en realidad, ¿deberíamos preguntarnos esto al final de "sujeto"?}}

La última pregunta relevante que cabe hacerse es por qué existen varios tributos directos distintos en lugar de uno solo, como sucede en prácticamente todos los sistemas comparados. 

La respuesta la tenemos ya medio construida, pues se trata del mismo patrón de fondo que explicaba por qué la teoría de la fuente seguía teniendo defensores frente a SHS, y por qué el propio SHS generaba el problema temporal que KALDOR trató de resolver. 

Fijémonos en algo que hasta ahora asumíamos implícitamente. Si $Y = C + \Delta W$ es la definición correcta de renta, entonces un impuesto sobre el patrimonio, por ejemplo, no mide una magnitud distinta de la que mide el impuesto sobre la renta. En realidad, mide la misma $\Delta W$ solo que observada como fotografía de balance en lugar de como acumulación de flujo a lo largo de un periodo. Entonces, si llevásemos la de finición SHS hasta sus últimas consecuencias, no deberíamos distinguir entre impuesto sobre la renta e impuesto sobre el patrimonio; ya que ambos son la misma medición de capacidad económica, tomada en dos unidades distintas (flujo anual frente a corte instantáneo). Que existan como dos tributos separados es ya, en sí mismo, una desviación deliberada respecto del ideal comprehensivo. 

Entonces, ¿por qué no tenemos una única figura de tributación directa? Las razones de esa desviación son, otra vez, las mismas tres que ya identificamos a propósito de la teoría de la fuente:
\begin{lnum}
    \item \textbf{Verificabilidad administrativa}. En primer lugar, se trata de una opción que facilita la gestión administrativa. Como es vidente, es más manejable un flujo anual reciente que reconstruir toda una trayectoria patrimonial acumulada;
    \item \textbf{Liquidez}. En segundo lugar, por la misma razón de liquidez que vimos antes. Gravar el patrimonio exige al contribuyente encontrar efectivo con el que pagar sobre un stock que puede no haber generado renta líquida ese año; y,
    \item \textbf{Amplitud de base: ¿Leviatán?}. Por último, de nuevo, el propio recelo frente al Leviatán. Una base que capturase de una sola vez toda la trayectoria vital de renta y patrimonio de un contribuyente sería, por su amplitud, la más peligrosa de explotar bajo el supuesto de Estado maximizador de recaudación.
\end{lnum}

Evidentemente esta fragmentación no es la única opción concebible y diferentes alternativas han sido tomadas en serio por la literatura. 

\subsubsection{Imposición acumulativa (VICKREY, 1947): ¿se trata de la solución definitiva?}
La más ambiciosa en esta dirección fue la de VICKREY (1947), quien propuso sustituir la liquidación anual del impuesto sobre la renta por una imposición acumulativa de renta vitalicia. Cada pago de impuesto realizado por el contribuyente se trataría como un depósito con interés en una cuenta abierta con la Hacienda Pública, cuyo saldo acumulado (incluidos intereses generados) se compensaría, al final, contra el impuesto que correspondería sobre la totalidad de la renta obtenida a lo largo de toda la vida del contribuyente hasta ese momento, no solo sobre la de un ejercicio artificialmente delimitado. 

La lógica de fondo es exactamente la generalización última de todo lo que hemos venido viendo. Considerando que SHS mide correctamente la capacidad económica, pero su problema fundamental es la distorsión intertemporal en el ahorro, 


solo cuando se aplica sobre el periodo relevante, y el periodo relevante para medir la capacidad económica real de una persona no es el año natural sino toda su vida, entonces la única base verdaderamente fiel a SHS no sería ni la renta anual ni el patrimonio como fotografía puntual, sino una renta vitalicia acumulada, que eliminaría de raíz no solo el problema temporal de KALDOR sino también la propia necesidad de fragmentar renta y patrimonio en tributos distintos. Ambos quedarían, sencillamente, absorbidos dentro de la misma cuenta acumulativa. 

La propuesta de VICKREY nunca se ha implementado a escala nacional en ningún país, y la literatura posterior que la ha revisado coincide en el diagnóstico (Tokyo Foundation, 2021). Aunque resolvería con elegancia la inequidad horizontal entre contribuyentes de renta vitalicia idéntica pero distribuida de forma distinta en el tiempo, exige un registro administrativo continuo de toda la trayectoria vital de renta de cada contribuyente, sujeto además a la dificultad añadida de la movilidad. ¿Qué ocurre si el contribuyente cambia de país de residencia a mitad de su trayectoria vital? En general, ha sido juzgada como una complejidad desproporcionada frente al beneficio de equidad que aportaría, insuficiente para justificar una reestructuración tan radical del sistema fiscal (BUCHANAN, 2006).



Con esta pieza final, la conclusión del Bloque 1 completo queda ya perfectamente cerrada, y conviene que la fijes con precisión porque es el resultado que vas a usar como premisa en todo lo que sigue: la fragmentación del objeto imponible —entre teoría de la fuente y renta-entrada, entre renta y gasto, entre renta periódica y patrimonio, entre tenencia y transmisión puntual— no es nunca una fragmentación arbitraria ni el resultado de un despiste del legislador, sino la manifestación repetida, en cada eje posible, de la misma tensión irreductible entre la comprehensividad que exige la justicia tributaria y la verificabilidad y liquidez que exige la administración práctica de cualquier sistema fiscal real. La renta vitalicia acumulada de Vickrey es el límite teórico de comprehensividad total; la fragmentación en múltiples figuras —renta anual, patrimonio, transmisiones, con sus propias sub-fragmentaciones internas que trabajarás en el Tema 19— es el punto de equilibrio práctico al que, en cada época y con la capacidad administrativa de que disponía, los sistemas fiscales comparados han ido convergiendo.

\clearpage
\section{El sujeto de gravamen}
\puntosclave{
    
}


Si hasta ahora hemos resuelto la diferentes vías para determinar la magnitud económica sujeta a gravamen (es decir, la base imponible), cada una con sus virtudes y dificultades, debemos ahora plantearnos una cuestión central: quién es el sujeto sobre el que se mide esa magnitud. 

La tensión de fondo es la misma en todo caso, el sujeto formal de un impuesto directo no tiene por qué coincidir con la unidad real donde efectivamente se experimenta el bienestar económico, se toman las decisiones de consumo y ahorro, o se genera el valor añadido. Esta discordancia entre unidad real y unidad conveniente se manifiesta en tres ejes que, aunque parezcan preguntas distintas a primera vista, son tres variantes de la misma tensión: si conviene sumar varias unidades reales en un solo sujeto formal (agregación, típicamente familiar), si conviene dividir una única unidad real en varios sujetos formales (desagregación, típicamente mediante la personalidad jurídica), y si el criterio para delimitar el alcance de esa unidad debe ser la pertenencia (residencia) o la mera localización física de la actividad (territorio). Empiezo por la agregación.





\subsection{Agregación de sujetos: ¿sí o no?}
La primera pregunta que nos importa será: ¿cuál debe ser la unidad cuya función única de bienestar debemos tener en cuenta a la hora de diseñar nuestro sistema fiscal? Por ejemplo, ¿es el hogar una unidad con una única función de bienestar? 

La respuesta será crucial, puesto que nos dirá si agregar tiene sentido como medición fiel de capacidad económica, o si es solo una simplificación administrativa que puede estar escondiendo desigualdades reales.

\subsubsection{Agregación horizontal: ¿por qué considerar el hogar la unidad básica de bienestar?}


\textcolor{red}{\textbf{OJO ->} Aquí hay que hacer una introducción explicando, bien y con un ejemplo, el \textit{impuesto al matrimonio}.}

\textcolor{red}{\textbf{OJO -> Además, las dos visiones que se presentan a continuación no se entienden muy bien.... veo dos problemas: (i) no se entiende muy bien la relación que tiene con el \textit{impuesto al matrimonio}, es decir, ¿se trata de justificar o no que el impuesto al matrimonio es algo \textit{soportable}?}; y, (ii) es poco pedagógico, se necesita más explicación para poder comprenderlo}

\paragraph{Argumentos a favor (SAMUELSON, 1956): Planificador familiar benevolente}



La justificación clásica para la agregación fue formula por SAMUELSON (1956) y desarrollada hasta su versión más influyente por BECKER (1981). Los autores suponen que, si el hogar se comporta como si maximizara una única función de utilidad compartida, con la renta total plenamente compartida entre ellos con independencia de quién la genere individualmente, entonces agregar no genera distorsión alguna. Ya que la unidad de bienestar real es el hogar, y medir la capacidad económica de la familia como un todo es, sencillamente, medir correctamente la unidad relevante. 

Un desarrollo complementario, formalizado por BERGSTROM (Teorema del niño malcriado, 1989), refuerza esta intuición desde otro ángulo. Incluso si algún miembro del hogar es puramente egoísta, la presencia de un cabeza de familia suficientemente altruista, que redistribuye transferencias internas según el bienestar del conjunto, basta para que incluso el miembro egoísta tenga incentivos a comportarse de forma que maximice el bienestar familiar agregado, porque sabe que cualquier ganancia que obtenga a costa de otro miembro será compensada, vía transferencia, de forma que termine perjudicándole a él mismo.

\paragraph{Argumentos en contra (CHIAPPORI, 1988, 1992): Heterogeneidad de preferencias}
La ruptura decisiva llega con CHIAPPORI (1988, 1992). El autor sustituye la función de utilidad única por el modelo colectivo, donde cada miembro del hogar conserva su propia función de utilidad individual y lo único que se asume es que el resultado observado es PARETO-eficiente, no que sea igualitario ni que refleje ningún reparto socialmente deseable.

Bajo estos supuestos, el resultado observado depende del poder de negociación relativo de cada miembro, que a su vez está determinado por variables externas al propio hogar. Por ejemplo, la renta potencial que cada cónyuge tendría fuera del matrimonio, las condiciones del mercado matrimonial de su entorno o las normas sociales predominantes sobre el reparto de responsabilidades domésticas.\footnote{
En relación al mercado matrimonial, (un modelo formalizado con antelación por McElroy y Horney, 1981, "Nash-Bargained Household Decisions", International Economic Review, y por Manser y Brown, 1980, "Marriage and Household Decision-Making: A Bargaining Analysis", International Economic Review)
}\textsuperscript{,}\footnote{%
    La evidencia empírica acumulada desde entonces (BROWNING y CHIAPPORI, 1998) tiende a rechazar sistemáticamente las restricciones que el modelo unitario impone sobre el comportamiento observado del hogar, y a confirmar las del modelo colectivo. El experimento natural más citado en toda esta literatura, y que probablemente ya conocías sin saber que pertenecía a este debate, es el cambio de política británico de los años setenta que trasladó el pago del subsidio por hijo del padre a la madre, sin alterar en absoluto la renta total del hogar, solo quién la recibía materialmente. Los estudios posteriores documentaron cambios significativos en el patrón de gasto familiar, con un desplazamiento hacia bienes asociados al cuidado infantil. Es decir, exactamente lo que el modelo colectivo predice que sí debería ocurrir si el poder de negociación de la madre, reforzado por el control de los fondos, desplaza la asignación familiar hacia sus propias preferencias; y exactamente lo que el modelo unitario predice que no debería ocurrir, ya que si el reparto interno fuera indiferente, cambiar el receptor formal del subsidio no debería alterar en nada el gasto.
}  

La implicación para el diseño fiscal es directa y nada trivial: si el modelo colectivo es el correcto, entonces la imposición familiar agregada, al tratar como irrelevante quién de los dos cónyuges genera la renta, no es una simplificación neutral, sino una decisión de diseño que puede estar reforzando estructuras de poder desiguales dentro del hogar, precisamente porque borra la señal —la titularidad fiscal individual de lo que cada uno gana— que en el modelo colectivo determina, en parte, la posición negociadora de cada miembro. La imposición individual, al preservar esa titularidad, no es solo una cuestión de neutralidad matrimonial en el sentido del trilema que ya trabajamos: es también, bajo esta lectura, un instrumento con efectos potenciales sobre el propio reparto interno de poder del hogar, no solo sobre la oferta de trabajo agregada de la pareja.

Esta conexión entre modelo de agregación y comportamiento real tiene, además, respaldo empírico específicamente fiscal y no solo doméstico. BICK y FUCHS-SCHÜNDELN (2018) estiman, con un modelo macroeconómico calibrado para 17 países de la OCDE, que la diferencia entre sistemas de imposición conjunta e individual explica una parte sustancial de la variación observada entre países en la oferta de trabajo de las mujeres casadas. El mismo mecanismo del tipo marginal elevado sobre el perceptor secundario que vimos con el trilema matrimonial.

Conviene no caer en la simplificación de que agregar es siempre malo. Existen razones defendibles a favor de la agregación que no dependen de aceptar el modelo unitario en su forma fuerte. La más sólida es la de las economías de escala en el consumo. Por ejemplo, dos personas que comparten vivienda no necesitan el doble de renta que una sola para alcanzar el mismo nivel de bienestar, lo que la literatura de escalas de equivalencia trata de capturar precisamente para medir correctamente capacidad económica por adulto equivalente, no por individuo aislado (por ejemplo, el de la OCDE). Esta es una razón de fidelidad a la capacidad económica real del hogar que no depende en absoluto de si el reparto interno es igualitario o no, y que conviene distinguir con precisión del argumento relativo al bienestar.

\subsubsection{Agregación vertical: ¿debemos agregar un \textit{genos} o no?}
\textcolor{red}{\textbf{LEER TODO y LUEGO VER SI ENTRA ESTO ->} El segundo problema guarda relación con una implicación mucho más sutil: ¿cómo podemos compatibilizar las fronteras temporales de lo intangible (riqueza/dinero) con las de lo tangibles (personas), manteniendo la equidad del sistema fiscal? Para ilustrar esta complicación recurramos al primer ejemplo. 

El agente A que ingresa $10.000€$ anuales, con sus ya mentadas preferencias ahorro-consumo, y que soporta una carga del $10\%$. Imaginemos, por simplicidad, guarda sus ahorros bajo el sofá. Tras una larga vida, fallece dejando $100.000€$ y nombra como heredero único a su nieto, Carlos. ¿Cómo debemos tratar esta transmisión de riqueza? 

Existen dos posibles vías, ambas incompatibles con algún principio impositivo.

\paragraph{Obligación genealógica: ¿podemos permitirnos la imposición nula?}
\textbf{Obligación familiar: imposición nula}. La primera resulta obvia. Nuestro agente A ya ha tributado por la generación de esa riqueza, por lo que sus herederos deberían poder recibirla sin ningún coste fiscal. Es decir, el impuesto se construye como una obligación \textit{familiar}, sobre todo la genealogía de ésta. Ahora bien, esto tiene dos implicaciones. La primera de ellas es que debemos considerar esa riqueza como ajena al circuito fiscal. Ya no existe a efectos fiscales. Evidentemente esto no supondría un peligro a corto plazo, pero a largo plazo puede acabar provocando que grandes cantidades de recursos queden absolutamente al margen del fisco. Por otro lado, imaginemos que Carlos tiene un amigo de la infancia, Juan, y ambos trabajan en la misma empresa percibiendo un salario de $10.000€$. Como es obvio, el año del fallecimiento de su abuelo, la capacidades económica de Carlos aumentará notablemente. Sin embargo, su carga fiscal será exactamente la misma que la de Juan, incumplimiento de facto la comprensividad fuerte de la definición (independencia de la fuente de ingresos) y, en consecuencia, la equidad horizontal del sistema (capacidad de pago).

\paragraph{Obligación personal: ¿distorsiona las decisiones individuales de los agentes?}
Por otro lado, tendríamos el extremo opuesto, aparentamente más justo. Si Carlos recibe en herencia los $100.000€$ de su abuelo, debería tributar por ellos en virtud del incremento de riqueza que esto supone. Es deir, el impuesto se construye como una obligación personal, sobre la vida de una persona concreta. Sin embargo, esta doble imposición (ya de por si discutible), tiene unas consecuencias claras sobre las decisiones de su propio abuelo: distorsiona, podemos suponer, sus decisiones de consumo-ahorro. Por tanto, resulta acabamos con una contradicción sobre los propios objetivos de la definición SHS: neutralidad del sistema, no debe alterar las decisiones económicas del contribuyente entre formas de obtener o materializar la renta.


El motivo de fondo, y la razón por la no es un defecto de diseño corregible, es que la riqueza económica no tiene fronteras naturales que coincidan con ningún periodo de observación finito. Cualquier frontera temporal que elijamos para liquidar cuentas va a conducir al problema planteado. Y en cada uno de esos cortes nos enfrentaremos al mismo dilema binario que acabamos de ver: o cuentas dos veces lo que cruza la frontera, o dejas que se escape sin contarlo nunca. Alargar el periodo de observación no elimina el problema: simplemente lo traslada a una frontera más lejana en el tiempo —del año a la vida, de la vida a la sucesión hereditaria—, sin que exista ningún periodo de observación finito, por generoso que sea, que no termine, tarde o temprano, topándose con su propia frontera y reproduciendo el mismo dilema. La única forma de eliminarlo por completo sería un periodo de observación infinito, que siguiera a la riqueza a través de todas las generaciones futuras sin cierre alguno —una construcción que ningún Estado, evidentemente, puede administrar—. Es, en el sentido más literal posible, un problema sin solución técnica: solo admite gestión parcial y compromiso, nunca resolución definitiva.}

\subsection{Desagregación del sujeto: ¿sí o no?}
Como sabemos, solo las personas físicas experimentan bienestar económico. Solo ellas consumen, ahorran, sufren la pérdida de una inversión o disfrutan de una ganancia. Una sociedad mercantil no come, no duerme, no tiene una función de utilidad que maximizar en ningún sentido psicológico. 

Sin embargo, la inmensa mayoría de los sistemas fiscales fragmentan el sujeto, sus fuentes y usos de renta, en diferentes figuras impositivas por razones diversas. Evidentemente, el caso paradigmático es el impuesto sobre sociedades que, como sabemos, trata a una sociedad como contribuyente autónomo, exigiéndole liquidar y pagar su propio impuesto sobre la renta, distinto y adicional al que pagarán sus accionistas cuando reciban el beneficio ya gravado. 

Por tanto, la última pregunta que debemos responder respecto al sujeto es la forma ideal en que se debe configurar la obligación que imponen las figuras de imposición directa. En particular, veremos dos vertientes del mismo problema. ¿Debe la obligación fragmentarse en diferentes obligaciones específicas o mantenerse como una única obligación general? ¿Debe la obligación ser personal, gravando al contribuyente como sujeto, o, por el contrario, ser real, gravando únicamente la renta dentro de nuestra soberanía fiscal? 



\subsubsection{Desagregación vertical: ¿?}
El debate en torno a la naturaleza misma que tienen las sociedades no está, ni de lejos, cerrado. Y, además, es la traducción de un debate jurídico mucho más antiguo que la propia doctrina alemana del siglo XIX libró con intensidad sorprendente. 

VON SAVIGNY, el gran sistematizador del Derecho civil alemán, defendió lo que se conoce como la Teoría de la Ficción. Solo el ser humano individual posee, por naturaleza, capacidad jurídica propia. Cuando el Derecho atribuye personalidad a una corporación, no está reconociendo nada que exista realmente, sino creando una ficción útil, una atribución artificial de derechos y obligaciones a una entidad que, en sí misma, carece de voluntad propia.\footnote{%
    El propio SAVIGNY llegó a vincular esa capacidad jurídica genuina con lo que llamaba el Volksgeist, el espíritu compartido de un pueblo, algo que ninguna corporación puede poseer.
} Frente a él, VON GIERKE desarrolló la \textbf{Teoría de la Realidad} (Genossenschaftstheorie). Sostuvo que una corporación, especialmente una con estructura organizativa compleja y duradera, desarrolla una voluntad propia, genuinamente distinta de la suma de las voluntades individuales de sus miembros. 

Lejos de quedar confinada, esta discusión se trasplantó Inglaterra a través de MAITLAND, traductor y divulgador de GIERKE, y desde allí a los Estados Unidos, donde terminó alimentando directamente el debate sobre si una corporación puede, o no, ser sujeto de derechos constitucionales propios, de responsabilidad penal autónoma, y (lo que aquí nos interesa) de capacidad contributiva propia a efectos fiscales.



El cierre que confirma que la pregunta sigue abierta: check-the-box y los desajustes híbridos

[Probable] Con este debate ya como contenido central —no como telón de fondo—, el material sobre Coase, Berle y Means y el régimen estadounidense de check-the-box que ya conocías adquiere ahora su función correcta: no como argumento autónomo, sino como la prueba empírica de que ninguna de las tres justificaciones a favor de la desagregación, ni tampoco la tradición de integración plena en contra, ha conseguido imponerse de forma limpia y definitiva. El propio hecho de que el ordenamiento estadounidense permita a determinadas entidades elegir voluntariamente si prefieren tributar como sujeto opaco o como transparente, con independencia de sus rasgos económicos reales de separación entre propiedad y control, confirma que ni siquiera el legislador más sofisticado ha encontrado un criterio objetivo capaz de zanjar, de una vez, el debate que Savigny y Gierke abrieron hace siglo y medio. Y cuando dos jurisdicciones distintas califican de forma distinta a la misma entidad —el fenómeno de los desajustes híbridos que ya trabajamos, y que la Acción 2 del proyecto BEPS de la OCDE trató de neutralizar—, lo que en realidad estamos viendo no es un simple problema técnico de coordinación internacional, sino la manifestación, a escala global, de que el mundo entero sigue sin ponerse de acuerdo sobre la pregunta más elemental de todas: si la sociedad mercantil es, de verdad, alguien a quien gravar, o solo el nombre que le damos al conjunto de personas físicas que, en último término, son las únicas que de verdad pagan.



\paragraph{Obligaciones específicas: }
Existen tres argumentos fundamentales a favor de tratar a la sociedad como sujeto autónomo. Las tres tradiciones distintas sostienen que, con independencia de cuál sea la verdadera naturaleza, existen razones sólidas para tratarla como sujeto tributario separado.

La primera es la Teoría del Beneficio. La corporación recibe del Estado un privilegio jurídico de valor económico muy real y nada trivial: la responsabilidad limitada. Esta ventaja permite a sus propietarios arriesgar únicamente el capital que han invertido, trasladando cualquier pérdida excedente sobre los acreedores de la empresa en lugar de sobre el patrimonio personal de sus accionistas. Es un privilegio que el Derecho no concede gratuitamente en ningún otro ámbito de la actividad económica, y que además genera un coste social identificable: externaliza parte del riesgo empresarial sobre terceros que no pueden protegerse frente a él con la misma facilidad con que lo hace un acreedor voluntario que negocia sus condiciones de antemano. Bajo esta lectura, el impuesto sobre sociedades no es un impuesto sobre la renta en sentido estricto, sino algo más próximo a una tarifa de uso por el disfrute de ese privilegio, exactamente el mismo tipo de argumento que ya trabajamos, en otro contexto y con otro objeto, a propósito del principio del beneficio frente a la capacidad de pago.

La segunda es la teoría de la retención en la fuente. Si la renta generada dentro de la sociedad no tributara mientras permanece retenida, cualquier accionista podría diferir indefinidamente su propia tributación personal simplemente absteniéndose de repartir dividendos, disfrutando mientras tanto de un crecimiento de su patrimonio libre de impuesto. El impuesto sobre sociedades funciona, bajo esta óptica, como un mecanismo de cierre, que impide que la envoltura jurídica se convierta en un refugio permanente frente a la tributación personal.\footnote{%
    Y esta función se vuelve especialmente insustituible cuando el accionista es no residente, porque en ese caso el Estado donde se genera la renta puede carecer, en la práctica, de cualquier otro instrumento eficaz para alcanzar a ese accionista directamente, dado que escapa a su jurisdicción personal.
} Por tanto, según esta segunda linea, la razón más sólida a favor de la desagregación no depende en absoluto de que la entidad sienta o no capacidad económica propia, sino de un problema puramente instrumental: sin un impuesto que grave la renta retenida dentro de la entidad, cualquier persona física podría diferir indefinidamente su tributación personal simplemente acumulando renta dentro de la envoltura jurídica en lugar de repartírsela a sí misma, disfrutando mientras tanto de un crecimiento libre de impuesto sobre esa renta ya generada pero formalmente no percibida.

La tercera es de naturaleza más administrativa que doctrinal. COASE (1937) explicó por qué existe la empresa como forma de organización, pero BERLE y MEANS (1932) añaden la pieza que aquí resulta decisiva:\footnote{%
    De bloques anteriores del temario (\authorfont{Ver Tema 3.A.15}), sabemos que COASE (1937) explicó que la empresa surge como una tecnología de organización, alternativa al mercado, que reduce los costes de transacción que supondría coordinar cada acto productivo mediante contratos individuales sucesivos. Esta perspectiva tiene una implicación directa: si la empresa es solo un dispositivo para ahorrar costes de contratación, no un sujeto de bienestar económico propio, entonces gravarla como sujeto autónomo corre el riesgo de confundir una tecnología organizativa con un sujeto de capacidad económica. Ahí es donde entran BERLE y MEANS (1932), dispuestos a solucionar esta aparente contradicción.
} en la gran sociedad de capital moderna, con miles o millones de accionistas dispersos, propiedad y control se separan de forma genuina, y esa separación es la razón práctica por la que resulta prácticamente inviable imputar directamente a cada accionista la parte proporcional de renta que le corresponde de una empresa cuya composición accionarial cambia, además, varias veces al día en cualquier bolsa de valores del mundo. Gravar en sede de la entidad no es, bajo esta tercera razón, una afirmación filosófica sobre si la sociedad es o no un sujeto real en el sentido de GIERKE; sino, sencillamente, el único punto del sistema donde la información necesaria para liquidar el impuesto está disponible con la fiabilidad y la inmediatez que la Administración necesita.

¿Qué criterio decide, entonces, cuándo conviene tratar a una entidad como opaca (sujeto autónomo) y cuándo como transparente (mera atribución a sus miembros)? La pista más reveladora es que la práctica comparada no trata igual a todas las formas de organización interpuesta. Las sociedades de capital son, casi universalmente, tratadas como opacas; mientras que las sociedades personalistas o de personas son, casi con la misma universalidad, tratadas como transparentes. La prueba más elocuente de que esta pregunta no tiene, ni siquiera dentro de un mismo ordenamiento, una respuesta única y cerrada es el régimen estadounidense conocido como check-the-box. Antes de existir, la calificación de una entidad como sociedad opaca o transparente dependía de un examen casuístico de cuatro factores (responsabilidad limitada, libre transmisibilidad de las participaciones, continuidad de vida jurídica independiente de sus miembros, y gestión centralizada), de modo que una entidad con tres o más de estos rasgos tributaba como sociedad opaca y con menos de tres, como transparente. El Tesoro sustituyó ese test material por un sistema de elección voluntaria. Determinadas entidades elegibles pueden simplemente marcar una casilla para decidir si prefieren tributar como entidad opaca o como transparente, con independencia de cuáles sean sus rasgos económicos reales de separación entre propiedad y control. Es difícil encontrar una prueba más directa de que la pregunta carece de una respuesta correcta única incluso dentro de un mismo sistema fiscal. Si existiera un criterio económico objetivo y verificable que determinara inequívocamente la respuesta, no tendría sentido dejarla al mero arbitrio contractual del contribuyente.

Hoy este elemento tiene una proyección internacional de primer orden. Cuando dos jurisdicciones distintas califican a la misma entidad de forma diferente, surge el fenómeno conocido como desajuste híbrido, que puede explotarse deliberadamente para lograr que un mismo gasto se deduzca en dos jurisdicciones a la vez, o que una misma renta no tribute en ninguna de las dos. La OCDE dedicó a este problema, dentro del proyecto BEPS (\authorfont{Ver Tema 4.B.18}), toda una Acción específica cuya sola existencia confirma que la pregunta está lejos de ser una cuestión meramente doctrinal: tiene consecuencias recaudatorias de primer orden precisamente porque distintos sistemas la han resuelto de forma distinta.


\paragraph{Obligaciones generales: }
Frente a estas tres justificaciones, existe una tradición que defiende exactamente lo contrario. Para estos autores, se debería eliminar el impuesto sobre sociedades como figura autónoma y atribuir directamente la renta societaria a los accionistas a medida que se genera. 

Conocida en la literatura anglosajona como teoría de la integración plena, tiene hitos institucionales muy concretos que conviene conocer. La Comisión CARTER (1966) canadiense, propuso ya un sistema de imputación con crédito pleno al accionista. Por otro lado, el propio Departamento del Tesoro de los Estados Unidos (\textit{Taxing Business Income Once}, 1992), estudió con detalle la viabilidad de gravar la renta empresarial una sola vez en lugar de dos. Por último, en un estudio de referencia elaborado por el profesor WARREN (1993), propuso convertir el propio impuesto sobre sociedades en un mecanismo de retención pura, acreditable en su totalidad contra la cuota personal del accionista, de modo que la carga final dependiera exclusivamente del tipo marginal de cada accionista individual y no de ningún tipo societario autónomo. WARREN y GRAETZ (2016) retomarían y actualizarían esta misma propuesta más de dos décadas después, señalando que Estados Unidos era ya, para entonces, el único gran país desarrollado que mantenía un sistema clásico sin ningún mecanismo de integración.

El obstáculo que ha impedido que ningún gran sistema fiscal adopte la integración plena para las grandes sociedades cotizadas no es doctrinal, sino exactamente el mismo problema administrativo que ya identificamos como tercera justificación a favor de la desagregación. La separación de propiedad y control hace prácticamente inviable atribuir la renta societaria a un accionariado disperso y cambiante con la inmediatez que el sistema exigiría. Es, en cierto modo, la misma dificultad que, en el apartado anterior sobre agregación intergeneracional, hacía que cualquier frontera temporal terminara topándose con su propio límite: aquí, cualquier intento de eliminar del todo la frontera entre sociedad y accionista termina topándose con el límite práctico de que, para las grandes corporaciones modernas, esa frontera es la única unidad sobre la que la Administración puede liquidar con la información y la fiabilidad que necesita.

\subsubsection{Desagregación horizontal: ¿residencia frente a territorio?}
Por último, debemos preguntarnos si conviene delimitar el alcance del sujeto tributario por su pertenencia a una comunidad política o por la mera localización física de la actividad que genera la renta. Es decir, debemos establecer el criterio de residencia o el criterio de territorio. O, dicho de otro modo, entender la imposición directa como una obligación personal o como una obligación real. 

La disyuntiva técnica es sencilla de enunciar. La obligación personal grava al residente por la totalidad de su renta, con independencia del país en que se haya generado. La obligación real grava únicamente la renta generada dentro del territorio del Estado, con independencia de la residencia de quien la obtiene. Además, esta disyuntiva no es solo una cuestión de reparto de soberanía fiscal entre Estados, sino que cada opción responde a un criterio de eficiencia económica distinto, formalizado desde los años sesenta (RICHMAN, 1963), bajo dos conceptos de neutralidad que conviene dominar con precisión.

\paragraph{Obligación personal: Capital Export Neutrality (CEN)}
La neutralidad en la exportación de capital (Capital Export Neutrality, CEN) se cumple cuando la carga fiscal total que soporta el capital de un residente es la misma, invierta donde invierta. Bajo CEN, la decisión de dónde invertir queda libre de distorsión fiscal, y el inversor elige el destino que ofrezca la mayor rentabilidad antes de impuestos, lo que asignaría el capital mundial a sus usos más productivos en términos globales. 

CEN se logra, técnicamente, mediante tributación por renta mundial combinada con un crédito por el impuesto ya pagado en el extranjero, de modo que el residente termine soportando, en todo caso, el tipo de su país de residencia. 

\paragraph{Obligación real: Capital Import Neutrality}
La neutralidad en la importación de capital (Capital Import Neutrality, CIN) exige, en cambio, que todo el capital invertido dentro de un mismo territorio soporte idéntica carga fiscal, con independencia de la residencia de quien lo aporta, y se logra mediante tributación exclusivamente territorial, exonerando la renta de fuente extranjera obtenida por los propios residentes. De este modo, un inversor doméstico y uno extranjero compitan en igualdad de condiciones dentro de ese mercado concreto.

El problema es que CEN y CIN son, en general, mutuamente incompatibles. Salvo que todos los países del mundo armonizaran sus tipos de gravamen sobre el capital, satisfacer una neutralidad exige sacrificar la otra. Un sistema construido sobre CEN es más eficiente desde la perspectiva de la asignación mundial del ahorro, pero menos competitivo para las empresas residentes que operan en mercados extranjeros de menor tributación, porque terminan pagando, vía crédito compensatorio, el tipo más alto de su país de origen. Por el contrario, un sistema construido sobre CIN protege la competitividad local de la inversión extranjera, pero puede incentivar a los propios residentes a desplazar su ahorro hacia jurisdicciones de baja tributación, precisamente porque esa renta extranjera queda exenta.

La literatura más reciente ha añadido un tercer criterio que conviene incorporar, porque resuelve una limitación que ni CEN ni CIN capturan bien. DESAI y HINES (2003) introdujeron la neutralidad en la propiedad del capital (CON), que requiere preguntase quién es el propietario más eficiente de un activo productivo dado. Esto resulta sobretodo relevante en un mundo de fusiones y adquisiciones transfronterizas, donde la productividad de un activo puede depender de quién lo controla y no solo de dónde está físicamente situado. 

CON se satisface, señalan los autores, tanto bajo un régimen de exención total de la renta extranjera adoptado universalmente como bajo un régimen de tributación mundial con crédito adoptado universalmente; lo que revela que la disyuntiva residencia frente a territorio no es, en realidad, tan determinante para la eficiencia global como CEN y CIN por separado sugerían, sino que lo decisivo es la coordinación entre sistemas, no la elección aislada de cada uno.\footnote{%
    No obstante, WEISBACH (2014) plantea una objeción metodológica sobre el conjunto de esta literatura de neutralidades. Como ninguna de éstas puede satisfacerse simultáneamente con las demás, ninguna proporciona, por sí sola, un criterio de diseño de política fiscal internacional completo. Es decir, permiten identificar dónde se produce una distorsión concreta, pero no indican cuál de las distorsiones posibles es preferible sacrificar, una decisión que exigiría remitirse a las razones últimas para gravar la renta del capital.
}

En definitiva, el resultado práctico de esta imposibilidad de optimizar simultáneamente es exactamente el mismo patrón de compromiso que vuelve repetirse en cada uno de los ejes anteriores. Ningún sistema fiscal comparado es puramente personal ni puramente territorial. La inmensa mayoría combina tributación mundial del residente con mecanismos de corrección de la doble imposición, junto con tributación territorial específica sobre determinadas rentas obtenidas por no residentes. ¿Por qué? Es, otra vez, la misma tensión irreductible entre comprehensividad ideal y viabilidad práctica que ha vertebrado todo el tema, ahora proyectada sobre el eje geográfico en lugar del temporal o del subjetivo interno al hogar o a la persona jurídica.

\clearpage
\section{La dimensión distributiva: cómo se construye la progresividad efectiva}
\puntosclave{
    
}

Con el objeto y el sujeto ya fijados, queda por resolver la última gran pregunta de este tema: una vez que sabemos qué gravamos y a quién, ¿cómo se traduce esa base en una carga tributaria efectivamente distinta según la capacidad económica de cada contribuyente? El Estado dispone, en realidad, de solo dos palancas para dar forma distributiva a un impuesto sobre la renta, y todo lo demás que vas a ver en este bloque —gastos fiscales, deducciones, créditos, escalas, tarifas fragmentadas por fuente— es instrumento al servicio de una de esas dos palancas, no una tercera cuestión autónoma. La primera palanca actúa sobre la base: excluir o reducir determinadas rentas del cómputo antes de aplicar cualquier tarifa. La segunda actúa sobre el tipo: aplicar una tarifa progresiva, proporcional, o fragmentada según el origen de la renta. Empiezo por la base, porque es ahí donde se libra el debate conceptual más profundo —qué es, en rigor, un "gasto fiscal", y si el propio concepto resiste el escrutinio al que ha sido sometido durante más de medio siglo—.

\subsection{Sobre la base: el gasto fiscal como concepto general}
[Seguro] El concepto nace con Stanley Surrey, entonces subsecretario del Tesoro estadounidense, en un discurso de 1967 y desarrollado con todo su aparato teórico en Pathways to Tax Reform: The Concept of Tax Expenditures (Harvard University Press, 1973): cualquier disposición tributaria que conceda una ventaja o subsidio a un grupo de contribuyentes o a una categoría de renta debe considerarse el equivalente funcional de un gasto público directo, con la única diferencia de que, en lugar de transferirse mediante una partida presupuestaria visible, se transfiere renunciando a una parte de la recaudación que correspondería según la estructura "normal" del impuesto. [Probable] La implicación de fondo, que da su verdadero peso conceptual a todo este subapartado, es que fijar qué es un gasto fiscal exige, necesariamente, haber fijado antes un benchmark o sistema tributario de referencia respecto del cual medir la desviación —y aquí es donde este bloque se conecta directamente con el Bloque 1 completo: la propia pregunta "¿qué es renta?" que dedicamos tres subapartados enteros a resolver es, en el fondo, la pregunta previa e indispensable para poder decir, más adelante, qué constituye un gasto fiscal y qué no. Sin haber resuelto primero qué contamos como renta —SHS puro, renta-fuente, algo intermedio—, la propia noción de "desviación" carece de referente.

[Seguro] Esta dependencia de un benchmark no siempre consensuado es, precisamente, el punto de la crítica más seria que el concepto de Surrey ha recibido nunca, y que conviene presentar con el mismo peso argumental que la propia idea de Surrey, porque el debate sigue vivo hoy. Boris Bittker, profesor de Yale y la voz crítica de mayor autoridad frente a Surrey, sostuvo —ya en su primera réplica ("The Tax Expenditure Budget: A Reply to Professors Surrey and Hellmuth", National Tax Journal, 1969) y con mayor desarrollo después— que todo el edificio del gasto fiscal descansa sobre un benchmark de "impuesto sobre la renta ideal" que es, en sí mismo, inherentemente subjetivo y no consensuado: no existe ningún acuerdo doctrinal cerrado sobre si, por ejemplo, la propia progresividad de la tarifa forma parte de ese "impuesto normal" o es ya, en sí misma, una desviación —y si algo tan central como la tarifa progresiva queda fuera del consenso sobre qué es "normal", advertía Bittker, el propio ejercicio de catalogar "desviaciones" se vuelve arbitrario, dependiente de qué economista o qué Administración fije el benchmark en cada caso. [Probable] Surrey nunca negó del todo esta objeción —de hecho, como recuerda la literatura posterior sobre su obra, aceptó que la tarifa progresiva en sí misma pertenecía al sistema "normal" mientras que los tipos preferenciales para categorías concretas de renta sí eran gasto fiscal—, pero la crítica de Bittker dejó una marca permanente: cualquier catálogo de gastos fiscales, en cualquier país, es el resultado de una elección de benchmark que podría, legítimamente, haberse trazado en otro lugar, y esa elección nunca es políticamente neutral, por mucho que se presente con el aparato técnico de la contabilidad presupuestaria.

[Seguro] Sentado el concepto y su crítica, conviene desarrollar en abstracto —sin el catálogo de casos concretos que corresponde a otros temas— por qué el legislador puede preferir la vía fiscal al gasto directo para perseguir el mismo objetivo de política pública, con dos razones ya bien identificadas en el propio documento original y que aquí generalizo. La primera es la disyuntiva entre desplazamiento e incentivo: si el Estado provee directamente un bien que el sector privado también estaría dispuesto a financiar parcialmente, el gasto directo puede desplazar la contribución privada preexistente, mientras que un incentivo fiscal, al reducir el precio relativo de la conducta subvencionada, puede atraer contribución privada adicional —aunque, como ya vimos con el ejemplo de las donaciones benéficas, también subvenciona inevitablemente la conducta que ya se habría producido sin el incentivo, el llamado efecto inframarginal, que es la fuente de ineficiencia más citada en toda esta literatura—. La segunda es la de soberanía del consumidor frente a información imperfecta: un incentivo fiscal respeta las preferencias individuales sobre cómo destinar los recursos subvencionados, mientras que el gasto directo impone las preferencias del planificador público, con el coste de que los individuos, actuando de forma descentralizada, pueden carecer de la información necesaria para asignar eficientemente ese gasto —el mismo problema de información asimétrica, aplicado ahora al donante particular en lugar de al comprador de un bien de segunda mano, que ya trabajamos con Akerlof en el Bloque 2 del Tema 15—.

[Seguro] Cierro este subapartado con la cuestión de forma jurídica —reducción de base frente a deducción o crédito en cuota— que ya conoces bien en su desarrollo, y que aquí basta con fijar en su formulación más general y su fundamento en la propia obra de Surrey. Una reducción en la base beneficia, en términos de ahorro fiscal absoluto, más a quien tributa a un tipo marginal más alto —el fenómeno que la literatura describe, de forma muy gráfica, como subsidio invertido (upside-down subsidy): cuanta más renta tiene el contribuyente, mayor es, en términos absolutos, el beneficio fiscal que recibe por idéntico comportamiento—, mientras que un crédito de importe fijo beneficia por igual a cualquier contribuyente con cuota suficiente para absorberlo. [Probable] Resulta revelador, y conviene que lo sepas porque cierra el círculo con la propia autoría del concepto, que el propio Surrey no fue neutral entre ambos instrumentos: como ha documentado la literatura reciente sobre su obra (Mehrotra y Zelenak, análisis de la trayectoria intelectual de Surrey, Tax Notes, 2022), su rechazo a los gastos fiscales estaba motivado, en buena parte, por el hecho de que la mayoría adoptaban precisamente la forma de deducciones —regresivas por diseño—, y consideraba los créditos reembolsables —que benefician igual, o incluso más en términos relativos, a la renta baja— sustancialmente menos objetables, en coherencia con su defensa paralela de la equidad horizontal como principio rector de cualquier reforma fiscal seria. Es, en cierto modo, la misma distinción que en el ámbito europeo había anticipado, desde una tradición intelectual distinta, Richard Titmuss con su concepto paralelo de "bienestar fiscal" (fiscal welfare, años cincuenta): el reconocimiento de que el sistema tributario reparte beneficios sociales de forma tan real como el gasto público directo, aunque de forma mucho menos visible y, con frecuencia, bastante menos progresiva.

\subsection{Sobre el tipo: la tarifa como instrumento alternativo}
Fijada la primera palanca —intervenir sobre la base—, queda la segunda: una vez determinada la base gravable, el Estado puede dar forma a la progresividad efectiva del impuesto mediante el diseño de la tarifa, con independencia de que la base se haya definido de forma comprehensiva o recortada por gastos fiscales. Y aquí, otra vez, la pregunta general se despliega en varias vertientes que conviene tratar como facetas de una sola cuestión —cómo repartir la carga entre tramos de renta y entre fuentes de renta—, no como preguntas autónomas.

[Seguro] La primera vertiente es la más elemental: tipo único frente a escala de tramos. Un impuesto de tipo único —el llamado flat tax— no está necesariamente reñido con la progresividad, y conviene precisar esta distinción con cuidado porque suele confundirse: lo que sí es constante, bajo un tipo único, es el tipo marginal —el que se aplica al último euro ganado—, pero si el diseño incorpora un mínimo exento suficientemente amplio, el tipo medio —la proporción de la renta total que efectivamente se paga— sigue creciendo con la renta, porque ese mínimo exento pesa proporcionalmente menos cuanto mayor es la renta total del contribuyente. Es exactamente el diseño que propusieron Robert Hall y Alvin Rabushka (The Flat Tax, Hoover Institution, 1985): un tipo único sobre una base próxima al consumo, combinado con una exención personal generosa, que preserva cierta progresividad por la vía del umbral y no por la vía de la escala. [Probable] Esta propuesta, puramente académica en su origen, encontró aplicación real y a gran escala precisamente donde el Bloque 0 predeciría que encontraría menor resistencia institucional: las economías postcomunistas de Europa del Este, que necesitaban reconstruir de cero su capacidad administrativa fiscal tras el colapso del sistema anterior —el mismo argumento de Besley y Persson sobre la construcción endógena de capacidad fiscal—, adoptaron el flat tax precisamente por su simplicidad administrativa: Estonia en 1994, seguida de las demás repúblicas bálticas y, de forma más mediática, Rusia en 2001 con un tipo único del 13%, un experimento seguido con enorme atención por la literatura de finanzas públicas comparada durante toda la década siguiente.

[Seguro] La segunda vertiente, mucho más consecuente para el resto del tema, es la de la tarifa fragmentada por fuente —el problema cedular o dual que ya trabajaste, en su aplicación concreta, en los Temas 16 y 17—: ¿por qué un sistema decide aplicar tarifas distintas según el origen de la renta —trabajo frente a capital— en lugar de una escala única sobre toda la renta agregada? [Seguro] La justificación de eficiencia más rigurosa procede de un principio muy anterior a cualquier debate sobre fiscalidad dual: la regla de la elasticidad inversa de Frank Ramsey ("A Contribution to the Theory of Taxation", Economic Journal, 1927), formulada originalmente para impuestos sobre bienes y trasladable, con las debidas cautelas, a la imposición sobre distintas fuentes de renta: la pérdida de eficiencia que genera un impuesto es menor cuanto menos elástica —menos sensible al propio impuesto— sea la base gravada, lo que sugiere gravar con mayor intensidad las bases más rígidas y con menor intensidad las más sensibles. El capital, por su movilidad internacional —ya trabajada con Zodrow y Mieszkowski—, es estructuralmente más elástico que el trabajo, que apenas puede desplazarse de jurisdicción con la misma facilidad; de ahí que la lógica de Ramsey, aplicada a esta disyuntiva, apunte hacia tipos más bajos sobre el capital que sobre el trabajo —exactamente el diseño dual que adoptaron los países nórdicos y que ya conoces con detalle—.

[Probable] Pero esta misma conclusión —tipos bajos, o incluso nulos, sobre el capital— tiene un desarrollo teórico mucho más radical dentro de la teoría de la imposición óptima dinámica, que conviene traer aquí como el extremo lógico de la lógica de Ramsey, sin desarrollarlo con el aparato matemático que corresponde propiamente al Tema 11: el resultado de Christophe Chamley (1986, Econometrica) y Kenneth Judd (1985, Journal of Public Economics), formulado de forma independiente por ambos y conocido como el resultado Chamley-Judd, demuestra que, en modelos dinámicos de crecimiento óptimo con horizonte infinito, el tipo impositivo sobre el capital debería converger a cero en el largo plazo —no simplemente ser bajo, sino literalmente nulo—, porque cualquier impuesto permanente sobre el capital distorsiona de forma acumulativa, periodo tras periodo, la decisión de ahorro, con un coste de eficiencia que crece sin límite a medida que se prolonga el horizonte temporal. [Seguro] Este resultado, que durante más de treinta años se consideró uno de los hallazgos más robustos y menos discutidos de toda la teoría de la imposición óptima, ha sido sometido a una revisión reciente y rigurosa que conviene que conozcas, porque es el ejemplo más actual de cómo ni siquiera los resultados matemáticamente más sólidos de esta disciplina están cerrados para siempre: Ludwig Straub e Iván Werning ("Positive Long-Run Capital Taxation: Chamley-Judd Revisited", American Economic Review, 2020) demuestran que la conclusión de tipo cero no se sigue, en realidad, de los propios modelos de Chamley y Judd bajo supuestos generales, sino solo bajo una condición muy específica sobre la elasticidad de sustitución intertemporal del consumo; para valores de esa elasticidad por debajo de uno —el rango que buena parte de la literatura empírica considera más plausible—, el tipo óptimo de largo plazo sobre el capital es, de hecho, positivo y significativo, pudiendo mantenerse elevado durante siglos antes de converger, si es que converge, hacia cero. Es, probablemente, el ejemplo más reciente y mejor documentado de todo este tema de cómo un consenso teórico aparentemente cerrado puede resultar, tras un examen más detenido, mucho más frágil de lo que parecía.

[Probable] La tercera vertiente, ya apuntada de forma intuitiva en el documento original a propósito del "cuello de botella" que las sociedades evitarían bajo una tarifa progresiva, explica por qué la tarifa societaria es, casi universalmente, proporcional y no progresiva: bajo una escala progresiva, cualquier discontinuidad en el tipo marginal —un salto de tramo— genera un incentivo a manipular la magnitud gravada para quedar justo por debajo del umbral, un fenómeno que la literatura empírica moderna denomina agrupamiento o bunching en los puntos de quiebre de la tarifa (Emmanuel Saez, "Do Taxpayers Bunch at Kink Points?", American Economic Journal: Economic Policy, 2010, la referencia metodológica de cabecera para detectar y medir empíricamente este fenómeno en cualquier tarifa discontinua). Una persona física, con una única renta personal, tiene escaso margen para fragmentar artificialmente su actividad económica en varias entidades para eludir ese salto de tramo; una empresa, en cambio, puede reorganizarse —fragmentarse en varias sociedades más pequeñas, como ya vimos con Garicano, Lelarge y Van Reenen a propósito de los umbrales regulatorios franceses— con relativa facilidad para mantenerse por debajo de cualquier umbral de tarifa progresiva, lo que explica por qué casi ningún sistema comparado arriesga esa distorsión adicional sobre la persona jurídica, reservando la progresividad, cuando existe en el ámbito societario, a mecanismos de tramos de tamaño —cifra de negocios, no tipo marginal sobre el beneficio— que ya conoces del Tema 17.

[Probable] Cierro el bloque completo con una nota breve, de menor calado teórico que las tres anteriores pero que conviene no perder de vista: el momento en que se cobra el impuesto —mediante retención en origen o mediante liquidación diferida al final del periodo— no es neutral respecto de cómo el propio contribuyente percibe la carga que soporta, un fenómeno que la literatura reciente de economía del comportamiento ha formalizado bajo el nombre de saliencia fiscal (tax salience): Raj Chetty, Adam Looney y Kory Kroft ("Salience and Taxation: Theory and Evidence", American Economic Review, 2009) muestran, con evidencia experimental y de campo, que los contribuyentes reaccionan de forma sistemáticamente más débil ante impuestos que no son visibles en el momento de la decisión económica —precisamente el caso de la retención en origen, descontada de la nómina antes de que el trabajador llegue siquiera a verla— que ante impuestos igualmente costosos pero explícitamente mostrados en el momento del pago. Es la formalización moderna, con metodología experimental rigurosa, de la misma intuición sobre "ilusión monetaria" que el documento original ya apuntaba de forma correcta pero sin este respaldo: la retención no es solo un instrumento de gestión de liquidez para el Estado, es también, y quizá sobre todo, un instrumento que reduce la resistencia política a la propia carga tributaria, precisamente por hacerla menos saliente ante quien la soporta.

\clearpage
\section{Comparación internacional: los tres debates proyectados en la práctica}
Cierro el tema con la pregunta que da sentido práctico a todo lo anterior: si los tres bloques precedentes muestran que ninguna de las grandes preguntas de la imposición directa tiene una respuesta teóricamente cerrada —ni el objeto, ni el sujeto, ni la forma de dar progresividad efectiva a la base—, ¿cómo se traduce esa falta de cierre teórico en la práctica de los sistemas fiscales realmente existentes? No organizo este bloque por países ni por cronología de reformas —eso degeneraría en un catálogo sin hilo argumental—, sino, exactamente como ya anunciamos, recorriendo las mismas tres preguntas que vertebraron los Bloques 1, 2 y 3, y mostrando, para cada una, qué respuestas distintas han dado los sistemas comparados y por qué. Vas a comprobar algo que conviene anticipar como conclusión: ningún país ha resuelto ninguna de las tres preguntas de forma pura. Cada sistema fiscal real es, sin excepción, un compromiso específico entre los polos teóricos que hemos ido enfrentando en cada bloque —y ese compromiso, más que cualquier "modelo nacional" cerrado, es el verdadero objeto de la comparación internacional.

\subsection{El objeto: fuente, entrada, gasto y patrimonio en la práctica comparada}
[Probable] Ningún sistema fiscal moderno aplica SHS en su forma pura, y los restos de la teoría de la fuente que anticipamos en el Bloque 1 no son un fenómeno residual: son visibles en la propia arquitectura de casi cualquier impuesto comparado. El caso más elocuente, y el que conviene retener por su valor pedagógico, es el propio impuesto de Pitt de 1799 que trabajamos en el Bloque 0: su estructura en cédulas —Schedule A a F, cada una gravando una fuente distinta de renta (tierra, valores públicos, actividad comercial, empleo, dividendos)— sobrevivió en el Reino Unido, con modificaciones, durante más de dos siglos, hasta su sustitución por un sistema más unificado ya en el siglo XXI —una persistencia institucional que desmiente cualquier idea de progreso lineal hacia la comprehensividad—. Francia, por su parte, mantiene todavía hoy un sistema organizado explícitamente en categorías de renta (catégories de revenus: traitements et salaires, bénéfices industriels et commerciaux, revenus fonciers, entre otras), cada una con sus propias reglas de determinación del rendimiento neto, en lugar de una definición unificada de renta seguida de un único cómputo — el resto conceptual de la teoría de la fuente más persistente entre las grandes economías europeas.

[Probable] La propuesta de Kaldor del gasto personal, que en el Bloque 1 identificamos como fracasada en su forma pura —India y Sri Lanka, breve ensayo y abandono—, sí ha triunfado, en cambio, de forma parcial y fragmentaria: el tratamiento EET de la previsión social —exención en la aportación, exención durante la acumulación, tributación solo al rescate— que ya conoces con detalle del Tema 16 está hoy generalizado en prácticamente todos los sistemas comparados de la OCDE, constituyendo, de facto, una isla de tributación tipo gasto dentro de sistemas que, en todo lo demás, siguen la lógica de renta. Es la confirmación empírica de algo que ya intuíamos: cuando la propuesta teóricamente más elegante resulta inviable a escala completa, sobrevive igualmente, pero fragmentada en aplicaciones sectoriales concretas, en lugar de desaparecer del todo.

[Probable] Y la fragmentación entre renta y patrimonio que cerró el Bloque 1 muestra, en comparación internacional, la misma variedad de compromisos que ya trabajaste con detalle en el Tema 19: apenas un puñado de países —Noruega, Suiza, España— mantiene un impuesto general y anual sobre el patrimonio neto, mientras que la inmensa mayoría de la OCDE, incluida Alemania desde 1997 y Francia desde 2018 —sustituido este último por un impuesto limitado a inmuebles—, ha abandonado esa vía, optando por dejar que las divergencias entre renta gravada y capacidad económica real se acumulen sin corrección periódica explícita, salvo en el momento de la transmisión.

\subsection{El sujeto: agregación, desagregación y territorio en la práctica comparada}

[Seguro] La agregación familiar, que en el Bloque 2 analizamos con el aparato de Becker y Chiappori, se resuelve de tres formas distintas en el panorama comparado, y el dato ya presente en tu documento original merece ahora la lectura que solo el aparato teórico completo permite dar: diecinueve países de la OCDE gravan a los cónyuges de forma individual —la opción coherente con la lectura del modelo colectivo, que preserva la titularidad fiscal de cada perceptor—; cinco —Francia, Alemania, Luxemburgo, Portugal y Suiza— ofrecen subvenciones al matrimonio mediante sistemas de división de renta (splitting), que suman la renta familiar y la reparten a partes iguales antes de aplicar la tarifa, beneficiando sistemáticamente a las parejas con reparto interno más desigual de la renta —exactamente el resultado que, bajo la lente de Chiappori, cabría esperar que reforzara, y no corrigiera, cualquier desequilibrio de poder de negociación interno—; y solo tres —Estados Unidos, Irlanda y Noruega— mantienen un sistema de imposición familiar pura y plenamente agregada, sin división de renta que module su efecto sobre la progresividad. [Probable] Este mapa no es estático: la propia evidencia de Bick y Fuchs-Schündeln que trajimos al Bloque 2 documenta que los países que han transitado desde la imposición conjunta hacia la individual en las últimas décadas —Suecia en 1971, pionera en esta transición, seguida más tarde por otros países nórdicos— experimentaron incrementos medibles en la participación laboral femenina, el mismo mecanismo del tipo marginal sobre el perceptor secundario que ya trabajamos en detalle.

[Probable] La desagregación mediante personalidad jurídica muestra, en cambio, una convergencia mucho mayor entre sistemas —la opacidad de la sociedad de capital frente a la transparencia de la sociedad de personas es casi universal, como ya vimos—, con la singularidad estadounidense del check-the-box como la excepción que confirma la regla: ningún otro sistema comparado ha llevado tan lejos la posibilidad de elección voluntaria entre ambos regímenes, lo que explica por qué buena parte de la planificación fiscal internacional más sofisticada de las últimas tres décadas ha operado, precisamente, explotando los desajustes híbridos que se producen cuando una entidad que Estados Unidos trata como transparente es tratada como opaca por la jurisdicción de la contraparte —el problema que la Acción 2 de BEPS trató de neutralizar—.

[Seguro] Y en el eje residencia-territorio, la evolución más reciente y significativa la protagoniza precisamente Estados Unidos, el país que durante más tiempo sostuvo el modelo más próximo a la neutralidad en la exportación de capital (CEN) —tributación mundial con crédito por impuesto extranjero—: la reforma fiscal de 2017 (Tax Cuts and Jobs Act) desplazó al sistema estadounidense hacia un modelo cuasi-territorial, exonerando buena parte de los dividendos de fuente extranjera repatriados por sus multinacionales, mientras introducía simultáneamente un impuesto mínimo sobre determinadas rentas intangibles de baja tributación en el extranjero (GILTI) — una combinación que no encaja limpiamente en ninguno de los tres modelos de neutralidad que trabajamos en el Bloque 2 (CEN, CIN, CON), sino que ilustra, una vez más, el patrón de compromiso pragmático que atraviesa todo este bloque comparado.

\subsection{La dimensión distributiva: tipos marginales y gasto fiscal en la práctica comparada}

[Seguro] Cierro con el dato comparado más elocuente de todo el tema, porque resume en una sola serie temporal el arco completo que hemos recorrido desde el Bloque 0. Recordarás que el tipo marginal británico alcanzó el 99,25% durante la Segunda Guerra Mundial. Esa progresividad extrema de posguerra se mantuvo, con matices, durante toda la llamada edad de oro del Estado del bienestar: en 1975, el tipo marginal máximo del Reino Unido era del 83%, y el de Estados Unidos, del 70% —niveles que, bajo la lógica de Mill y Edgeworth que trabajamos en el Bloque 0, se justificaban explícitamente como aplicación del sacrificio marginal igual—. [Seguro] La contrarreforma que tu documento original ya identificaba, asociada a Thatcher y Reagan, se tradujo en una caída de aproximadamente cuarenta puntos porcentuales en poco más de una década: el tipo marginal británico cayó al 60% en 1980 —el primer presupuesto de Thatcher— y al 40% en 1988; el estadounidense cayó al 50% en 1982 y hasta el 28% con la Tax Reform Act de 1986, el nivel más bajo desde antes de la Segunda Guerra Mundial. [Probable] Desde entonces, los tipos marginales máximos se han estabilizado, en la mayoría de la OCDE, en una banda mucho más estrecha que la de los años setenta —entre el 40% y el 50% en la mayoría de las grandes economías—, sin volver nunca a los niveles de posguerra ni tampoco descender de forma sostenida por debajo de ese nuevo suelo, lo que sugiere que la contrarreforma de los ochenta no fue un episodio pasajero sino un nuevo equilibrio duradero entre los mismos dos extremos —sacrificio igual frente a Leviatán— que enmarcamos teóricamente en el Bloque 0.

[Probable] El volumen del gasto fiscal, por su parte, no sigue esta misma trayectoria descendente: los datos que ya aportaba tu documento original —cerca de 80.000 millones de euros en España en 2020, y una cifra equivalente a casi el 40% de la recaudación proyectada del impuesto sobre la renta en Estados Unidos en 2014— muestran que, mientras los tipos se recortaban drásticamente desde los años ochenta en el marco de la doctrina de "ensanchar bases, bajar tipos" que ya trabajamos en los Temas 16 y 17, el gasto fiscal no se ha reducido en proporción equivalente en la mayoría de los sistemas comparados —confirmando, con datos agregados y no solo con el argumento teórico de Bittker, que el catálogo de desviaciones respecto de cualquier benchmark normativo tiende, una vez creado, a persistir con independencia de las reformas que afecten a la otra palanca distributiva, la de la tarifa.

Con este recorrido comparado queda cerrado el Tema 12 completo. El argumento que lo ha vertebrado, desde el Bloque 0 hasta aquí, es uno solo, repetido en cada uno de sus niveles: la imposición directa sobre la renta no es la aplicación de un ideal teórico cerrado, sino una sucesión de compromisos históricamente contingentes entre exigencias que nunca se satisfacen simultáneamente —comprehensividad y verificabilidad, equidad y eficiencia, capacidad de pago y recelo frente al Leviatán, fidelidad a la unidad real de bienestar y conveniencia administrativa—. Cada uno de los sistemas fiscales que hemos comparado en este último bloque es, precisamente, uno de esos compromisos posibles, y ninguno agota las opciones que la teoría deja abiertas.


























\clearpage
\section{Imposición directa: Marco conceptual}


Los impuestos gravan la capacidad contributiva de los individuos y pueden tener naturaleza directa, cuando gravan la renta o la riqueza, o indirecta, cuando gravan el empleo de esa renta o riqueza, es decir, el gasto.

En general los impuestos directos tienen en consideración las características individuales del contribuyente por lo que pueden ajustar la tributación a estas características. Los impuestos indirectos por el contrario se establecen sobre las transa1.:1.:iones con independencia de las circunstancias personales del consumidor.

Realizamos esta separación entre empresas y hogares por claridad e>..positiva, si bien el SEC 201 O no la realiza. Así:

• Los impuestos corTientes sobre l a renta, el patrimonio, etc., comprenden todos los pagos obl igatorios sin contrapartida, en efectivo o en especie, recaudados periódicamente por la:, administraciones públicas y por el resto del mundo sobre la renta y el patrimonio de las unidades institucionales. así como algunos impuestos periódicos que no se exigen ni sobre dicha renta ni sobre dicho patrimonio.

• Los impuestos sobre el capital son impuestos que j,,rravan a intervalos i rregulares y muy poco frecuentes el valor de los activos o el patrimonio neto ele las unidades instinicionales o el valor de los activos transferidos entre unidades institucionales como resultado de sucesiones, clonaciones ínter vivos u otras trnnsferencias.

Es decir. a efectos del SEC. la diferencia importante es la periodicidad más que la unidad institucional ( cont1ibuyente) que realice el pago o si este recae sobre la renta o el patrimonio. En nuestro caso detallaremos un poco más a la hora de analizar la imposición directa, como ya hemos sefialado.


l l . I l l . Progrcsividad y redistribución

Los impuestos directos suelen presentar una base imponible limitada y circunscrita a unos deseos del l egislador y menor a l a de los indirec1os. Esto se cumple generalmente, ya que los indirectos recaen sobre el consumo y los directos principalmente sobre la renta, cuya base en el conjunto de la economía es más reducida.

En general. aquellos sistemas tributarios basados en una mayor proporción de imposición directa suelen ser más progresivos. Los impuestos directos suelen contener una configuración progresiva en su diseño, especialmente en la base imponible, y. por tanto, al exigirse además tipos  impositivos distintos en func ión ele la renta del consumidor, en la práctica suponen una adaptación de la carga tributaria según la capacidad económ ica del contribuyente. Es más, al recaer también sobre e.1 patrimonio este tipo de tributos incorporan un efecto redistributivo de l a riqueza entre todos los contribuyentes.

111. La imposición directa 
Los sistemas fiscales responden a dos preguntas fundamentales a la hora de diseñar la imposición directa:

• ¿Qué rentas gravar? Es decir, la delimitación del hecho y base imponible. Ü\c'sde el punto de vista de la eficiencia. la imposición directa al ser una de las formas de impo::.ición más distorsionantes debería recaer sobre las rentas que menos pe��j udican la::. decisiones de los agentes, por ejemplo. gravar el patrimonio y no la renta de los trabajadores.
• ¿Sobre qué unidad de renta? Es decir, la determinación del sujeto pasivo del cual se grava la renta obtenida. 


\subsection{¿Cuál debe ser el objeto de tributación directa?}
La primera cuestión que tenemos que tratar, una de las más controvertidas, es la determinación del objeto sujeto a gravamen. Por ejemplo, los impuestos sobre la renta no se determinan simplemente tomando toda la renta acumulada durante el año. Pero, ¿por qué? ¿es está un objeto fruto del consenso teoríco? ¿es este objeto óptimo para tomarlo como referencia? 

\subsubsection{Teorías sobre el objeto: ¿?}
Como podemos intuir, está muy lejos de ser una cuestión firme y libre de debate. 














En la mayoría de sistemas de imposición sobre la renta, existe una variedad de exenciones, deducciones y créditos fiscales que hacen que la base imponible sea menor que la renta total. ¿Cuáles son las justificaciones de estos agujeros en la base imponible? ¿Cuáles son sus implicaciones para la equidad del sistema fiscal? Analicemos brevemente la base del impuesto sobre la renta teóricamente apropiada y examinamos cuidadosamente las fuentes de desviación de los impuestos sobre la renta respecto de ese ideal teórico y sus implicaciones.

El punto de referencia que los economistas de las finanzas públicas utilizan para definir la renta es la definición de renta integral de HAIG-SIMONS, que define los recursos imponibles como la capacidad de un individuo para pagar impuestos. Esta capacidad de pago es igual al consumo anual potencial de un individuo, es decir, el consumo total del individuo durante el año, más cualquier aumento de su stock de riqueza. ¿Tiene sentido la definición de Haig-Simons como objetivo del diseño de la base imponible? Recuérdese que existen dos aspectos para medir la equidad de un sistema fiscal. La equidad vertical se alcanza cuando los contribuyentes de renta alta pagan una proporción mayor de su renta en impuestos, y la equidad horizontal se alcanza cuando contribuyentes idénticos pagan la misma cantidad en impuestos independientemente de sus elecciones.

Utilizar la definición de renta de HAIG-SIMONS para determinar la base sobre la que se pagan los impuestos mejora la equidad vertical porque quienes tienen más recursos pagan más impuestos, aunque obtengan esos recursos a través de un canal no gravado (como ocurre con los pagos en especie). Así, si usted y yo tenemos los mismos salarios en efectivo, pero yo dispongo de una vivienda proporcionada por mi empleador, yo debería pagar más impuestos.

El enfoque de HAIG-SIMONS también mejora la equidad horizontal al garantizar que las personas que son iguales en términos de sus recursos subyacentes paguen la misma cantidad de impuestos independientemente de la forma en que elijan recibir o gastar sus recursos. Bajo el sistema actual, si yo elijo recibir mi remuneración en forma de salarios y usted elige, en cambio, recibir parte de su remuneración en forma de pagos en especie proporcionado por el empleador, usted paga menos impuestos que yo. Pasar a un enfoque de HAIG-SIMONS, que trata todas las formas de remuneración como renta, mejoraría la equidad horizontal.

Así, adherirse a una definición de HAIG-SIMONS podría mejorar enormemente la equidad de un sistema fiscal tanto en la dimensión de la equidad vertical como en la de la equidad horizontal. Sin embargo, aplicar una definición de HAIG-SIMONS en la práctica es difícil y, en ocasiones, muy controvertida.\footnote{%
    ¿Por qué? Al margen de las que se expondrán a continuación, existen también, al menos otras dos dificultades para aplicar una definición de HAIG-SIMONS en el sistema fiscal:
\begin{la}
    \item \textbf{¿Cómo calcular la capacidad de pago de un individuo?}. Supongamos que dos individuos tienen la misma renta en un año cualquiera, pero que uno sufre un gran incendio en su casa y tiene que gastar el $20\%$ de su renta para reparar los daños. En ese año, esta persona tendrá un menor consumo y menores aumentos de riqueza, o una menor capacidad para pagar impuestos, y esto debería reflejarse ajustando su base imponible. El deseo de tener en cuenta los gastos que no están asociados con el consumo deseado es la justificación de una de las principales deducciones de la renta imponible permitidas por el código fiscal, la deducción por pérdidas de bienes y siniestros;
    \item \textbf{¿Cómo computar los costes derivados de la actividad económica?}. Debido a que la definición integral de renta se refiere únicamente al incremento neto de los recursos durante el período, cualquier coste legítimo de hacer negocios debería deducirse de la renta de una persona. Sin embargo, a menudo surgen dificultades al definir los costes empresariales legítimos (comidas de negocios, transporte, publicidad, etc.). A modo de ejemplo, la publicidad sería totalmente deducible a efectos de la definición de HAIG-SIMONS, pero las comidas no. ¿Por qué? Porque el consumidor obtiene cierto valor de consumo, y este valor de consumo debería incluirse en su renta de HAIG-SIMONS. En teoría, solo debería poder deducir de su renta el coste total del almuerzo menos el propio valor de consumo derivado de ese almuerzo.
\end{la}
}




La justificación clásica de las desviaciones respecto de una definición integral de renta es la posibilidad de que reducir los impuestos sobre determinadas actividades genere beneficios externos para la sociedad. 

Analicemos dos desviaciones importantes respecto de la definición de HAIG-SIMONS que normalmente se justifican por el hecho de que es probable que el mercado privado proporcione en cantidad insuficiente algún bien o actividad: las donaciones benéficas y los gastos en vivienda.

\subsubsection{Subvención fiscal: ¿Por qué no aplicar la definición de renta integral?}
Un excelente ejemplo de la aplicación de la justificación basada en los beneficios externos es que las donaciones a organizaciones benéficas pueden deducirse de la renta imponible. Debido al problema del free rider (cuando los costes son privados pero los beneficios son públicos, los individuos proporcionarán bienes públicos en cantidad insuficiente), es probable que el mercado privado proporcione un apoyo benéfico insuficiente para muchos bienes públicos.

Supongamos, por ejemplo, que al gobierno le preocupa que el sector privado no esté proporcionando fondos suficientes para construir refugios para las personas sin hogar, lo cual constituye un caso clásico de bien público. Una forma de abordar este problema sería subvencionar las donaciones benéficas destinadas a las personas sin hogar con el fin de aumentar el apoyo del sector privado, y un medio de subvencionar las donaciones benéficas consiste en permitir que los individuos deduzcan de su renta imponible la cantidad que donan a organizaciones benéficas. Al hacerlo, el gobierno hace que las donaciones benéficas sean baratas en relación con otros tipos de consumo, que deben financiarse con dólares después de impuestos. El sistema fiscal permite a quienes detallan sus deducciones restar de la renta imponible las contribuciones a organizaciones benéficas, hasta cierta cuantía.

Supongamos que Ana está decidiendo si comprar una taza de café o donar esa cantidad a un refugio para personas sin hogar. Si compra la taza de café, tendrá que pagar el impuesto sobre la renta sobre ese euro a un tipo impositivo $t$. Por tanto, para obtener la cantidad requerida para la taza de café ($x$), Ana tendrá que ganar primero $x/(1 − t)$. Por ejemplo, si su tipo impositivo es del $50\%$, tendrá que ganar $2x$ dólares. Sin embargo, si Ana gasta $x$ en la donación, no paga ningún impuesto sobre esa cantidad. Así que puede utilizar $x$ de sus ingresos para dar $x$ al refugio. Por tanto, el precio relativo de las donaciones benéficas, en relación con otros tipos de consumo, es $x/[x/(1 − t)]$, $(1 − t)$. Con un tipo impositivo del $50\%$, el precio relativo de las donaciones benéficas es de solo $0.5x$ dólares porque Ana tiene que ganar $2x$ para obtener una taza de café, pero solo tiene que ganar $x$ para donar $x$ a una organización benéfica. Este tratamiento fiscal genera un beneficio (hace que donar a organizaciones benéficas sea mucho más atractivo) a costa de desviarse del estándar, puesto que la parte de la renta imponible que se dona a organizaciones benéficas no está gravada, de modo que los impuestos se aplican a una medida de la capacidad de pago que no es integral.

Sin embargo, existe otro enfoque que el gobierno podría adoptar para apoyar la provisión del bien público: podría proporcionar él mismo el bien. En lugar de inducir indirectamente a los individuos privados a aumentar sus contribuciones benéficas concediéndoles una desgravación fiscal, el gobierno podría simplemente gastar su propio dinero para mejorar los refugios para personas sin hogar. Entonces, ¿por qué el gobierno elige desviarse del estándar de HAIG-SIMONS en lugar de aumentar el gasto público? Existen al menos dos posibles razones.

\paragraph{Desplazamiento vs. incentivos: Efectos marginales e inframarginales}
En primer lugar, la provisión pública puede desplazar las contribuciones privadas al bien público. Si el gobierno subvenciona los refugios para personas sin hogar, lo más probable es que disminuya la cantidad de donaciones benéficas privadas destinadas a esos refugios.

Sin embargo, cuando el gobierno subvenciona fiscalmente las donaciones benéficas, puede atraer, o aumentar, las contribuciones privadas. Esto ocurre porque la subvención fiscal del gobierno reduce el precio relativo de las donaciones benéficas. Un precio relativo más bajo de las donaciones benéficas aumenta las donaciones privadas, tanto mediante efectos de sustitución (el precio de las donaciones benéficas ha disminuido) como mediante efectos renta (como Ana está pagando menos impuestos, es más rica y puede donar más a organizaciones benéficas).

Al mismo tiempo, cuando el gobierno subvenciona fiscalmente una conducta, como las donaciones, también está concediendo una desgravación fiscal a quienes ya habrían apoyado a las personas sin hogar incluso sin esta desgravación fiscal. Esto puede considerarse horizontalmente equitativo porque la desgravación fiscal se aplica a todos los que hacen donaciones a organizaciones benéficas. Pero aumenta significativamente el coste de este agujero en la base de HAIG-SIMONS: el gobierno subvenciona no solo las nuevas donaciones a organizaciones benéficas, sino también las antiguas donaciones a organizaciones benéficas que habrían existido incluso sin esta desgravación fiscal. Por eso, Cuando los economistas analizan el impacto de desgravaciones fiscales como la correspondiente a las contribuciones benéficas, a menudo distinguen entre los impactos marginales e inframarginales de estas desgravaciones fiscales.\footnote{%
    Por ejemplo, supongamos que hay un millón de euros de donaciones privadas a las personas sin hogar sin ninguna desgravación fiscal. El gobierno permite entonces a los individuos deducir las donaciones benéficas de su renta imponible, a un tipo impositivo del $50\%$, y esto eleva las donaciones a las personas sin hogar a un millones y medio de euros. Por tanto, el gobierno ha fomentado $500.000$ euros de nuevas donaciones mediante esta desgravación fiscal, logrando su objetivo de ampliar las donaciones privadas. Sin embargo, para quienes habrían donado de todos modos el millón, ahora también existe una desgravación fiscal de $500.000$ euros que recompensa lo que ya estaban haciendo.
}

Lo que determina la eficiencia de una desgravación fiscal diseñada para fomentar un comportamiento es la proporción de la desgravación fiscal que va a quienes cambian su comportamiento frente a la proporción que va a quienes cuyo comportamiento no se ve afectado. Las desgravaciones fiscales más eficientes en términos de costes tienen grandes impactos marginales que proceden de quienes cambian su comportamiento. En ese caso, un gasto público relativamente pequeño puede generar grandes impactos marginales. Las desgravaciones fiscales menos eficientes en términos de costes tienen pequeños impactos marginales y grandes costes en términos de ingresos debido a quienes ya estaban realizando el comportamiento subvencionado. En ese caso, puede haber gastos públicos muy elevados sin que haya un gran impacto marginal.

Efectos de las subvenciones fiscales frente al gasto directo 

Por tanto, existe una disyuntiva que debe considerarse al decidir cuál de estos dos enfoques utilizar, el gasto directo o las subvenciones fiscales. Si el gobierno gasta el dinero directamente, aumenta los recursos públicos disponibles para las personas sin hogar pero, mediante el desplazamiento, reduce potencialmente los recursos privados disponibles. Si el gobierno proporciona una desgravación fiscal, induce nuevas donaciones, pero también gasta dinero para subvencionar las donaciones existentes. Esta disyuntiva puede resumirse con la siguiente pregunta: si el gobierno dispone de un euro de ingresos que quiere gastar en causas benéficas como los refugios para personas sin hogar, ¿cuál es la mejor manera de gastar ese dólar? Matemáticamente, el gobierno debería utilizar una desgravación fiscal en lugar de gasto directo si: el aumento de las donaciones benéficas por dólar de desgravación fiscal es menor a $1 -$ la reducción de las donaciones benéficas por dólar de gasto público.

Si se cumple esta desigualdad, el gobierno puede aumentar la cantidad total de donaciones benéficas subvencionando fiscalmente la beneficencia privada en lugar de gastar directamente sus ingresos en bienes o servicios benéficos.\footnote{%
    Existe un gran número de estudios que han tratado de estimar el efecto de las subvenciones fiscales sobre las donaciones benéficas. La conclusión general de estos estudios es que la elasticidad de las donaciones benéficas con respecto a su subvención es aproximadamente −1: por cada reducción del $1\%$ en el precio relativo de las donaciones benéficas, la cantidad donada aumenta un $1\%$. Esto significa que el aumento de las donaciones benéficas (el efecto marginal de la subvención fiscal) es igual a los ingresos fiscales perdidos por las desgravaciones fiscales concedidas a las donaciones existentes (el efecto inframarginal de la subvención fiscal); lo que corresponde al ejemplo anterior.
}

\paragraph{Soberanía del consumidor vs. información imperfecta}
La subvención fiscal también puede preferirse al gasto público directo por razones de soberanía del consumidor. Cuando el gobierno proporciona el gasto directamente, impone sus preferencias sobre cómo se gastan los fondos. Por el contrario, al ofrecer subvenciones fiscales a los individuos privados para que donen como deseen, el gobierno respeta directamente las preferencias de sus ciudadanos.

La desventaja de esta provisión descentralizada es que el sector privado puede no disponer de los mecanismos apropiados para garantizar una distribución eficiente del gasto privado. La mayoría de los individuos realizan muy poca investigación antes de hacer, por ejemplo, sus donaciones benéficas. Un estudio reciente encontró que dos tercios de todos los donantes estadounidenses declaran no realizar ninguna investigación sobre sus contribuciones benéficas cada año, y solo el $5\%$ realiza dos o más horas de investigación, en promedio, los estadounidenses dedican aproximadamente cuatro veces más tiempo a investigar compras de televisores (cuatro horas) y ocho veces más tiempo a investigar compras de ordenadores (ocho horas) que a inversiones más costosas en organizaciones benéficas (una hora).

Este enfoque tiene el potencial de conducir a donaciones que no proporcionan mucho beneficio al destinatario previsto, sino que, en cambio, únicamente recompensan a la organización que recauda los fondos.\footnote{%
    Un ejemplo dramático de esto es la U.S. Navy Veterans Association, que fue fundada en 2002 para proporcionar apoyo a veteranos de la Marina necesitados y que, para 2010, tenía delegaciones en 41 estados y había recaudado 100 millones de dólares. O quizá no. Resulta que todo esto era información falsa proporcionada por un individuo que utilizó teleoperadores directos para estafar este dinero a los donantes y quedárselo para sí mismo. Este fraude fue descubierto en 2010, y el culpable fue capturado dos años después. Quizá lo más inquietante sea que la mayor parte de lo que hacía esta «organización benéfica» era perfectamente legal: durante el juicio posterior, uno de los recaudadores de fondos con ánimo de lucro de esta Asociación declaró que el $90\%$ de todas las donaciones se retenía como comisión por recaudación de fondos, sin repercusiones legales.
} Por supuesto, también existen ineficiencias en la provisión pública directa de bienes públicos. La cuestión importante, para la que no existe una respuesta clara en este momento, es si el sector privado o el sector público es más eficaz a la hora de transformar cada euro de gasto benéfico en resultados beneficiosos.

\begin{specialblock}{\textbf{Ejemplo.-} Subvenciones por la compra de vivienda (España, hasta 2023)}
    Un segundo ejemplo de una desviación respecto de HAIG-SIMONS que está potencialmente justificada por motivos de externalidades es la subvención fiscal a la propiedad de la vivienda. Supongamos que usted está buscando un lugar donde vivir y encuentra una vivienda que le gusta con un valor de alquiler de mercado de $1.000€$ al mes. El propietario le ofrece o bien alquilársela por $1.000€$ al mes o bien vendérsela por $100.000€$. Para financiar esa compra de $100.000€$, tendría que obtener un préstamo para vivienda, una hipoteca, con pagos de intereses de $1.000€$ al mes (en este ejemplo). Su tipo impositivo es del $50\%$ y gana $4.000€$ al mes.

    Un verdadero enfoque de HAIG-SIMONS para definir su renta en cualquiera de los dos casos consistiría en añadir a su renta el valor neto de consumo que obtiene de vivir en esa vivienda y restar el coste que supone para usted hacerlo. Tanto si alquila como si es propietario, está consumiendo cada mes $1.000$ en servicios de vivienda, lo que se denomina el valor de alquiler imputado de la vivienda. Y, en cualquiera de los dos casos, está pagando $1.000€$ por el derecho a vivir en esa vivienda. Por tanto, su carga fiscal según HAIG-SIMONS no se vería afectada por su decisión de alquilar o comprar.
    
    El sistema fiscal español (actual) no incluye el valor de alquiler de la propia vivienda en la renta imponible. No obstante, el impuesto sobre la renta sí permite a los individuos deducir los intereses hipotecarios de su renta imponible, pero no les permite deducir los pagos de alquiler. Debido a esta deducción, existe una ventaja financiera en comprar en lugar de alquilar. ¿Por qué subvencionar la propiedad de la vivienda? La justificación más común que se ofrece para esta subvención a la propiedad de la vivienda es que la propiedad de la vivienda tiene externalidades positivas que el alquiler no tiene.\footnote{%
        Como escriben GLAESER y SHAPIRO (2002) para Estados Unidos: \textit{Para sus defensores, la deducción de los intereses de las hipotecas sobre viviendas es la piedra angular de la sociedad [estadounidense]. La propiedad de la vivienda da a las personas un interés en la sociedad y las induce a preocuparse por sus barrios y ciudades. Al subvencionar la propiedad, la deducción induce a las personas a invertir y, después, a tener un interés en nuestra democracia. La propiedad hace que las personas voten a favor de inversiones a largo plazo en lugar de transferencias a corto plazo}.
    } Por ejemplo, GLAESER y SHAPIRO (2002) revisaron un amplio conjunto de evidencia empírica que respalda la existencia de estas externalidades positivas: en relación con el alquiler, la propiedad de la vivienda está positivamente correlacionada con el activismo político y la conexión social, y los propietarios cuidan mejor sus propiedades, lo que conduce a un mayor valor de las viviendas circundantes (una externalidad financiera positiva para sus vecinos).\footnote{%
        ¿Consigue sus objetivos la subvención fiscal por la compra de una vivienda? Existe un amplio conjunto de evidencia que sugiere que esta subvención fiscal a la propiedad de la vivienda aumenta los gastos en vivienda, con una elasticidad del gasto con respecto a las subvenciones fiscales de aproximadamente $1$. Al mismo tiempo, no existe evidencia de que esta subvención haga que más personas compren viviendas. Por ejemplo, para Estados Unidos, la tasa de propiedad de la vivienda se ha mantenido esencialmente constante desde la década de 1950, en torno al $65\%$; sí aumentó bruscamente a finales de la década de 1990 y principios de la de 2000, pero volvió a invertirse durante la Gran Recesión. Dado que la subvención fiscal no parece aumentar la propiedad de la vivienda, pero sí aumenta los gastos en vivienda, parece que la subvención fiscal está \textbf{induciendo a los individuos a gastar más en viviendas} que habrían comprado de todos modos, incluso sin la subvención fiscal. Este hallazgo sugiere que la justificación para subvencionar fiscalmente la propiedad de la vivienda es débil por motivos de externalidades positivas.
    
    A partir de la evidencia de que la deducción de los intereses de las hipotecas sobre viviendas no tiene beneficios externos, no existe una justificación clara para esta desviación respecto de la definición integral de renta de HAIG-SIMONS. Sin embargo, la deducción de las hipotecas sobre viviendas sigue siendo una de las disposiciones más populares del código fiscal, y los políticos rara vez critican o intentan limitar esta disposición.
    }
    
    No obstante, si los responsables de las políticas quieren subvencionar la propiedad de la vivienda, existen formas mucho más eficientes y equitativas de hacerlo, como el crédito para compradores de primera vivienda analizado en GALE et al. (2006), que subvencionaría directamente la compra de la primera vivienda mediante un crédito fiscal en lugar de permitir una deducción por los intereses hipotecarios pagados.
\end{specialblock}


\subsubsection{¿Cómo diseñar la desviación? Deducción vs. Créditos}
Otra cuestión importante que surge cuando un sistema fiscal se desvía de la definición integral de renta de Haig-Simons es si aplicar esas desviaciones en forma de deducciones fiscales o créditos fiscales. Las deducciones fiscales permiten a los contribuyentes reducir su renta imponible en una determinada cantidad (por ejemplo, la cantidad correspondiente a las donaciones benéficas o a la deducción de los intereses de las hipotecas sobre viviendas); por tanto, las deducciones reducen el precio del comportamiento en cuestión a $(1 − t)€$, donde $t$ es el tipo impositivo marginal. Los créditos fiscales permiten a los contribuyentes reducir en una determinada cantidad el importe de los impuestos que deben al gobierno (por ejemplo, la cantidad que gastan en el cuidado de los hijos). Si el gasto de un individuo es inferior al importe del crédito, el crédito fiscal reduce a cero el precio del comportamiento en cuestión.

\paragraph{Consideraciones de eficiencia }

¿Qué enfoque es preferible? La respuesta no está clara desde el punto de vista de la eficiencia. Por ejemplo, considérese sustituir la actual deducibilidad de las donaciones benéficas por un crédito fiscal del $100\%$ para las donaciones benéficas de hasta $1.000€$. Para quienes actualmente donan menos, el crédito proporciona un incentivo mucho mayor para aumentar las donaciones hasta ese nivel porque es gratuito (los pagos de impuestos disminuyen en un euro por cada euro donado). Una vez que una persona dona más de $1.000€$, ya no existe ningún beneficio adicional derivado del crédito fiscal. Por otro lado, la deducción continúa proporcionando cierto incentivo para donar incluso más allá del límite de 1.000 dólares (después de que se haya agotado el incentivo del crédito). 

La disyuntiva a la que se enfrenta el gobierno es entre un sistema que subvenciona parcialmente todas las donaciones (la deducción) y uno que subvenciona completamente algunas donaciones y otras no las subvenciona en absoluto (el crédito). Qué política, deducción o crédito, es más eficiente viene determinado por dos consideraciones. La primera es la naturaleza de la demanda del bien subvencionado. En la realidad, las elasticidades de la demanda pueden diferir dependiendo de la magnitud del cambio de precio. Para algunos bienes, la demanda individual puede ser muy elástica en respuesta a grandes reducciones del precio, pero no muy elástica en respuesta a pequeñas reducciones del precio; en tales casos, los créditos pueden provocar un aumento mayor de las donaciones benéficas (por ejemplo) que las deducciones, porque los créditos conducen a mayores reducciones del precio.

En segundo lugar, los responsables de las políticas deben decidir lo importante que es alcanzar algún nivel mínimo del comportamiento. En el caso de las donaciones benéficas, por ejemplo, no existe ninguna razón obvia por la que $1.000€$ deban ser el objetivo, por lo que puede ser mejor simplemente subvencionar a los individuos para que donen tanto como deseen. Sin embargo, con algunos comportamientos, el gobierno puede querer subvencionar algún nivel mínimo de provisión, pero no subvencionar una provisión particularmente generosa. Por ejemplo, con la subvención a la vivienda, podemos querer subvencionar algún nivel básico de consumo de vivienda, pero no viviendas particularmente grandes que no tengan beneficios externos evidentes.

\paragraph{Consideraciones de equidad }
Desde el punto de vista de la equidad vertical, los créditos fiscales son más equitativos que las deducciones. El valor de una deducción aumenta con el tipo impositivo de una persona (y, por tanto, con su renta), lo que hace que las deducciones sean regresivas (los importes de las deducciones son mayores, como proporción de la renta, para los contribuyentes de renta más alta). Los créditos, por otra parte, están disponibles por igual para todos los niveles de renta, de modo que son progresivos (los importes de los créditos son menores, como proporción de la renta, para los contribuyentes de renta más alta). Esta diferencia se acentúa aún más por la naturaleza de las deducciones detalladas de nuestro sistema fiscal. La tasa de utilización de deducciones detalladas entre los grupos de renta alta es muy elevada: el $96\%$ de los hogares con una renta superior a $200.000$ detallan sus deducciones. Sin embargo, solo el $6\%$ de quienes tienen rentas inferiores a $20.000$ euros lo hacen. Debido a que las deducciones normalmente solo están disponibles para quienes detallan sus deducciones, serán utilizadas en mayor medida por los grupos de renta más alta que sí las detallan, un hecho que reduce aún más su equidad.

\subsubsection{Gastos fiscales: ¿Un debate académico?}
El resultado final del conjunto existente de desviaciones respecto de la definición integral de renta de HAIG-SIMONS es un conjunto de gastos fiscales. En España, el conjunto de gastos fiscales (beneficios, deducciones, exenciones, créditos y tipos reducidos) suman un total de casi $80.000$ millones de euros (2020). O, de forma equivalente, casi el doble del déficit presupuestario observado para ese mismo año fiscal (2019).\footnote{%
    Para obtener más información: \href{https://www.funcas.es/articulos/los-gastos-fiscales-de-los-principales-impuestos-espanoles/}{\textit{Los gastos fiscales de los principales impuestos españoles}. FUNCAS}.

En Estados Unidos, por ejemplo, el mayor gasto fiscal individual es la exclusión de las primas del seguro médico proporcionado por el empleador de la renta imponible, que costó al gobierno federal $212,8$ mil millones de dólares en ingresos del impuesto sobre la renta no percibidos en 2014. Otros grandes gastos incluyen la exclusión de la renta neta imputada del alquiler (75,5) y los planes de aportación definida de los empleadores (79,7 mil millones de dólares), así como la deducción de los intereses hipotecarios (101,5 mil millones de dólares). La cantidad total de gastos fiscales del impuesto sobre la renta fue algo inferior al $40\%$ de los ingresos fiscales totales proyectados para 2014 y más de tres veces superior al déficit federal proyectado para 2014.
} 

Por tanto, el debate sobre la base imponible no es académico; los agujeros en la base imponible de HAIG-SIMONS permiten que grandes sumas de dinero escapen del Tesoro cada año.

\subsection{¿Cuál debe ser la unidad de imposición?}
Otro aspecto importante de la base imponible ha estado ausente de nuestro análisis del principio de HAIG-SIMONS: ¿cómo debería elegir el gobierno la unidad apropiada de imposición? Es decir, ¿cómo debería repartir el gobierno la carga fiscal entre los individuos que pertenecen a la misma familia? ¿Deberían aplicarse los impuestos sobre la renta total de la familia o únicamente sobre las rentas de los individuos?

\subsubsection{Trilema de la unidad de imposición: Impuesto al matrimonio}
Supongamos que usted fuera contratado por el gobierno federal para diseñar un sistema fiscal que tuviera tres objetivos: (i) \textbf{progresividad}, los tipos impositivos marginales aumentan a medida que aumentan las rentas familiares; (ii) \textbf{equidad horizontal entre familias}, las familias con rentas iguales pagarían impuestos iguales (situaciones objetivamente análogas son tratadas iguales); y, (iii) \textbf{equidad horizontal entre estados matrimoniales}, las cargas fiscales serían neutrales respecto al matrimonio, independientemente de si dos individuos deciden casarse (situaciones subjetivamente análogas son tratadas iguales). Todos ellos parecen objetivos que merece la pena perseguir. Sin embargo, existe un problema: es literalmente imposible alcanzar los tres objetivos al mismo tiempo. 

Considérese un sistema fiscal con un tipo impositivo del $10\%$ sobre la renta de hasta $20.000€$, un tipo marginal del $20\%$ hasta $80.000€$ y un tipo marginal del $30\%$ sobre cualquier renta superior a $80.000€$. Consideremos ahora dos parejas exitosas. La primera pareja, Carlos y Ángel, obtiene la mayor parte de su renta de Ángel, que gana $140.000€$ al año, mientras que Carlos gana solo $10.000€$. Jesús y María, por otra parte, contribuyen por igual a su pareja, ganando cada uno $75.000€$ al año. Ambas parejas tienen la misma renta total de $150.000€$. Tenemos dos opciones sobre cómo gravar a estas parejas:
\begin{la}
    \item \textbf{Base personal}. Si el gobierno grava a las parejas de forma individual, calcularía la carga fiscal de cada individuo y después sumaría las cargas para hallar la factura fiscal de la familia. Este enfoque daría lugar a una factura fiscal de $33.000€$ para Carlos y Ángel, pero de solo $26.000€$ para Jesús y María. Por tanto, este enfoque viola nuestra segunda condición: las familias con rentas iguales no están pagando impuestos iguales;
    \item \textbf{Base familiar}. La alternativa es gravar a las parejas sobre una base familiar. Bajo este sistema, el gobierno sumaría primero las rentas de los individuos para obtener la renta familiar y después calcularía un impuesto familiar basado en esa renta familiar. Este enfoque da lugar a facturas fiscales de $35.000€$ para ambas familias, de modo que se cumple la segunda condición. Sin embargo, un sistema de este tipo viola nuestra tercera condición, porque ambas parejas pagan ahora más impuestos como familias ($35.000€$) de lo que pagarían como individuos que viven juntos ($33.000€$ o $26.000€$). Es decir, ahora existe un impuesto sobre el matrimonio, un aumento de las cargas fiscales totales de dos individuos simplemente por casarse.
\end{la}

Estos problemas surgen debido a la primera condición, la progresividad. Si gravamos sobre una base individual, los tipos marginales crecientes significan que los individuos con rentas más altas pagan una proporción mayor de su renta en impuestos. Por tanto, las familias con una distribución más igualitaria de la renta pagarán menos impuestos, violando nuestra segunda condición. Los tipos marginales crecientes también significan que, cuando los individuos combinan sus rentas formando familias, una proporción mayor de la renta familiar total está sujeta al tipo impositivo marginal más alto. Por tanto, las familias pagarán más impuestos cuando estén casadas que cuando estén solteras, violando nuestra tercera condición.\footnote{%
    Volviendo a nuestro ejemplo, cuando gravamos sobre una base individual, Ángel paga el $23\%$ de su renta en impuestos, Carlos paga el $10\%$, y Jesús y María pagan cada uno el $17\%$. Sumando los impuestos individuales que debe cada familia, Carlos y Ángel pagan una proporción mayor de su renta total ($22\%$) que Jesús y María ($17\%$), violando la segunda condición. Si pasamos entonces a la imposición familiar, los $10.000€$ de ingresos de Carlos, que antes estaban gravados a un tipo bajo del $10\%$, ahora están gravados a un tipo mucho más alto del $30\%$. Bajo el impuesto individual, tanto Carlos como Ángel se beneficiaban de que los primeros $10.000€$ de renta estuvieran gravados únicamente al $10\%$, pero, con un impuesto familiar, Ángel agota ese tramo de $10.000€$ con sus primeros $10.000€$ de ingresos. Por tanto, los ingresos de Carlos se añaden por encima de la renta de Ángel, de modo que quedan dentro del tramo del tipo impositivo marginal más alto. Un fenómeno similar ocurre con Jesús y María: mientras que antes ambos se beneficiaban de una imposición baja sobre los primeros $10.000€$ de ingresos, ahora solo uno de ellos lo hace. Esto provoca el impuesto sobre el matrimonio, violando la tercera condición.
} La única manera de satisfacer tanto la segunda como la tercera condición sería tener un impuesto proporcional, lo que violaría la primera condición (progresividad).

¿Debería preocuparnos los impuestos sobre el matrimonio? Una razón para preocuparse es la cuestión de la equidad horizontal señalada anteriormente. La otra razón para preocuparse es que la sociedad podría querer realmente fomentar el matrimonio, no desalentarlo mediante impuestos sobre el matrimonio. Una vez más, si el gobierno utilizara un incentivo fiscal para fomentar el matrimonio, tendría que considerar los efectos marginales e inframarginales de este incentivo: ¿influyen realmente los impuestos en las decisiones de casarse, o una política de este tipo simplemente concede una gran desgravación fiscal a quienes se habrían casado de todos modos?

La respuesta parece ser, en gran medida, lo segundo; la mayor parte de la investigación sobre este tema sugiere que los impuestos ejercen poco efecto sobre la decisión de casarse. Por tanto, fomentar el matrimonio no es un argumento sólido a favor de los grandes costes en términos de ingresos que supondría eliminar los impuestos sobre el matrimonio impuestos a algunas familias (y aumentar las subvenciones al matrimonio para otras).

Por último, la razón por la que podría preocuparnos no es en absoluto el matrimonio, sino el elevado tipo impositivo marginal sobre los perceptores secundarios de renta.\footnote{%
    En nuestro ejemplo, pasar de un impuesto individual a un impuesto familiar ha elevado el tipo impositivo marginal sobre los ingresos de Carlos del $10 \to 30\%$. Un tipo impositivo tan elevado podría reducir la oferta de trabajo de los perceptores secundarios de renta, haciendo, por ejemplo, que Carlos no trabajara en absoluto.
} Una propuesta para abordar este problema (que es mucho menos costosa que aumentar la deducción para quienes presentan una declaración conjunta) consiste en introducir una deducción para el perceptor secundario de renta que permitiría a las familias deducir cierta cantidad de los ingresos de un perceptor secundario. Bajo una política de este tipo, por ejemplo, las parejas podrían deducir de la imposición los primeros $10.000€$ de ingresos de un perceptor secundario, lo que eliminaría los tipos impositivos más altos sobre el trabajo de Carlos bajo la imposición familiar en comparación con la imposición individual. Esta deducción reduciría enormemente tanto el impuesto sobre el matrimonio como la distorsión de la oferta de trabajo del perceptor secundario. Al mismo tiempo, sin embargo, reduciría la equidad vertical porque las familias con dos perceptores de renta tienden a tener más dinero que las familias con un solo perceptor.

\paragraph{¿Es ineludible el impuesto al matrimonio? Deducciones}
Obsérvese que la tercera condición se formuló en términos de tener neutralidad respecto al matrimonio, no en términos de no tener impuestos sobre el matrimonio. De hecho, podríamos tener un sistema sin impuestos sobre el matrimonio proporcionando deducciones muy grandes a las parejas casadas en relación con los contribuyentes solteros. Supongamos que complementáramos nuestro sistema fiscal con una deducción para las parejas casadas de 20.000 dólares al año, sin ninguna deducción para los solteros. En este caso, cada pareja tendría ahora una renta imponible de 130.000 dólares y, por tanto, pagaría 29.000 dólares en impuestos. Bajo este sistema, Carlos y Ángel tendrían una subvención al matrimonio (una reducción de su factura fiscal por casarse), mientras que Jesús y María tendrían un impuesto sobre el matrimonio. Incluso podríamos aumentar la deducción a 40.000 dólares, y entonces ambas familias tendrían una subvención al matrimonio.

La cuestión no es que el gobierno no pueda eliminar los impuestos sobre el matrimonio; puede hacerlo. La cuestión es que no existe ningún conjunto de deducciones que pudiéramos establecer que hiciera que el sistema de imposición basado en la familia fuera neutral respecto al matrimonio, en lugar de proporcionar subvenciones a algunos matrimonios y gravar a otros. Dadas las diferencias entre estas parejas y los demás objetivos de la política fiscal, no existe forma de diseñar un sistema que genere un impuesto/subvención al matrimonio igual a cero para ambas parejas. Esta es la esencia del problema al que se enfrentan los sistemas fiscales basados en la familia.

De las naciones industrializadas de la OCDE, $19$ gravan individualmente a los maridos y a las esposas (violando el objetivo de equidad horizontal entre estados matrimoniales), y cinco (Francia, Alemania, Luxemburgo, Portugal y Suiza) ofrecen subvenciones al matrimonio a prácticamente todas las parejas mediante la imposición familiar con división de la renta. Esto significa que la renta familiar se suma y después se divide por igual, bien entre el marido y la esposa (en Alemania, como en Estados Unidos entre 1948 y 1969) o bien entre todos los miembros de la familia, incluidos los hijos (en las otras cuatro naciones), violando no obstante la noción de neutralidad entre estados matrimoniales. Bajo la imposición familiar con tipos impositivos marginales crecientes, las familias con una renta distribuida de manera más igualitaria (Jesús y María) pagan menos impuestos que aquellas con rentas más desiguales (Carlos y Ángel). Por tanto, permitir que las familias dividan su renta por igual, junto con otras disposiciones generosas para las familias, conduce a bonificaciones por matrimonio en la mayoría de los casos en estas naciones. Solo tres naciones tienen un sistema puro de imposición familiar: Estados Unidos, Irlanda y Noruega.

















































































\clearpage
\section{Imposición directa: Periódicos}

La imposición sobre la renta tiene una serie de elementos comunes que debemos dejar claros.

• Se entiende como rentas aquellas en sentido amplio: rentas, beneficios o ganancias de capital efectivas o presuntas. 
• Su sujeto pasivo puede ser personas fisicas, los hogares, las sociedades y las instituciones sin fines de lucro.

El impuesto sobre la renta es un impuesto que ha sabido adaptarse a los cambiantes circunstancias de la estructura económica y a los diversos encuadres institucionales. Dichas características han permitido que el impuesto sobre la renta sea una figura central en la recaudación a lo largo de gran parte del S. XX.

La implantación oficial del primer impuesto sobre la renta se produjo en Gran Bretaña en 1799, como impuesto clave para sufragar gastos bélicos y mitigar tensiones sociales. De esta manera en su base se concebía como un instrumento de doble filo. Por un lado, un instrumento recaudador y, por otro, uno redistributivo. Es por ello que la imposición sobre la renta es más compleja en estructura y objetivos.

Así pues, el impuesto a la renta muestra una variada estructura de tipos y bases, desde los más complejos modelos integrales hasta los más simples de tasa uniforme, y exhibe una gran diversidad de exenciones e incentivos. Esta estructura es representativa del dificil equilibrio que se debe alcanzar entre los objetivos tributarios en la aplicación de la imposición sobre la renta.

Así la suficiencia, eficiencia y equidad en la imposición sobre la renta actuaran en sentido contrario y será dificil alcanzar un equilibrio dentro del mismo tributo, que además suele funcionar como pilar central del sistema tributario.

Por este motivo y por la dificultad de alcanzar un diseño adecuado en el impuesto sobre la renta, este ha sufrido diversos vaivenes desde ser abolido a reinstaurado en diversos países a lo largo de los siglos XIX y XX.

·'Las conti11uas guerras de consolidaciá11 de los Estados nacionales y de expansión imperial, la presiá11 de los movimientos políticos opuestos al capitulismo industrial. así como los cambios sociales )' tecnolóKicos, hicieron necesario consolidar y tramformar el tributo.'' Podemos señalar que la imposición sobre la renta ha sufrido cuatro grandes cambios:

La introducción de un impuesto técnico y no discrecional. que incluía una estructura progresiva y un sistema de gestión administrativo. tanto en la Ley de presupuesto británica de 1909 como en la Ley de impuesto a la renta federal de los Estados Unidos. de junio de 1913.

• La masificación del impuesto se dio por primera vez en EE.UU. en la admin istraciónRooscvelt. alrededor de la segunda guerra mundial. en para lelo con la ampliación de la participación democrática y de la acción del Estado mediante programas de bienestar social (welfare state). Sustentado por el éxito del New Deal y la reconstrucción de Europa occidental, así como por la interpretación sesgada del keynesianismo. el crec imiento de los gastos públicos generó voracidad fiscal. Esta presión sometió al impuesto a una espiral al alza de los tipos y a la vez a un juego de complt:ias exenciones y tratamientos especiales que desnatura lizaron su estructura.

• La tercera fase fue la "contrarreforma·· iniciada por las reformas de la primera ministra Thatcher. ·reforzadas por las de la administración Reagan. en los a1ios 1980, que devolvieron la composieión del tributo a sus orígenes mediante una fuerte reducción de tasas y ampliación de bases, sin afectar significativamente ni el rendimiento fiscal ni la carga sobre los factores.

• La cuarta reformulación, a inicios de los años 1990, basadas en el modelo dual de los países nórdicos y el impuesto de tasa uniforme (tlat tax) de las economías en transición (del socialismo al capitalismo) del ex bloque soviético adaptaron el impuesto a la renta para hacer frente a la competencia internacional por el ahorro y la inversión, manteniendo (parcialmente) su carácter progresivo.



\subsection{Renta de las personas físicas: IRPF}
Se trata de aquel conjunto de impuestos que recaen sobre las personas físicas. Se recauda de manera periódica, su hecho imponible es una manifestación de riqueza asociado a una renta y tiene en cuenta las características particulares del contribuyente.

La estructura de este tipo de tributos es de especial relevancia pues presenta mucha heterogeneidad. Puede ser, por ejemplo, un tributo sencillo de estrecha base y altos tipos como en Francia, o un impuesto dual como en los países nórdicos. También puede ser un impuesto de base estrecha y tipo bajo como en Estados Unidos.

IV. ti. Estructura del tributo.
La estructura general de un tributo sobre la renta. independientemente de que nación o sistema lo aplique es la siguiente:

+ RENTA BRUTA
- Gastos deducibles.
= RENTA NETA
- Reducciones de la renta
- Bonificaciones
= BASE IMPONIBLE
- Reducciones personales
= BASE LIQU IDABLE
X TIPO DE GRAVAMEN
= CUOTA INTEGRA
- Deducciones.
= CUOTA LIQUIDA
- Retenciones y pagos a cuenta.
= RESULTADO

Hasta este momento hemos hablado de la determ inación de la base imponible pues es lo más importante de este tipo de tributos pues lo define y establece su naturaleza. Sin embargo, hay otra serie de elementos relevantes que son trascendentales a la hora de definir este tipo de impuestos y que repasaremos ahora.

\subsubsection{Desarrollo histórico: ¿Cuándo se ha desarrollado el impuesto sobre la renta?}
Debemos partir de una idea de tributo primitiva aun hoy presente en países en vías de desarrollo. un impuesto no técnico y arcaico en el que los tipos son excesivamente elevados y el nivel de fraude y elusión también. De esta manera se intenta captar algo de recaudación en una base imponible impredecible e inestable. Este tipo de estructuras hace que los primeros cambios fuesen una revolución y vayan asociados al desarrollo de un sistema tributario.

El primer sistema técnico buscaba establecer una base imponible conocida y estable. sobre la que los agentes tributasen. En este sentido la renta que tributaba al principio era la del trabajo. Así pues, la única renta que se tenía en cuenta era la imposición sobre el trabajo y el resto de rentas (mobiliarias. i nmobi liarías o de actividades económicas) quedaban fuera del marco de un tributo sobre la ren ta. Hay que recordDr que el desarrollo de los tributos Vél en paralelo al desarrollo del gasto, pues estos surgen para atender una necesidad. En el caso de los tributos primitivos se buscaba atender a unas necesidades puntuales (Guerras) o de mantenimiento de un Estado de tamaño pequeño y poco organizado. Con el desarrollo de los estados se hizo necesario un im puesto más estructurado y ello ocurrió en la revolución industrial. donde el trabajo estaba más concentrado y la tributación por primera vez podía re alizarse de manera sistemática, sobre los trabajadores. Este tipo de tributación tuvo como contraprestación el inicio del estado del bienestar y por ello supuso en términos de supervivencia un éxito. La concentración de los trabajadores y el diseño a priori a través de retenciones hacían que el trabajador cayese en la ilusión monetaria y la tributación fuese relativamente indolora para el contribuyente.

El siguiente paso fue la tributación de otras fu entes de renta, que podemos diferenciar en dos tipos:

• Las de las actividades económicas y • las provenientes del patrimonio (rendimientos mobiliarios, inmobiliarios y ganancias patrimoniales). A priori estas rentas no podrían estar sujetas a retención por su natural. y por ello al requerir declaración eran más complejas de gravar. A su vez el diseño del tributo con respecto a las rentas provenientes de actividades económicas era de especial trascendencia pues podría provocar el traspaso de tributación de los hogares a las empresas y viceversas. Por todo dio pasamos a analizar los principales elementos definidores de estos tipos de tributos.

IV.I . Definición d e l a base i m ponible.
La imposición sobre la renta y la riqueza parten de la clave de una definición de base imponible genérica y de carácter casi universal. que es:
" Toda renta y toda riq11e::a ohtenida por los hogares ".
La definición técnica de la misma depende de muchas decisiones de política económica. pero de manera general podemos decir que:
• Base Imponible de l a Renta = Rentas del trabajo + Beneficios empresariales extraordinarios + Rentas del capital.
• Base Imponible de l a Riqueza = Activos - Pasivos. 
U n elemento clave de los impuestos de tipo directo es que tienen en cuenta las característ icas de los contribuyentes a lo largo de todo el diseño del tributo. Es por ello que unos elementos importantes a resaltar son todos aquellos que puedan afectar a la determinación de la base imponible.

\subsubsection{¿Cuáles son las desviaciones contempladas? Gastos fiscales}

Elementos de política económica que afectan a la determinación de la Base Imponible de la Renta:



• Reducciones.
Son reducciones a la renta neta que están totalmente desvinculadas de ser gastos y que se i ntroducen por motivos técnico tributarios o por motivos de política económica. Por ejemplo, en el caso de España todo contribuyente tiene derecho a una reducc ión de 2000€ sobre su renta del trabajo, sea este multimillonario o no. De hecho, los defensores de la renta básica. sei'ialan que esta ya existe en esta reducción y que trasladarla del lado del ingreso al gasto pennite m ejorar su progresividad y capacidad distributiva, pero que el impacto presupuestario ya está asumido en par1e. En otros países se introducen reducciones por residencia o por sector de trabajo. Corno en China donde los funcionarios tienen grandes ventajas fiscales.

• Bonificaciones.
Las bonificaciones son una forma de reducción proporcional que se aplica a la renta, sea esta bruta o neta, según se decida por motivos meramente de política tributaria. Unas bonificaciones es no incluir un porcentaje de la renta en el cómputo de la base. Esta es la principal diferencia con la reducción. Por ejemplo, es posible que ciertas fuentes de rema sean bonificada en gran parte porque se considera que provienen de una fuente que debe ser premiada: alquiler para vivienda de larga duración (bonificado al 60\%) vs alquiler vacacional (sin bonificación). Otros países. generalmente aquellos que producen una materia prima, para compensar el efecto de aranceles o cuotas a la exportación establecen bonificaciones a todos los que trabajan en el sector e:-: portador.



• Reducciones personales.

Este tipo de reducciones es de gran trascen dencia pues son aquellas que pemliten una adaptación a las características propias del contribuyente a la vez que establecen un primer paso de progresividad. Entre estas reducciones podemos establecer una distinción:
o Aquellas que no tienen en cuenta las características del contribuyente y se aplican a todos los contribuyentes independientemente de sus características y que pueden variar o no en virtud de algún criterio económico. Como puede ser el mínimo personal que se aplica a todos los contribuyentes.
o Aquellas que tienen en cuenta las características del contribuyente. por ejemplo, el número de hijos o dependientes a cargo o si estos tienen alguna minusvalía o no. Esta reducción aplica por primera vez en el impuesto los criterios principios de justicia tributaria pues la reducción sobre el total de base para una persona con una renta millonaria será menor que la de una persona que no lo es, tanto en media como en el margen. Así pues. por el mero hecho de existir estas reducciones ya existe progresividad. A su vez el hecho ele que estas tengan en cuenta las condiciones particulares de cada contribuyente aplican el principio de equidad ho1izontal al tributo.

• Tipo de gravamen.

El segundo elemento determinante es el tipo de gravamen y las características de este. En este punto existen infinidad de modelos que podemos resumir en dos:

• Impuestos sobre la renta de tipo único. Que confían la progresividad del impuesto. si este lo tiene, a la determinación de la base. Como ha ocurrido en algunos países del este de Europa.

• Impuestos sobre la renta de tipo mLiltiple que a su vez divididos en dos:
o De tipo múltiple y que aplica un tipo difer ente a cada tipo de renta. Este tipo de
estructura se aplica generalmente a los tributos sobre la renta de tipo dual, que
veremos más adelante.
o De tipo múltiple y que aplica una escala a la base imponible. Este tipo de estructuras fí an la progresividad al tipo donde la base imponible se divide y dependiendo del tamaño de esta tributará por un numero de escalas determinadas. En este caso es muy importante recordar la diferencia entre el tipo marginal, último tipo de la escala a la que se tributa, y tipo medio, que es el tipo efectivo de la base imponible después de pasar por todos los niveles de la escala. De manera que el tipo marginal será siempre superior al medio. El número de tramos de la escala y su definición son un elemento importante de política tributaria. Véase el ejemplo.

\textbf{ESCALA A Y B: IMAGEN ICEX-CECO}

• Deducciones.
Las deducciones son medidas de pol ítica económica que premian determinadas actuaciones
o circunstancias personales del contribuyente y que implican una reducción directa de la carga
tributaria pues se reducen directamente de la cuota. En algunos casos este tipo de deducciones
se pueden convenir en monetizables ex ante. como si fuesen una subvención. y se las suele
denominar impuestos negativos. La peculiaridad de estas con respecto a la subvención es que
sólo se pueden beneficiar los contribuyentes.
Así pues. caben dos tipos ele deducción:
o Por una circunstancia personal. Deducción por nacimiento de hijos.
o Por una actuación fiscalmente fovorecida. Deducción por inversión en vivienda
habitual.

Las deducciones por lo general son consideradas beneficios fiscales. en el sentido que suponen una menor recaudación a la esperada debido a una decisión de política económica.

• Retenciones y pagos a cuenta.
Conceptualmente se tratan de una liquidación anticipada de la deuda tributaria previa al devengo. El devengo de este tipo de tributos suele ser al final del año económico (en el mundo occidental es el 31 de diciembre). Sin embargo. con antelación y en virtud de información t1ibutaria cierla se pueden realizar liquidaciones anticipadas para suavizar el impacto en la liquidez y solvencia de los contribuyentes, y como no. dotar de fondos y liquidez al Estado antes del devengo del tributo. La existencia de un sistema de retención o pago a cuenta y de una liquidación anual suele ser indicador de un sistema tributario moderno. Los sistemas menos modernos. se basa solamente en un sistema de retención sin liquidación anual. lo cual genera una perfecta ilusión monetaria sobre los contribuyentes y reduce en gran medida la imagen de presión fiscal. Pese a ello, lo que los contribuyentes no son conscientes es que este hecho puede implicar de facto una carga tributaria mayor, al no existir devoluciones. y una mayor desigualdad. por la naturaleza de las rentas sujetas a retención.

Otra razón por la que existen las retenciones es por la gestión de liquidez de las autoridades públicas. La existencia de retenciones permite que no se tenga que esperar al devengo del tributo para que se ingrese en las arcas públicas el dinero. El ingreso periódico de tributos supone un ingreso recmTente y esencial para las autoridades pues mejora su posición de liquidez y permite el realizar desembolsos con la misma periodicidad, como nóminas o pensiones. Así pues. la existencia de este tipo de estructuras tributarias es imprescindible para el diseño de un Estado.

Las rentas sujetas a retención tienen la característica de ser periódicas y determinables en el momento. Estas rentas suelen provenir de dos fuentes:
o Trabajo. Una proporción de la renta obtenida por el trabajador.
o Actividades económicas. Asociadas a una presentación de cuentas u otros tributos.

Otro tipo de rentas (inmobiliarias. ganancias patrimoniales) no presentan retención sea por su falta de regularidad o por motivos de técnica tributaria. Si bien desde el desarrollo de las tecnologías de la inform ación ciertas rentas que se sujetaban difícilmente a retención se sujetan con mayor facilidad. como los rendimientos mobiliarios. Esto ocurre asi generalmente en los países europeos. Las retenciones y los pagos a cuenta como parte del tributo pueden y deben en principio tener en cuenta las condiciones particulares de cada individuo, si bien sólo tienen en cuenta una óptica del mismo, la del pagador de una renta. En este sentido si un contribuyente tiene más de un pagador, la retención de cada uno se hará como si esa fuese la única renta. lo que en térm inos generales de progresiviclad supondrá el aplicar un tipo inferior al que debería apli carse si se tuviese en c uenta la renta universal.


I V.1 .1. 1. ])os ti pos dual y plano

En un primer momento los tributos sobre la renta eran de tipo plano o único, como en algunos países de E uropa del este o en desaiTollo. Es decir. todas las rentas tributaban al mismo tipo o escala de tipos. Sin diferenciar su origen o potenciales efectos sustitución. La existencia de una escala única progresiva supone que todas las rentas tributan por la misma escala y los contribuyentes con tipos marginales altos tienen incentivos para un potencial problema de trasvase hacia otras f iguras como un impuesto sobre la renta y riqueza de las empresas, pues si los tipos o escalas de estas estructuras son menores, todos los empleados propios ( o autónomos) tenderán a desviar su tributación a la tributación empresarial.

Es por ello que existe e l sistema tributario dual. Este sistema reproduce la carga tributaria de las empresas para aquellas rentas que tiene naturaleza similar pero que obtienen las personas físicas. Es decir. que se da un doble sistema tributario para las personas físicas. Una faceta tradicional, que suele cubrir las rentas del trabajo, de los rendimientos inmobiliarios o las · ganancias patrimoniales, que pasa por una escala de IRPF tradicional. sin tener grandes reducciones. Por otro lado, otra que afecta principalmente a las actividades económicas o rendimientos inmobil iarios, que utiliza la política de deducción ele gastos y que suele ir gravada a un tipo menor y próximo a la de la tributación de las empresas. Estos sistemas duales se introdujeron por primera vez en los países nórdicos.

IV. IV. La progresividad y la rcgresividad en el tributo.

La progresividad o regresividad del tributo vendrá determ inada en este tributo por el diseño del m ismo, pese· a que con general idad estamos hablando de tributos progresivos al tener en cuenta las condiciones particulares de cada contribuyente. Las principales características que le otorgan progresividad, que a sensu contrario de no existir generan regresividad, son:
• Escala de gravamen progresiva.
• Existencia de mínimos exentos y familiares.
• Existencia de reducciones en la base en función de las fuentes de renta y el contribuyente.
• Existencia de deducciones que tienen exclusivamente las características del contribuyente.

La no existencia de estas características llevaría a que se concibiese el impuesto como uno regresivo, próximo a un impuesto de suma fija.

Como no podría ser de otra forma la existenci a de estas medidas genera mayores distorsiones en el sentido económico. Un ejemplo clarificador, es la deducción por vivienda habitual, la existencia de este tipo de deducciones introduce un sesgo pro compra de vivienda que afecta al equilibrio que alcanza una econom ía y puede introducir. como en el caso de España un desequilibrio que será dificilmente ajustable.

\subsubsection{¿Cómo se determina la renta sujeta a gravamen? }
• Determinación de la renta neta. 
La renta por la que tributa el contribuyente, es una renta Neta, no Bruta. Por lo tanto, hay una serie de gastos que son fiscalmente deducibles. La determinación de qué gastos son fiscalmente deducibles o no depende del legislador de cada sistema tributario pero una norma más o menos uni versal es la siguiente:
•·sera deducible todo gasto en el c¡ue se haya incurrido para generar la renta·•. 
Definición genérica que hace que dependiendo del sistema se pueda incluir un gasto u otro, por ejemplo, en el caso español se deducen los gastos de seguridad social o determinados gastos sindicales de la renta del trabajo, mientras que en el caso de los EE. UU. son deducibles muchos más gastos como el seguro médico privado. educación. alqui ler etc . . .

\subsection{Renta de las personas jurídicas: Impuesto de Sociedades (IS)}
Una sociedad es una forma de organización de la actividad económica en la que la propiedad suele estar representada por certificados de acciones transferibles. Las sociedades son entidades jurídicas independientes y, como tales, personas jurídicas artificiales. Una sociedad puede celebrar contratos, poseer bienes, contraer deudas, demandar y ser demandada. Y, al igual que cualquier otra persona, una sociedad debe pagar impuestos sobre su renta. ¿O no? El presidente Ronald Reagan una vez mencionó: \textit{Probablemente me arrepentiré de haber dicho esto, pero ¿cuándo vamos a tener el valor de señalar que, en nuestra estructura fiscal, el impuesto sobre sociedades es muy difícil de justificar?} 

¿Tiene sentido disponer de un sistema fiscal especial para las sociedades? Sin duda desde un punto de vista jurídico, las sociedades son personas. Pero, desde un punto de vista económico saemos que esta noción no tiene ningún sentido. Como sabemos: solo las personas reales pueden pagar un impuesto. Si es así, ¿por qué debería estar sujeta la actividad societaria a un impuesto especial? ¿Por qué no gravar simplemente las rentas de los propietarios de la sociedad mediante el impuesto sobre la renta personal?

\subsubsection{¿Por qué gravar la renta de las personas jurídica?}
En primer lugar, contrariamente a la opinión que acabamos de exponer, las sociedades (especialmente las muy grandes) son realmente entidades distintas. Las grandes sociedades tienen miles de accionistas, y los directivos de dichas sociedades están controlados solo de manera muy laxa, si es que lo están en absoluto, por los accionistas-propietarios. La mayoría de los economistas estarían sin duda de acuerdo en que la propiedad y el control están separados en las grandes sociedades, y esto crea importantes problemas para comprender cómo funcionan las sociedades. No obstante, de ello no se sigue que la sociedad deba ser gravada como una entidad separada.

Una segunda justificación de la imposición sobre las sociedades es que la sociedad recibe una serie de privilegios especiales de la sociedad, el más importante de los cuales es la responsabilidad limitada de los accionistas: el impuesto sobre sociedades puede considerarse una \textbf{tarifa de usuario por este beneficio}. Sin embargo, no hay ninguna razón para creer que los ingresos pagados se aproximen a los beneficios recibidos. En cualquier caso, ¿por qué deberíamos considerar que las leyes que permiten una forma eficiente de que los individuos agrupen su capital constituyen un beneficio que requiere un pago? Las leyes que permiten otros tipos de contratos no se consideran de esta manera.

La existencia de un tributo sobre las empresas está íntimamente asociado a la idea de una personalidad jurídica y de una responsabilidad limitada. Si repasamos el periodo histórico hasta que no encontramos ambas c a racterísticas no podemos hablar de un impuesto sobre la renta de las empresas moderno, lo cua l sería a mediados del siglo XX en Europa y a principios e;:n EE. U U.

Las figuras anteriores eran más próximas a extracciones coactivas de renta a las empresas existentes, sin ningún viso tributario sino más bien contributivo (usted debe aportar una cantidad X, sin mayor justificación). Algunos pa íses, sin embargo, como Alemania, mantiene estructurasde tributación arcaicas, basadas en los ingresos de las empresas y su sector. Esto sería una suerte de impuesto a la actividad económica de la empresa, pero alejado de un tributo moderno. La existencia de un tributo moderno está asociada a su vez a la existencia de una contabilidad regulada y de la existencia de determinados controles externos a la empresa. Esta cara cterísticaes la principal diferencia y bastante considerable entre la tributación de las empresas y de las personas. La existencia de una contabilidad y de unos criterios formales aleja a este tipo de figuras de la imagen del tributo sobre las rentas de las personas físicas.

\paragraph{Beneficio retenido: Integridad del impuesto sobre la renta personal}
Sin embargo, el más sonado (y acertado) de todos las justificaciones es que el impuesto sobre sociedades protege la integridad del impuesto sobre la renta personal. 

Supongamos que la parte de Ana de los beneficios de una sociedad durante un determinado año es de $10.000€$. Según la convención estándar de los economistas para definir la renta, estos $10.000€$ son renta tanto si el dinero resulta ser retenido por la sociedad como si se paga a Ana. Si los $10.000€$ se pagan, se gravan en una cuantía que depende de su tipo impositivo sobre la renta personal. En ausencia de un impuesto sobre sociedades, los $10.000€$ no generan ninguna obligación tributaria si son retenidos por la sociedad. 

Por tanto, a menos que se grave la renta de las sociedades, Ana puede reducir su obligación tributaria acumulando renta dentro de la sociedad. Por supuesto, el dinero será gravado cuando finalmente se pague, pero, mientras tanto, los $10.000€$ completos crecen al tipo de interés antes de impuestos: \textbf{los impuestos diferidos son impuestos ahorrados}.

De esta forma, los impuestos que paga y debe pagar cualquier sociedad son:

\eqblock{
\begin{aligned}
    T = ([\text{Ingresos} − \text{Gastos}] \cdot t) − \text{Créditos fiscales}
\end{aligned}
}{}

Como podemos suponer, los ingresos son lo que la empresa obtiene por la venta de bienes y servicios en el mercado. Los gastos, por su parte, serán los característicos de la propia actividad económica. Y, por último, los créditos fiscales será aquellos gastos fiscales estipulados por la legislació tributaria. 


\subsubsection{¿Cuáles son las desviaciones contempladas? Gastos fiscales}

El primero son los costes de flujo de caja de desarrollar la actividad empresarial. 


Estos costes comprenden todos los gastos realizados durante el último año en servicios o bienes. 


El segundo componente son los pagos de intereses, es decir, los pagos realizados a quienes prestan dinero a la empresa. 

\paragraph{Tratamiento de la financiación: Intereses deducibles}
Otro elemento importante es la deducibilidad de gastos financieros. Tradicionalmente se considera un gasto el interés pagado por la existencia de recursos ajenos. Así pues. existe un sesgo en favor de este tipo de endeuda.miento. Para eliminar dicho sesgo pro deuda se pueden realizar dos tipos de actuaciones.
• La limitación en los gastos a aplicar como un ajuste al resultado contable o
• El establecimiento de una reducción por financiación propia (ACE. allowance of capital equity). Es decir. se establece un porcentaje de reducción como gasto ante ampliaciones de capital y finalización propia. Este tipo de reducciones son un elemento también determinante a la hora de fijar la base imponible. Que se ha incorporado del sistema tributario de EE.UU. recientemente a Europa.

Cuando las sociedades piden prestado, los pagos de intereses a los prestamistas se excluyen de la renta imponible. Una vez más, la justificación es que los costes empresariales deberían ser deducibles. Sin embargo, cuando las empresas financian sus actividades mediante la emisión de acciones, los dividendos pagados a los accionistas no son deducibles de los beneficios de la sociedad.

Tanto los costes de flujo de caja como los pagos de intereses son deducibles de los beneficios de la empresa durante el período en que se producen; la cantidad gastada cada año en estas categorías se resta de los beneficios de ese mismo año al calcular la carga tributaria de la empresa. 

Como sabemos, las empresas tienen tres opciones principales para financiar la inversión: pueden utilizar beneficios retenidos, pueden aumentar su deuda pidiendo prestado o pueden emitir capital (participaciones de propiedad). Además, si una empresa elige la vía del capital, tiene dos opciones sobre cómo devolver los beneficios a los inversores: puede repartir dividendos o puede reinvertir los fondos e intentar aumentar el valor de las acciones, dando lugar a ganancias de capital. Por ahora, dejemos de lado el caso menos interesante de financiar la inversión mediante beneficios retenidos y centrémonos en el impacto de los impuestos sobre la financiación de la inversión mediante nuevo capital.


Supongamos que una empresa necesita 10 dólares para una inversión que generará 1 dólar de renta societaria cada año. La empresa quiere financiar esta inversión mediante deuda o capital y, en cualquiera de los dos casos, devolver ese dólar al inversor.  Gráficamente:

\imagenfit{Fig1.png}{Financiación: Incidencia conductual derivada de la deducibilidad de los intereses}{GRUBER}

El impacto de los impuestos sobre la decisión de la empresa acerca de cómo financiar esa inversión (ignorando una vez más, por ahora, la opción de financiarla mediante beneficios retenidos). En el primer nodo, la empresa decide entre asumir deuda o emitir capital para financiar la inversión de 10 dólares. Si la empresa se financia con deuda (la rama superior), pagará el dólar en intereses a los tenedores de bonos. En este caso, la empresa no pagará impuestos sobre el dólar de renta societaria porque los pagos de intereses son deducibles de la renta de la empresa a efectos del impuesto sobre sociedades: el dólar de renta societaria se compensa con 1 dólar de pagos de intereses al calcular los impuestos. Por tanto, el tenedor de bonos que ha prestado dinero a la empresa recibe el dólar completo, pero después tiene que pagar impuestos personales sobre ese dólar al tipo del impuesto sobre la renta personal aplicable a los intereses, t_{int}, de modo que el tenedor de bonos recibe 1\$ \times (1-t_{int}).

Alternativamente, la empresa puede emitir capital para financiar la inversión. En este caso, cuando la empresa obtiene el dólar de la inversión, no tendrá el pago de intereses que proporcione una deducción compensatoria a efectos de la imposición sobre sociedades, por lo que tendrá que pagar impuestos sobre sociedades sobre ese dólar. Por tanto, solo 1\$ \times (1-t_c), donde t_c es el tipo del impuesto sobre sociedades, llegará a los accionistas.

Si la empresa utiliza financiación mediante capital, tiene que tomar una decisión adicional: ¿cómo debería hacer llegar el dinero a los accionistas? Con la deuda, no hay elección: el dólar simplemente se paga como interés. Con el capital, la empresa tiene la opción de pagar un dividendo o retener los beneficios y reinvertirlos. Si paga el dólar como dividendo, entonces el accionista tiene que pagar impuestos sobre dividendos sobre ese dólar. Por tanto, el accionista recibe finalmente 1\$ \times (1-t_c)\times(1-t_{div}), donde t_{div} es el tipo impositivo personal sobre la renta procedente de dividendos. Esta característica se denomina a menudo doble imposición de los dividendos, ya que los dividendos son gravados tanto por el sistema del impuesto sobre sociedades como por el sistema del impuesto personal.

Si, en cambio, la empresa reinvierte el dólar, el accionista posiblemente obtendrá una ganancia de capital cuando venda la acción. Supongamos que cada dólar de reinversión aumenta el precio de la acción en 1 dólar. En este caso, el accionista recibe 1\$ \times (1-t_c)\times(1-t_{cg}), donde t_{cg} es el tipo impositivo efectivo sobre las ganancias de capital. El tipo impositivo efectivo sobre las ganancias de capital es inferior al tipo impositivo legal sobre las ganancias de capital para la mayoría de los inversores debido al conjunto de beneficios de los que disfrutan las ganancias de capital que analizamos en el Capítulo 23, incluida la imposición en el momento de la realización en lugar de en el momento del devengo y la actualización de la base fiscal en el momento del fallecimiento. El tipo impositivo efectivo sobre las ganancias de capital también ha sido tradicionalmente muy inferior al tipo impositivo sobre los dividendos. Desde 2004, los tipos legales sobre las ganancias de capital y los dividendos son iguales, del 15 %, pero el tipo impositivo efectivo sobre las ganancias de capital sigue siendo inferior al tipo impositivo sobre los dividendos.

Como debería quedar claro a partir de la Figura 24-7, la estructura del impuesto sobre sociedades puede desempeñar un papel fundamental en las decisiones de una empresa sobre cómo financiar sus inversiones. La financiación mediante deuda evita la imposición societaria de los beneficios, pero no permite a los individuos aprovechar los bajos tipos impositivos sobre las ganancias de capital. La financiación mediante capital, y el pago de dividendos, parece ser la vía menos eficiente fiscalmente de todas. Por tanto, esta figura plantea dos preguntas sobre la financiación empresarial.


\begin{specialblock}{\textbf{NOTA.-} ¿Por qué no se financian únicamente deuda?}
    Una empresa tiene dos medios para financiar una inversión. Si aumenta su deuda, los intereses que paga para financiar su inversión están exentos de la imposición sobre sociedades. Si aumenta su capital, paga el impuesto sobre sociedades. Entonces, ¿por qué las empresas no financian simplemente toda su inversión mediante deuda y evitan así la imposición sobre sociedades?

    Se han dado varias respuestas específicas a esta pregunta, pero todas se reducen a la distinción fundamental entre deuda y capital: la deuda exige pagos fijos, mientras que el capital no. Una empresa puede elegir si paga dividendos sobre el capital, pero no puede elegir si realiza sus pagos de intereses: debe realizarlos. Si la empresa no realiza sus pagos de intereses, incumple sus deudas y puede verse obligada a declararse en quiebra. En una quiebra, todos los individuos que han invertido en una empresa pueden presentar reclamaciones sobre los activos que le quedan. Los tenedores de bonos son los principales reclamantes; son los primeros en la fila para reclamar cualesquiera activos que puedan quedar. Los titulares de capital, por el contrario, solo tienen derecho a los recursos que queden después de que los tenedores de bonos hayan recibido lo que les corresponde.
    
    Imagine que pudiera elegir el sistema de calificación de este curso y que tuviera dos opciones. Bajo una opción, realiza una hoja de problemas cada semana y, si alguna vez suspende una hoja de problemas, suspende el curso. Bajo otra opción, realiza una hoja de problemas cada semana y es la nota media de todas sus hojas de problemas la que determina su nota final. La mayoría de nosotros elegiríamos el segundo sistema de calificación debido a su mayor flexibilidad.
    
    Del mismo modo, la flexibilidad adicional del capital puede elevar su valor lo suficiente como para compensar su desventaja fiscal. La financiación mediante capital puede proporcionar una «zona de amortiguación» frente al riesgo de quiebra; con capital, si las cosas van mal, la empresa puede simplemente no pagar un dividendo en lugar de verse obligada a declararse en quiebra. Los directivos en particular, que a menudo se preocupan más por reducir el riesgo cotidiano de perder su empleo que por maximizar el rendimiento para los propietarios distantes de una empresa, podrían preferir el capital a la deuda.
    
    Esta preferencia por el capital, a pesar de sus desventajas fiscales, se ve reforzada por un tipo particular de problema de agencia dentro de la empresa: el conflicto de intereses entre los tenedores de deuda y los titulares de capital. Los tenedores de deuda son partes que obtienen un rendimiento fijo (pago de intereses) independientemente de lo bien que le vaya a la empresa, siempre que esta no quiebre. Los titulares de capital, por el contrario, reciben un rendimiento que está vinculado al desempeño de la empresa: la parte de los beneficios que se reinvierte en la empresa aumenta de valor a medida que aumenta el valor de la empresa. En la Figura 24-7, nosotros  supusimos que cada dólar reinvertido en la empresa generaba únicamente un dólar de valor. En realidad, cada dólar de inversión obtiene una cantidad incierta. Los titulares de capital sufren cuando una inversión empresarial obtiene malos resultados, pero se benefician cuando obtiene buenos resultados.

    Para los titulares de capital, sin embargo, existe un límite inferior a cuánto pueden perder cuando a la empresa le va mal: todo lo que pueden perder es su inversión inicial si la empresa quiebra. Por otra parte, no existe ningún límite superior a cuánto pueden ganar cuando a la empresa le va bien. Por tanto, cuando los titulares de capital poseen solo una pequeña proporción de una empresa, querrán asumir riesgos excesivos porque obtienen un rendimiento mayor pero soportan relativamente poco riesgo; dado que poseen una pequeña parte de la empresa, no tienen tanto dinero que perder. La posibilidad de asumir riesgos excesivos es un problema porque los titulares de capital toman las decisiones sobre a quién contratar y, potencialmente, qué proyectos elegir.
    
    Considérese una empresa que actualmente tiene un capital de 1 millón de dólares y una deuda de 5 millones de dólares, como se muestra en el panel superior de la Tabla 24-1, lo que da un valor total de la empresa de 6 millones de dólares. El pago de intereses de la empresa sobre la deuda en cada período es de 500.000 dólares. La empresa tiene actualmente unos beneficios de 600.000 dólares al año, de modo que cada año puede pagar los intereses que debe a sus tenedores de deuda y pagar los 100.000 dólares restantes de beneficios como dividendo a sus titulares de capital.
    
    Supongamos que, además de este flujo constante de beneficios, se ofrece a la empresa un proyecto que tiene un 50 % de probabilidades de generar 3 millones de dólares y un 50 % de probabilidades de perder 6 millones de dólares. Si la inversión fracasa, la empresa quebraría porque tendría que vender sus 6 millones de dólares en activos (5 millones de dólares de deuda y 1 millón de dólares de capital) para pagar los 6 millones de dólares de deuda derivados de la pérdida. Desde la perspectiva de la empresa en su conjunto, este no es un buen riesgo que asumir. El rendimiento esperado de esta inversión, las probabilidades de ganar multiplicadas por el valor si la empresa gana, más las probabilidades de perder multiplicadas por el valor si la empresa pierde, es
    
    (0,5\times3\text{ millones de dólares})+(0,5\times-6\text{ millones de dólares})
    =-1,5\text{ millones de dólares},
    
    de modo que a la empresa le conviene más no realizar la inversión porque, en promedio, pierde dinero.
    
    Pensemos ahora en el proyecto desde la perspectiva de los titulares de capital. Si esta inversión tiene éxito, reciben los 3 millones de dólares completos porque obtienen todos los beneficios derivados del mejor desempeño de la empresa. Si la inversión fracasa y la empresa quiebra, pierden el millón de dólares que han invertido, además de su pago de dividendos de 100.000 dólares en cada período futuro. A un tipo de interés del 10 %, el valor actual descontado de este flujo de pagos de dividendos es de 1 millón de dólares, de modo que los titulares de capital pierden un total de 2 millones de dólares si la empresa quiebra.²⁰ Desde la perspectiva de los titulares de capital, existe un 50 % de probabilidades de que pierdan 2 millones de dólares y un 50 % de probabilidades de que ganen 3 millones de dólares, lo que da un valor esperado para los titulares de capital de 500.000 dólares:
    
    [0,5\times(-2\text{ millones de dólares})]+[0,5\times3\text{ millones de dólares}]
    =500.000\text{ dólares}.
    
    Por tanto, votarán a favor de emprender este proyecto.
    
    ¿Qué ocurre con los tenedores de deuda? No obtienen nada si esta apuesta tiene éxito porque la empresa ya gana lo suficiente para realizar sus pagos de intereses. Pero si la apuesta fracasa y la empresa quiebra, no solo pierden sus 500.000 dólares anuales en pagos de intereses, sino también su inversión inicial de 5 millones de dólares. Por tanto, desde su perspectiva, se trata de una apuesta con un 50 % de probabilidades de ganar cero (si tiene éxito, siguen recibiendo simplemente su pago de intereses de 500.000 dólares) y un 50 % de probabilidades de perder 5 millones de dólares en activos y 5 millones de dólares de valor actual descontado de los futuros pagos de intereses. Para los tenedores de bonos, el valor esperado de esta inversión es de −5 millones de dólares.
    
    En este caso, existe un claro conflicto de intereses entre los titulares de capital y los tenedores de deuda. Los titulares de capital, que son propietarios de la empresa y tienen control sobre quienes toman las decisiones, si no sobre las propias decisiones, apoyarían emprender este proyecto, mientras que los tenedores de deuda no lo harían. Esto constituye un problema porque la proporción de capital es muy pequeña. Si la proporción de capital fuera de 5 millones de dólares y la proporción de deuda fuera de 1 millón de dólares, los titulares de capital no apoyarían este proyecto. Esto se ilustra en el panel inferior de la Tabla 24-1 (suponiendo que la empresa paga ahora 500.000 dólares en dividendos a los titulares de capital y 100.000 dólares en pagos de intereses a los tenedores de bonos). En esta situación, los titulares de capital no quieren arriesgarse con la nueva inversión porque tienen demasiado que perder: el rendimiento esperado de la inversión para ellos es ahora de −3,5 millones de dólares. Como antes, los tenedores de deuda tampoco quieren realizar esta inversión arriesgada porque su valor esperado sigue siendo negativo. En este caso, los intereses de los tenedores de deuda y de los titulares de capital están ahora alineados (ambos pierden dinero), y el proyecto no se lleva a cabo.
    
    La idea fundamental aquí es que, a medida que aumenta la proporción de la financiación de la empresa que corresponde a deuda, aumenta el potencial de este conflicto de intereses. Este potencial aumenta porque, a medida que aumenta la proporción de deuda, los tenedores de deuda soportan una proporción cada vez mayor de los proyectos que terminan en quiebra, mientras que los titulares de capital tienen un riesgo cada vez menor al realizar una apuesta.
    
    Ahora imagine que trabaja para un banco que está considerando prestar su dinero a una pequeña empresa. ¿Preferiría prestar a una empresa con un 50 % de capital y un 50 % de deuda o a una con un 10 % de capital y un 90 % de deuda? Siguiendo la lógica del ejemplo que acabamos de analizar, preferiría prestar a la empresa con un 50 % de capital porque una mayor cantidad del capital de los propietarios está en riesgo, por lo que tomarán decisiones más sensatas. Le preocuparía que, si presta a la empresa con solo un 10 % de capital, esta asuma riesgos disparatados que tengan un valor esperado negativo para usted como tenedor de deuda.
    
    Como resultado de este problema de agencia, los bancos (y otros prestamistas) cobrarán tipos de interés más altos por los préstamos a las empresas a medida que aumente su proporción de financiación mediante deuda. Estos tipos de interés más altos compensan la ventaja fiscal de la deuda, de modo que las empresas se enfrentan ahora a una disyuntiva: más deuda significa una mayor ventaja fiscal, pero potencialmente mayores costes de financiación porque los bancos temen conceder préstamos a empresas fuertemente endeudadas.

    El análisis anterior sugiere que el impacto de la imposición sobre sociedades en la estructura financiera de las empresas —y, en particular, en la dependencia de la deuda— no está claro. Por supuesto, esta cuestión no puede abordarse simplemente comparando si las empresas que se enfrentan a mayores cargas del impuesto sobre sociedades tienen niveles de deuda más elevados. ¡Los propios impuestos sobre sociedades son función de diversas decisiones de las empresas, incluida la estructura financiera! Una empresa con más deuda tendrá pagos de impuestos más bajos, por lo que una regresión de los niveles de deuda sobre los pagos del impuesto sobre sociedades producirá una estimación sesgada.

    HEIDER y LJUNGQVIST (2012) abordaron este problema utilizando múltiples cambios en los impuestos estatales sobre sociedades en Estados Unidos. Durante el período 1990–2011, encontraron 38 casos de estados que modificaron sus tipos del impuesto sobre sociedades. Compararon el efecto de estos cambios en los tipos impositivos sobre las empresas de esos estados con el efecto sobre otras empresas del mismo sector en estados cercanos. De este modo, podían utilizar empresas comparables para identificar cómo afectan los cambios fiscales a la estructura financiera de las empresas.
    
    Sin embargo, es posible que estos cambios en los impuestos estatales sobre sociedades respondan ellos mismos a los factores que determinan el apalancamiento de las empresas. Por ejemplo, una perturbación económica negativa en un estado puede llevar a las empresas a elevar los tipos del impuesto sobre sociedades para compensar la caída de los beneficios empresariales; al mismo tiempo, durante los períodos de recesión, las empresas pueden necesitar pedir más préstamos, aumentando sus ratios de deuda. Los autores presentaron un enfoque ingenioso para abordar esta preocupación: examinaron empresas que se encuentran justo cerca de la frontera entre estados que modifican los impuestos y estados que no los modifican. Tales empresas se enfrentan a condiciones económicas similares, pero los recortes fiscales solo afectan a algunas de las empresas y no a otras.
    
    Utilizando este enfoque empírico, los autores encontraron que existe un efecto muy considerable de los cambios en los tipos estatales del impuesto sobre sociedades sobre la forma en que se financian las empresas. Los aumentos de impuestos conducen a un mayor uso de la deuda: el aumento medio del impuesto estatal sobre sociedades en su muestra, del 13 %, condujo a un aumento del 4,5 % de la deuda como proporción de la financiación de las empresas. Al mismo tiempo, encuentran que este efecto es asimétrico: los recortes en la imposición sobre sociedades no conducen a reducciones del apalancamiento de las empresas. Esta sorprendente asimetría sugiere que intervienen más factores que la simple optimización por parte de las empresas de su estructura financiera con respecto a la imposición.
\end{specialblock}

\paragraph{Tratamiento de la inversión: Depreciación}
La otra cara de la moneda, y componente final de los gastos, es la depreciación de las inversiones en capital, es decir, el ritmo al que las inversiones en bienes de capital pierden valor con el paso del tiempo. Cuando la empresa compra maquinaria o un edificio nuevo, está invirtiendo en un bien que proporcionará servicios durante muchos años. Por ello, ningún código tributario permite deducir el coste total de la máquina en el período en que se adquiere. En su lugar, se permite a la empresa deducir de sus impuestos las amortizaciones fiscales, que son deducciones fiscales diseñadas para aproximar el ritmo al que la inversión en capital pierde valor. 

Las empresas aplican estas amortizaciones fiscales cada año durante un determinado número de años, distribuyendo así el coste del activo a lo largo de varios ejercicios. La determinación adecuada de las amortizaciones fiscales es una cuestión muy importante para la política del impuesto sobre sociedades. 

¿Cómo deberían tratarse fiscalmente estas adquisiciones de bienes de larga duración? En principio, el código tributario debería permitir que las empresas dedujeran la depreciación económica, es decir, el deterioro del valor de la máquina en cada período, como un gasto. La máquina se irá desgastando gradualmente con el tiempo y, después de un tiempo, la máquina carecerá de valor, debiendo ser sustituida. Por tanto, el coste para la empresa de utilizar esa máquina durante un año es igual al precio de compra menos su valor al cabo de un año, porque esa cantidad constituye el coste para la empresa de utilizar la máquina durante ese año y, por tanto, es el gasto apropiado que debe deducirse de los beneficios. Del mismo modo que utilizar un año de trabajo cuesta el salario correspondiente a ese año, utilizar un año de una máquina cuesta la reducción de valor que experimenta la máquina durante ese año.

El problema de aplicar la depreciación económica en la práctica es que la verdadera tasa de depreciación económica normalmente no se observa y varía considerablemente entre unos activos y otros. La depreciación económica de un edificio de oficinas puede ser bastante reducida porque el desgaste del edificio es gradual. Sin embargo, la depreciación económica de una máquina industrial puede ser mucho mayor debido al desgaste ocasionado por el proceso productivo. Como consecuencia de estas incertidumbres, el código tributario ha adoptado una serie de calendarios de amortización para diferentes categorías de activos. 

Un método consiste en tomar la vida útil típica del activo y distribuir el importe de la depreciación de manera uniforme entre todos los años de vida del activo: esquema de amortización lineal. En otros casos, la depreciación puede producirse con mayor rapidez. Para esos casos, el gobierno puede ofrecer calendarios de amortización acelerada o, alternativamente y dentro de un mismo horizonte temporal, el gobierno podría permitir una amortización más concentrada en los primeros años, con deducciones por depreciación más elevadas al comienzo de la vida útil del activo. En el caso extremo, el gobierno puede permitir que las empresas imputen inmediatamente como gasto las inversiones en capital físico y deduzcan el coste íntegro del activo en el mismo año en que se realiza la inversión. El punto fundamental para modelizar el impacto del sistema tributario sobre las decisiones de inversión es que el valor de las deducciones por depreciación aumenta cuanto más rápidamente pueden aplicarse. Ello se debe a que el valor actual descontado (VAD) de cualquier beneficio fiscal es mayor cuanto antes se obtiene dicho beneficio.\footnote{%
    Pongámos un ejemplo. Supongamos que la empresa de nuestro ejemplo puede endeudarse a un tipo de interés del $10\%$, de modo que utiliza ese $10\%$ como tasa de descuento para realizar los cálculos del valor actual descontado. Consideremos, en primer lugar, la depreciación económica (que, en este caso, coincide con la amortización lineal). El valor actual descontado del conjunto de deducciones por depreciación es: $10.000+\frac{10.000}{1,1}+\frac{10.000}{(1,1)^2}+\cdots+\frac{10.000}{(1,1)^9}=67.590$. Supongamos ahora, en cambio, que la empresa puede aplicar una amortización lineal acelerada, de modo que puede amortizar 20.000€ al año durante cinco años. El valor actual descontado de ese conjunto de deducciones por depreciación es: $20.000+\frac{20.000}{1,1}+\frac{20.000}{(1,1)^2}+\frac{20.000}{(1,1)^3}+\frac{20.000}{(1,1)^4}=83.397$. Al acelerar el ritmo al que la empresa puede amortizar su máquina, el gobierno ha permitido que la empresa deduzca más de 15.000€ adicionales de su renta imponible en términos de valor actual. En el caso extremo de la imputación inmediata como gasto (expensing), el valor actual descontado de la deducción por depreciación asciende a 100.000€, porque toda la deducción se realiza en el primer período.
}


\paragraph{Tratamiento de la inversión: Créditos fiscales}
El último componente para calcular la carga tributaria de las sociedades son los créditos fiscales que reducen el importe del impuesto sobre sociedades. 

El más conocido de ellos es el crédito fiscal por inversión (ITC). El ITC permite a las empresas deducir un determinado porcentaje de sus gastos anuales de inversión que cumplen los requisitos establecidos del importe de impuestos que deben pagar. El ITC fue una característica periódica del sistema estadounidense del impuesto sobre sociedades, aunque dejó de estar vigente después de 1986 (cuando las empresas podían recibir un crédito equivalente al 6 %–10 % de sus gastos de inversión, dependiendo de la vida útil del activo). No obstante, otros créditos fiscales siguen ocupando un lugar destacado en el código tributario; como informa KOCIENIEWSKI (2011), el fabricante de videojuegos Electronic Arts pudo aprovechar suficientemente los créditos fiscales existentes (y otras lagunas fiscales) como para pagar casi ningún impuesto sobre 1.200 millones de€ de beneficios.



Después de estos ajustes se suele realiwr una serie de reducciones o compensaciones. Un fenómeno que ocurre en este tipo de tributos es que cabe la posibilidad de detraer, de la determinación de la base imponible, bases imponibles de ejercicios anteriores que son negativas. En este sentido se llevan hacia adelante "loss carryforward'' las pérdidas de ejercicios pasados. Esto tiene especial relevancia en los casos en los que se supera una crisis pues se están generando una ser.ie de bases negativas que harán que el tributo no suponga una carga en la fase expansiva. 



Sobre esto el legislador puede establecer otra serie de reducciones que considere necesarias a la hora de determinar la base del impuesto y que pueden ir desde aportaciones a fondos de previsión como generación de empleo o inversiones en determinadas regiones.

Señalar en este punto que este proceso de paso de un resultado contable a una base imponible es vital para la de finición del tributo y su potencial recaudador. Caben dos opciones. Que la base del tributo vaya en sincronización con los resultados contables, lo que a priori parece deseable. pues se está consiguiendo gravar a la capacidad económica y fuente de renta planificada. O que en caso contrario la base imponible se distancie mucho del resultado contable, lo cual genera problemas de interpretación del tributo.

V. IV. Deduccion es, retenciones y pagos a cuenta.

Como ya hemos seifalado la progresividad en este tributo, de existir, puede introducirse en las deducciones y boni ficaciones. que se establecerán de acuerdo a decisiones de política eco nómica. como puede ser invertir en un territorio determinado. en unas actividades determinadas o tener en cuenta unas peculiares caract erísticas de las empresas. Por ejemplo en Francia con respecto a los territorios de ultramar.

El elemento verdaderamente importante aquí son los pagos fracci onados, un tipo específico de retención que implica, sin embargo, un proceso de liquidación tributario completo. Es decir, las empresas presentarán una declaración simplificada y pagarán por los beneficios fiscales hasta el momento de presentar el pago fraccionado. Este sistema es acumulativo y se presen taran cuantos pagos considere la norma teniendo en cuenta los datos anteriores. Es decir, el segundo pago descontará el primero e incorporará nueva información y así sucesivamente. De esta forma se adelanta en el tiempo el pago del impuesto sin necesidad de esperar a que las empresas cietTen su contabilidad.

En algunos países se establece un pago mínimo que consiste en que independientemente de las normas del tributo las empresas pagarán un xo/o de su resultado contable como pago fraccionado. Si esta medida solo existe en el pago fraccionado, estaremos hablando de un adelanto de la empresa al Estado pues conllevará una devolución al representar la declaración anual. Si, sin embargo, se incorpora dicho pago mínimo a la cuota de la declaración en este caso no se trata de un adelanto si no de un suelo no1mativo que indica que como poco las empresas deben pagar una ¡ cantidad x% de sus beneficios al Estado.







\subsubsection{¿Cómo se determina la renta sujeta a gravamen? Tratamiento del beneficio extraordinario}

V.J. La renta sujeta a tributación.

Nuestro primer paso consiste en definir qué renta es a quella que tributa: la de la empresa a nivel nacional o la global o mundial. Esta pregunta no es baladí especialmente cuando se está considerando una economía tremendamente glob·a lizada. La respuesta es sencilla y lógic a :

La empresa tributara por su renta global.

Esta definición nos acerca a la idea y diferencia entre PIB y PNB. Como se repite en innombrables ocasiones en los temas de contabilidad pública y de macromagnitudes, el concepto de na cionalidad es imprescindible para el reparto de rentas y cómo no de los tributos, que giran sobre rentas, deben fundamentar su carga en un concepto de nacionalidad. Si bien esta nacionalidad no es idéntica a la de la contabilidad p ública, pues habrá empresas, que son filiales de otras extranjeras, que, por el mero hecho de estar radicadas en territorio español, mediante una persona jurídica o establecimiento perma nente. tendrán la consideración de nacionales.

En este sentido tenemos que partir del concepto de resultado contable para definir la renta del tributo. Una empresa en esta figura resume el conjunto de sus actividades de carácter universal independientemente de dónde se haya producido el gasto o el ingreso.

Debemos hacer ahora una pequei'ia pausa para hablar de los grupos ) la consolidación. Este concepto primordial para este tipo de impuestos introduce una complejidad técnica impo11ante y simplemente señalaremos su existencia y consecuencias. Un grupo fiscal es un conj unto de personas jurídicas que tienen relación en su capital social. estableciéndose relaciones de control. En este sentido tendremos en cada grupo una matriz y un conjunto de filiales que dependen de esa matriz. La matriz controla las fi liales y sobre las cuentas funcionan como única empresa. La existencia de grupos fiscales hace que exista un fenómeno que se denomina consolidación. La consolidación es la conversión de un conjunto de cuentas ··N+n'· en una cuenta única ·•111·· aplicando una serie de convenciones contables que lo que establece finalmente es un resultado contable consolidado para todo el grupo. La consolidación por lo tanto es una suerte de neto o clearing, que hace, siguiendo no1mativa contable y fiscal. compensar saldos de las empresas del grupo (positivos y negativos) para ofrecer un resultado contable del grupo. Qué decir que este tipo de operaciones son altamente ventajosas para los grupos consolidados y que, en algunas jurisdicciones (México) se ha decidido eliminar dicha práctica.

El esquema del impuesto es

Resultado contable.
+/- Ajuste al resultado contable
-Bases imponibles negativas previas y otros ajustes temporales.
= Base Imponible
X Tipo de gravamen
= Cuota i ntegra
+/- Bonificaciones
+/- Deducciones
=Cuota líquida positiva
-Retenciones
-Pagos a cuenta
= A ingresar/ devolver

V.I J. La determinación de la base imponible.

Partiendo de un resultado contable o resultado contable consol i dado se realizan una serie de ajustes al mismo. Lo que hacen este tipo de nOJmas es adaptar un resultado contable al deseo de l egislador tributario. Pongamos por simplicidad el ejemplo de las amo11izaciones. En la contabilidad privada 110 existe l imitación a la amortización de manera que una empresa decide que quiere amortizar la inversión '·y·· en un ejercicio. Esto, aunque legitimo contablemente y acorde a la idea de imagen fiel, es totalmente erróneo según criterios contables fiscales pues la norma fiscal establece que como mucho se puede amortizar un 30\% cada ejercicio. De esta fonna se sumará al menos un 70\% del valor amortizado al resultado contable para la determinación de la base imponible. Estos ajsutes al resultado contable son muy importantes y determinarán qué tipo de base imponible se desea obtener. En este sentid0 las ideas generales que rigen este tipo de ajustes son:
• Los gastos se deben ajustar a lo efectivamente realizado, no a lo contablemente reflejado.
• Se deb,m excluir determinados ingresos por considerar que existe doble imposición. 
• Se deben ajustar los resultad0s a determinados criterios de imputación temporal, generalmente más duros que los puramente contables y <lcordes al devengo del impuesto. 



V. lll. El tipo de gravamen.

Sobre la base imponible así determinada se aplicará el tipo de gravamen. Este suele ser único. es decir no estamos ante una escala progresiva si no que se aplicará un tipo nominal a toda la base imponible. Esto es así porque una escala conllevaría generar un cuello de botella al que las empresas dedicarían mucho esfuerzo en no superar. cosa que las personas físicas no pueden hacer
con tanta facilidad. Así pues. la progresividad de este tributo es reducida y solo provendrá de la articulación de determinadas deducciones o gravámenes compleme111arios asociado�� al volumen de la cifra de negocio. resultados o base imponible.

Sin embargo. no es raro que se establezcan tipos diferentes entre tipos de empresas diferentes, que buscan favorecer a la equidad horizontal. En este sentido se pueden establecer estos tipos en función del tipo de empresa, sector en el que trabaja (financiero o no), el tamaño o de su antigüedad.

La diferenciación en función de la antigüedad es óptima en el sentido de que es temporal y hace que las empresas no dependan de un beneficio fiscal para asegurar su rentabilidad, si bien es cieno que no tiene en cuenta las necesidades de cada empresa y que una empresa por sus periodos de maduración puede necesitar más tiempo que otra.

La otra opción es establecer diferencias en función del tipo de empresa (crédito. industrial, fondos de pensiones etc ... ). La razón para este tipo de diferenciaciones es simple y llanamente de política económica, un país puede preferir que sus bancos paguen más que otro tipo de empresas o viceversa. Los países exportadores de materias primas pueden optar por que sus industrias exportadoras no paguen o que sean las únicas que paguen beneficiando al resto de empresas.



V.V. Las m ermas a la base im po nible y a la recaudación. E l problema de la internacionalización.

Dos problemas surgen con respecto a la tributación de las personas jurídicas.

El primero es temporal, pues los aj ustes de hases imponibles negativas provenientes del pasado o deducciones que no se han podido aplicar suponen que existirán periodos de génesis de estos saldos negativos. generalmente en las crisis y que luego se i rán aplicando a medida que la recuperación de las cuentas ocurre. En este sentido, la crisis golpeará dos veces sobre la recaudación tributaria. una primera cuando ocurre y otra segunda cuando ocurre la recuperación, pues las empresas podrán utilizar este sistema de traslado de saldos negativos a futuro para reducir su factura fiscal. Este sistema es un sistema casi universal que existe en muchos países que simplemente cuenta con dos tipos de posibles limitaciones.
• Temporal. los saldos negativos tienen un pe1iodo de validez o
• Cuantitativa. sólo se puede aplicar una cantidad de dichos saldos cada vez.

Un debate en torno a este tipo de diseí'íos tributarios es el de justicia tributaria pues las personas jurídicas disponen de estos mecanismos, mientras que las personas físicas en general no. El segundo problema está asociado a la internacionalización ele la empresa. A medida que la empresa está más presente en el exterior, las bases imponibles también lo están y es cada vez más difícil asegurar una carga justa en estos tributos sin caer en la doble imposición. En este sentido, se han desarrollado por parte de determinadas personas j urídicas esquemas que se denominan de planificación fiscal agresiva que consisten básicamente en trasladar bases imponibles a otros territorios donde la tiscalidad es más laxa (en ningún caso estamos hablando de paraísos fiscales, si no de territorios simplemente donde los tipos son más bajos, la problemática de paraísos fiscales es ajena a este tema). El traslado de bases imponibles se ha denomi nado BEPS (Base Erosion and Profit Shifting) y consiste en usar mecanismo de transferencia (generalmente en precios) o simplemente de facturación. para trasladar los resultados contables a otras jurisdicciones y así la carga fiscal. Sin entrar en detalles de cómo frenar este fenómeno, sei'ialar que la OCDE ha desaiTollado un ambicioso plan para acabar con este tipo de comportami entos tributarios y favorecer que las empresas paguen los tributos allá donde tienen su sede y se genera la actividad, y no solo de acuerdo a la primera. Dentro de este plan se busca la participación de países en desarrollo en un esfuerzo de desarrollar un estándar internacional global.


Es cierto que no gravar la renta de las sociedades crea oportunidades para la elusión del impuesto personal. Pero un impuesto especial sobre las sociedades no es la única manera de incluir los beneficios acumulados en las sociedades. Al final de este capítulo analizamos un método alternativo que muchos economistas consideran superior.



\textbf{TIPO}

Las sociedades tributan sobre sus beneficios netos (beneficios menos gastos) de acuerdo con una escala de tipos impositivos aproximadamente progresiva, mostrada en la Figura 24-3. Las empresas más pequeñas, con beneficios netos inferiores a 50.000€ anuales, pagan un tipo impositivo del 15 %. Posteriormente, los tipos aumentan a medida que crecen los beneficios netos de la empresa, alcanzando un máximo del 39 % para las empresas con beneficios netos comprendidos entre 100.000 y 335.000€, y entre 1,5 millones y 1,83 millones de€, antes de descender al 35 % para las empresas con beneficios de 1,833 millones de€ o superiores. La inmensa mayoría de las grandes sociedades está sujeta al tipo impositivo del 35 %. (El «abultamiento» (bubble) consistente en aplicar un tipo superior del 39 % entre 100.000 y 335.000€ y entre 1,5 y 1,83 millones de€ se introdujo para que tanto el tipo impositivo medio como el tipo marginal de las empresas situadas por encima de esos niveles fueran del 35 %.)





\subsection{Tributación sobre el ahorro: Patrimonio personal}
Cuanto mayor sea la riqueza de un individuo, mayor será su capacidad de pago, permaneciendo iguales las demás cosas. Por lo tanto, los individuos ricos deberían pagar impuestos más altos. 

Los impuestos sobre la riqueza son aquellos que gravan a intervalos regulares y frecuentes el valor de los activos o el patrimonio neto de las unidades institucionales, en este sentido se trata de impuestos que giran sobre un stock que las familias o las empresas poseen. Expresado en términos contables el patrimonio es el balance resultante de activos menos pasivos. o pat rimonio neto. Generalmente estamos hablando de un impuesto que recae sobre las personas físicas pues el patrimonio de las personas jurídicas es un activo de las contribuyentes representado a través de las acciones de las compaiñías. Es por ello. que el contribuyente será en casi todos los sistemas una persona física. Si bien existen opciones de tributación al patrimonio de las empresas como la introducida en Catalui'ia sobre los activos no productivos de las mismas.

Debemos recordar que. en paralelo a este diseño de la tributación sobre la renta. la riqueza se ha venido gravando de manera histórica. La riqueza es la manifestación de capacidad económica más fácilmente gravable. tiene un componente inmóvil. el inmobiliario; aunque no toda ella es inmóvil. A su vez es fác ilmente identificable e individualizable por parte de .las autoridades. lo que hace que sea fácil establecer un sistema de tributación recurrente sobre la misma. lo que lo diferencia de la tributación sobre el capital.

Como se puede intuir, vamos a analizar la tercera y última figura tributaria periódica: la imposición sobre la riqueza/ahorro. En este sentido, muchos autores defiendes que la imposición sobre la riqueza, o rentas del ahorro, ayuda a corregir ciertos problemas (inevitables) que surgen en la administración de un impuesto sobre la renta. Recuérdese que, en puridad, todas las ganancias de capital (realizadas o no) pertenecen a la base imponible de un impuesto integral sobre la renta (HAIG-SIMON). En la práctica, a menudo es imposible gravar las ganancias de capital no realizadas. Al gravar la riqueza de la cual estas ganancias pasan a formar parte, quizá esta situación pueda remediarse. Ahora bien, es cierto que la riqueza en un momento dado incluye la suma de las ganancias y pérdidas de capital de todos los años anteriores. Sin embargo, no hay razón para creer que el rendimiento de un impuesto anual sobre la riqueza se aproxime a los ingresos que se habrían generado mediante la tributación anual completa de las ganancias de capital no realizadas.

En prinicipio, la base bruta sujeta a gravamen por un impuesto de estas características debería ser sencillo: el conjunto de bienes y derechos de contenido económico de los que la persona física sea titular. 

Es decir, como hemos indicado antes, el resultado del balance o patrimonio neto de un contribuyente, que podría reflejarse en algo parecido.

\textbf{ACTIVO | PASICO: ESQUEMA ICEX-CECO}

VI.U. La base imponible y el tipo el e gravamen.

Este tipo de impuestos son se ncillos en la definición de la base imponible, esta es el patrimonio. Con todas las exenciones o no sujeciones que el legislador considere y reducido en la forma que el mismo sei'iale. Es decir, por ej emplo. se suelen excluir de la medición las empresas f'amiliares, pues estas son una fuente de renta: en el pasado, se excluían los aperos de labranza por ser útil de trabajo.

\subsubsection{¿Por qué gravar las rentas derivadas del ahorro?}
Supóngase que una persona ha acumulado un enorme tesoro de oro que no produce renta. ¿Debería ser gravada sobre el valor del tesoro? Algunos creen que, mientras esta persona estuviera sujeta al impuesto sobre la renta mientras el tesoro se estaba acumulando, no debería ser gravado de nuevo. Otros sostendrían que el oro per se genera utilidad y debería estar sujeto a impuestos.\footnote{%
    Quizá el principal problema del argumento de la capacidad de pago es que incluso las personas ricas tienen un componente sustancial de su riqueza en capital humano. Sin embargo, no hay forma de valorar el capital humano excepto por referencia a la renta que produce, lógica que nos lleva de nuevo a la renta como la base apropiada.

E l impuesto del patrimonio es un impuesto que ha sido muy debatido recientemente por dos motivos. E l primero es que es un tributo que grava el patrimonio corno acumulación de riqueza y puede argumentarse que es un doble gravamen sobre la misma fuente de capacidad económica.

El segundo es que es un impuesto de poca afectación en número de contribuyentes, con faci l i dad de gestión y que favorece en una manera muy significativa la equidad del sistema. Los debates ideológicos hacia ambos lados son inagotables y los argumentos también. En este punto sólo vamos a sefíalar los principales puntos a favor o en contra.
}

Entonces, ¿por qué debe tributar la riqueza de forma independiente? Existen como mínimo tres argumentos en favor de este tipo de impuestos.


\paragraph{Concentración de riqueza: }
En primer lugar, la tributación de la riqueza reduce la concentración de la riqueza, lo cual es deseable social y políticamente. (\textcolor{red}{Como vimos en el Capítulo 12}) Aunque es difícil medir la renta con precisión, las mejores estimaciones sugieren que, salvo excepciones, la distribución de la renta es bastante desigual. La calidad de los datos sobre la riqueza es aún menor. La información que existe sugiere que la distribución de la riqueza es muy desigual. Una preocupación relacionada es que una distribución altamente concentrada de la riqueza conduce a la corrupción de los procesos políticos democráticos. Los escépticos responden que, si la concentración de poder es el problema, entonces no hay justificación para gravar incrementos de riqueza de $1$ millón de euros, $10$ millones de euros o incluso $50$ millones de euros. Como señala STEIN (1997), [...] \textit{se necesita mucho más dinero que eso para generar poder en Estados Unidos hoy}. STEIN observa además que existen fuentes de influencia distintas del dinero: \textit{Oprah Winfrey tien[e] más poder que cualquier persona megarrica hoy}. ¿Debería Oprah enfrentarse a un impuesto especial porque es poderosa?

El segundo argumento sostiene que los impuestos sobre la riqueza son necesarios para evitar una concentración excesiva de riqueza y de poder en la sociedad en manos de unas pocas dinastías adineradas. Muchos autores han sostenido que la existencia de un impuesto sobre la riqueza es fundamental para contribuir al mantenimiento de una sociedad meritocrática, garantizando que las personas sean recompensadas por su talento y no únicamente por la posición social en la que nacieron.\footnote{%
    En términos muy contundentes, Frank Keating, presidente del American Council of Life Insurers, criticó los intentos de derogar este impuesto afirmando: \textit{Estoy institucional e instintivamente en contra de enormes bloques de riqueza heredada. No creo que necesitemos al Vizconde de Enron ni al Duque de Microsoft}.
} PIKETTY y SAEZ (2006) analizaron la evolución de la proporción de ingresos correspondiente a los niveles más altos de la distribución de la renta (por ejemplo, el $0,1\%$ superior) en una amplia muestra de países a lo largo del siglo XX. Encontraron que, en los países con sistemas tributarios progresivos tanto sobre las rentas del capital como sobre la riqueza, se produjo una enorme reducción de la proporción de ingresos controlada por los perceptores situados en la cúspide de la distribución. Este resultado respalda la idea de que la tributación de la riqueza y del capital desincentiva la concentración de la riqueza a largo plazo.\footnote{%
    En los últimos años, sin embargo, se ha producido de nuevo un fuerte aumento de la proporción de ingresos correspondiente a las familias con mayores ingresos en algunos países —especialmente en Estados Unidos—, de modo que la proporción de la renta correspondiente al $0,1\%$ de las familias más ricas ha vuelto a situarse en niveles similares a los de comienzos del siglo XX. En otros países —como Francia y Japón— no se ha producido una \textit{reconcentración} comparable de la riqueza.
}

En este sentido, no obstante, también hay quien sostiene que el impuesto sobre sucesiones representa una doble imposición: los ingresos ya fueron gravados cuando se obtuvieron, ya fuera en el mercado de trabajo o mediante intereses sujetos a tributación, y posteriormente vuelven a gravarse cuando los hijos los reciben tras el fallecimiento de sus padres. Este resultado es criticado por dos motivos. El primero es la equidad horizontal: ¿por qué quienes pagan el impuesto sobre sucesiones deberían tributar una segunda vez por esos mismos ingresos? El segundo es la eficiencia: el hecho de que el patrimonio vaya a soportar una tributación tan elevada puede reducir el incentivo para ahorrar con el fin de dejar una herencia a los hijos y, por tanto, distorsionar las decisiones de ahorro. Si al fallecer se va a perder hasta el $48\%$ de lo ahorrado, es posible que una persona decida gastar ese dinero en un crucero en lugar de ahorrar para sus hijos. No obstante, este argumento presenta tres problemas:
\begin{la}
    \item \textbf{Característica generalizada}. En primer lugar, la doble imposición constituye una característica generalizada de los sistemas tributarios tanto en Estados Unidos como en el resto del mundo. Cuando una persona utiliza su renta después de impuestos para comprar caramelos en la tienda de la esquina y paga un impuesto sobre las ventas por esa compra, está siendo gravada dos veces. ¿Por qué debería eliminarse o reducirse una forma concreta de doble imposición mientras que las demás permanecen? 
    \item \textbf{¿Efectos sobre el ahorro?}. En segundo lugar, la doble imposición no reduce necesariamente el ahorro, porque existen tanto efectos renta como efectos sustitución. Por una parte, las recompensas asociadas al trabajo y al ahorro disminuyen como consecuencia de la doble imposición, por lo que la persona puede decidir trabajar o ahorrar menos (efecto sustitución). Por otra parte, si una persona valora la herencia que desea dejar a sus hijos, una mayor tributación puede llevarla a trabajar o ahorrar más para mantener el valor de esa herencia después de impuestos (efecto renta); y,
    \item \textbf{¿Cuándo existe doble imposición?}. Por último, el argumento de la doble imposición se complica debido a su interacción con la actualización de la base de las ganancias de capital en el momento del fallecimiento. Los ingresos del trabajo efectivamente tributan dos veces para quienes pagan el impuesto sobre sucesiones: una vez cuando se obtienen y otra al fallecimiento. Del mismo modo, los intereses y otras rentas del capital devengadas sujetas a tributación también tributan dos veces: una cuando se perciben los rendimientos y otra al fallecimiento. Sin embargo, las rentas procedentes de las ganancias de capital pueden llegar a tributar solo una vez, en el momento del fallecimiento. Si los individuos transmiten sus activos de capital a sus hijos, no existe tributación sobre las ganancias de capital acumuladas porque la base del activo aumenta cuando este se transmite. Como consecuencia, en ausencia de un impuesto sobre sucesiones, las rentas derivadas de ganancias de capital transmitidas entre generaciones podrían llegar a escapar por completo de la tributación.
\end{la}

Más aún, ¿es reprochable la doble imposición? Cabe recordar que el desprecio hacia la doble imposición carece, verdaderamente, de sustento jurídico. Se trata, en esencia, de su conexión al principio \textit{non bis in idem} por una falsa analogía. El derecho tributaria no es derecho penal, no si quiera derecho sancionador. En este sentido parece más bien el éxito de una guerra cultural vencida por los oradores más neoliberales que un verdadero argumento jurídico que pueda vertirse en contra de esta figura impositiva.\footnote{%
    La evidencia sobre el impacto del impuesto sobre sucesiones en el comportamiento es revisada por KOPCZUK (2010). Este autor concluye que el tamaño de los patrimonios hereditarios responde de forma moderada a la tributación sucesoria: cada incremento del $10\%$ en el tipo del impuesto sobre sucesiones conduce a patrimonios entre un $1\%$ y un $2\%$ más pequeños. Esto puede reflejar tanto una reducción real del tamaño de los patrimonios como los esfuerzos por trasladar activos hacia formas no sujetas al impuesto sobre sucesiones. Por tanto, aunque unos impuestos sucesorios más elevados pueden afectar al ahorro y a la oferta de trabajo, es probable que dichos efectos sean reducidos.
}

\paragraph{Equidad vertical: Progresividad del Sistema Fiscal}
El primero es que constituye un medio extremadamente progresivo de recaudar ingresos públicos. Como sostuvo William Gates Sr. (el padre del multimillonario de Microsoft), se necesita un impuesto sobre sucesiones para \textit{proteger nuestra democracia de una mayor acumulación de riqueza hereditaria} (GATES y COLLINS, 2002).\footnote{%
    SUMMERS, subsecretario del Tesoro en 1997, declararía ante un proyecto de ley que buscaba derogar dicho impuesto que: \textit{cuando se trata de [reducir] el impuesto sobre sucesiones, no existe otra justificación que el egoísmo. En términos de fondo, este argumento sobre el impuesto sobre sucesiones es tan malo como puede ser}.
}

La mayoría de académicos señalan que, debido a los elevados niveles de exención, el impuesto afecta únicamente a una parte muy pequeña de los ciudadanos más ricos en el momento de su fallecimiento. Por ejemplo, en Estados Unidos, la proporción de patrimonios hereditarios que pagan el impuesto siempre ha sido bastante reducida y, en 2016, se estimaba que únicamente el $0,2\%$ de todos los patrimonios hereditarios (aproximadamente $5.400$ patrimonios) pagaría el impuesto sobre sucesiones. Y, sin embargo, en 2014 este impuesto recaudó el $0,6\%$ de los ingresos totales del gobierno, una cantidad superior al gasto conjunto destinado por el gobierno a la Food and Drug Administration, los Centers for Disease Control and Prevention y la Environmental Protection Agency. Más aún, GALE y SLEMROD (2001a) también señalan que los herederos de estos patrimonios suelen presentar características de ingresos y riqueza similares a las de sus padres, de modo que el carácter progresivo del impuesto sobre sucesiones no depende de que la carga tributaria se calcule sobre el fallecido o sobre el heredero.

Dejando de lado la cuestión normativa de si el gobierno debería o no perseguir una distribución más igualitaria de la renta, analicemos la cuestión positiva de si es probable o no que un sistema eficaz de tributación de las sucesiones alcance este objetivo. Ciertamente, el supuesto predominante es que lo haría: \textit{Desde sus comienzos, el impuesto sobre sucesiones fue considerado como un contrapeso a una concentración indebida de la riqueza} (GALE y SLEMROD, 2000). Sin embargo, existen varias razones por las que gravar los legados podría resultar contraproducente y crear una distribución menos igualitaria de la renta:
\begin{la}
    \item \textbf{Efectos sobre el ahorro: Disminución de salarios}. Si el impuesto sobre sucesiones reduce el ahorro, habrá menos capital. Esto conduce a un salario real más bajo para el trabajo y, bajo ciertas condiciones, a una proporción menor de la renta destinada al trabajo. En la medida en que la renta del capital esté distribuida de manera más desigual que la renta del trabajo, el efecto es aumentar la desigualdad;
    \item \textbf{Efectos sobre la equidad intergeneracional}. Dentro de una generación, es probable que la mayoría de los individuos transfieran riqueza únicamente a otros que están en peor situación que ellos. Tales transferencias tienden claramente a aumentar la igualdad. Reducir tales transferencias voluntarias bien podría conducir a una mayor desigualdad. Supóngase que los padres cuyas capacidades de obtención de ingresos son mucho mayores que la media tienen hijos cuyas capacidades de obtención de ingresos están más cerca del nivel medio. (Este fenómeno se conoce como regresión hacia la media). Los padres acomodados, que desean compensar a sus hijos por su menor capacidad de obtención de ingresos mediante legados, tienden a disminuir la desigualdad entre generaciones. A la inversa, reducir tales transferencias aumenta la desigualdad intergeneracional;
    \item \textbf{Efectos sobre la desigualdad en el consumo}. Una preocupación relacionada es que el centro de atención de la política debería ser la desigualdad del consumo en lugar de la desigualdad de la riqueza. En la medida en que el impuesto sobre sucesiones anime a las personas ricas a gastar más dinero mientras están vivas, empeora entonces la desigualdad del consumo. 
\end{la}

Como vemos, desde un punto de vista positivo, el efecto de la tributación de las sucesiones sobre la desigualdad es ambiguo. La investigación empírica no ha determinado si domina el efecto que aumenta la igualdad o el que la disminuye.

Ahora bien, en esta misma linea de argumentación favorable debida a su carácter progresivo también podría argumentarse en favor de abolir la tributación de sucesiones y donaciones mediante la tributación por el aumento de renta que supone la incorporación de estos bienes. Muchos autores sostienen que es moralmente inapropiado gravar a los individuos en el momento de su muerte.\footnote{%
    Aunque, como ellos mismos señalan, la muerte no constituye ni una condición necesaria ni una condición suficiente para que se apliquen los impuestos sobre las transmisiones patrimoniales: las transmisiones entre personas vivas pueden dar lugar al impuesto sobre donaciones, y el $98\%$ de quienes fallecen no paga ningún impuesto sobre sucesiones. Además, la muerte constituye un momento natural para recaudar impuestos, porque es cuando resulta más sencillo comparar los recursos acumulados a lo largo de toda la vida por distintos individuos.
} Entonces, ¿por qué no gravar la incorporación de esos bienes mediante el impuesto sobre la renta? Después de todo, constituyen adiciones al consumo potencial y, según la definición HAIG-SIMONS, son por tanto renta para el receptor. Sin embargo, siempre ha existido una fuerte aversión a incluir las herencias y las donaciones en la base del impuesto sobre la renta. Tales ingresos simplemente no se perciben como pertenecientes a la misma clase que los procedentes de salarios e intereses. 

\begin{specialblock}
    
\end{specialblock}


\paragraph{Obligaciones peculiares: }
Los impuestos sobre la riqueza son pagos que los poseedores de riqueza deben realizar por beneficios que reciben del gobierno. Como dijo el presidente Theodore Roosevelt, \textit{el hombre de gran riqueza tiene una obligación peculiar con el Estado porque obtiene ventajas especiales de la mera existencia del gobierno}. 

Se podría argumentar, por ejemplo, que un objetivo principal del gasto en defensa es proteger (de enemigos extranjeros) nuestra riqueza existente. Si es así, quizá un impuesto sobre la riqueza sea un método justo para financiar la defensa. Además, el gobierno realiza ciertos gastos que probablemente benefician especialmente a los poseedores de riqueza. Si el Estado construye y mantiene una carretera que pasa por mi tienda, entonces me confiere un beneficio por el cual debería pagar. Aunque la idea de basar los impuestos en los beneficios tiene cierto atractivo, no está claro que ningún impuesto sobre la riqueza factible pueda alcanzar este objetivo. Un abogado que defendía el argumento a favor de gravar la propiedad preguntó retóricamente: \textit{¿[N]o es cierto que alguien con una casa dos veces mayor recibe el doble de beneficio de los servicios de policía y bomberos prestados a la propiedad?} (HAGMAN, 1978). 

Contrariamente a lo que aparentemente creía, la respuesta es probablemente no. El valor para un hogar determinado de la mayoría de los servicios proporcionados por el gobierno local depende de factores distintos del tamaño de la casa. Por ejemplo, el valor de la educación depende del número de hijos. Incluso el valor de los servicios de bomberos y policía depende de cuántos muebles haya en la casa y de cuánta protección de seguro se haya adquirido. Si la tributación según los beneficios es el objetivo, un sistema de tarifas de usuario por los servicios públicos sería más apropiado que un impuesto sobre la riqueza.

En resumen, los impuestos sobre la riqueza, ya sea como un añadido a un sistema de impuesto sobre la renta o a un sistema de impuesto sobre el consumo, han sido racionalizados tanto sobre la base de la capacidad de pago como sobre la base de los beneficios. Ambos conjuntos de argumentos muy controvertidos.





















El último argumento en contra del impuesto sobre sucesiones se basa en razones de cumplimiento y equidad. Existen diversas formas mediante las cuales los contribuyentes con mayores conocimientos fiscales pueden evitar legalmente el pago del impuesto sobre sucesiones.

Uno de los métodos más utilizados consiste en que los padres constituyan fideicomisos (trusts) para sus hijos. Un fideicomiso es un acuerdo jurídico mediante el cual una persona concede a un tercero (el fiduciario o trustee) el control de determinados activos, con la condición de que dichos activos se utilicen en beneficio de personas concretas, como sus hijos.

Supongamos que Kanga obtuvo 5 millones de dólares de sus inversiones y desea transmitir esa riqueza a su hijo Roo. Si entrega ese patrimonio a Roo mientras ella aún vive, deberá pagar el impuesto sobre donaciones; si espera hasta su fallecimiento, deberá pagar el impuesto sobre sucesiones. Sin embargo, Kanga puede legalmente depositar esos 5 millones de dólares en un fideicomiso administrado por Tigger, quien actúa como fiduciario. Siempre que Kanga especifique que Tigger solo puede utilizar ese dinero en beneficio de Roo, los fondos depositados en ese fideicomiso nunca estarán sujetos a tributación, aunque Roo reciba plenamente sus beneficios (suponiendo, por supuesto, que Tigger sea realmente digno de confianza). Los padres también pueden constituir fideicomisos de seguros, mediante los cuales las pólizas de seguro de vida se incorporan a un fideicomiso para que sus hijos reciban las indemnizaciones del seguro libres de impuestos.

Otro método popular (y legal) para evitar los impuestos sobre las transmisiones patrimoniales consiste en conceder a los herederos participaciones en una nueva empresa familiar. Si Jesús Gates III hubiera sido padre en 1986, cuando Microsoft salió a bolsa, podría haber entregado a cada uno de sus hijos un 1 % de las acciones de Microsoft sin generar un impuesto elevado sobre donaciones (porque las acciones apenas tenían valor en los primeros momentos de la empresa). Si lo hubiera hecho, sus hijos serían hoy muy ricos y habrían evitado los elevados impuestos sobre las transmisiones patrimoniales asociados a esa riqueza (aunque habrían tenido que pagar el impuesto sobre las ganancias de capital si hubieran vendido las acciones durante su vida). Los fideicomisos y las transmisiones de acciones constituyen, por tanto, dos mecanismos mediante los cuales las personas bien asesoradas jurídicamente pueden evitar legalmente el pago de los impuestos sobre las transmisiones patrimoniales.

Como consecuencia de estas alternativas, algunos califican el impuesto sobre sucesiones como un «impuesto voluntario», ya que únicamente terminan pagándolo quienes no poseen los conocimientos suficientes para evitarlo. Este constituye un resultado injusto que vulnera el principio de equidad horizontal. Algunos sostienen que, si un impuesto no puede aplicarse de forma equitativa, no debería aplicarse en absoluto. Una vez más, los costes derivados de estas desigualdades horizontales deben compararse con los beneficios de la recaudación obtenida mediante este impuesto: ¿es la injusticia tan importante como para justificar la eliminación de casi 20.000 millones de dólares de ingresos públicos anuales?

Un objetivo implícito del impuesto sobre sucesiones es gravar la riqueza al menos una vez por generación. Sin embargo, las personas pueden evitar el impuesto de varias maneras. Muchas de ellas implican establecer fideicomisos, que son acuerdos mediante los cuales una persona o institución conocida como fiduciario posee el título legal de unos activos con la obligación de utilizarlos en beneficio de otra parte. Como ejemplo del uso de fideicomisos para evitar el impuesto sobre sucesiones, considérese el problema al que se enfrentan unos padres que poseen pólizas de seguro de vida que nombran a sus hijos como beneficiarios. Los ingresos procedentes de las pólizas de seguro se incluyen en la sucesión bruta de los padres. Sin embargo, los padres pueden establecer un fideicomiso de seguros y asignar la póliza de seguro al fideicomiso. Dado que los padres ya no poseen la póliza, esta queda fuera de su sucesión, y sus hijos reciben el beneficio completo del seguro de vida.

Otra técnica relativamente sencilla y popular consiste en conceder a los propios herederos acciones de una sociedad de propiedad restringida. Específicamente, supóngase que Mickey constituye su empresa como sociedad y posee todas las acciones. Durante su vida, Mickey realiza donaciones de una parte sustancial de las acciones —pero menos de la mitad— a sus herederos, Morty y Ferdy. Si las transferencias se producen relativamente pronto en la vida de la empresa, las acciones no valen mucho, por lo que se incurre en poca o ninguna deuda por el impuesto sobre donaciones. Debido a que Mickey posee la mayoría de las acciones de la empresa, sigue al frente de la compañía y controla efectivamente el valor de las acciones transferidas. Si la empresa de Mickey prospera, para cuando él muera, las acciones de Morty y Ferdy pueden ser extremadamente valiosas. Mickey ha conseguido así transferir una riqueza sustancial a sus herederos y proteger la transferencia del impuesto sobre donaciones y sucesiones. ¿Qué ocurre con las acciones que Mickey todavía posee al morir? Existen otras técnicas más complicadas para protegerlas.

En resumen, existen muchos métodos disponibles para realizar transferencias intergeneracionales de riqueza sin soportar ningún impuesto y sin perder el control efectivo de la propiedad durante la propia vida. Muchas de estas técnicas de elusión son complicadas y costosas. Como se señaló en un informe del Comité Económico Conjunto del Congreso, las obligaciones del impuesto sobre sucesiones «dependen de la habilidad del planificador de la sucesión, en lugar de la capacidad de pago» [Joint Economic Committee, 2006, p. 34]. En efecto, las únicas personas que pagan el impuesto son aquellas que descuidan realizar la planificación apropiada. Sin embargo, incluso en los casos en los que el impuesto no genera ingresos, puede crear cargas excesivas y/o costes de cumplimiento para las personas que modifican su comportamiento para evitarlo.






A nivel SEC 20 10 se encuadrarían dentro de cada una de las categorías anteriores. pero hemos decidido separarlo por tener principios recaudatorios diferentes.



Un aspecto i mportante de este tipo de tributos es su complementariedad con los de la renta, pues si rven mutuamente de fuente de control. Un activo genera rentas y la acumulación de rentas es un patrimonio. En este sentido y sea corno fuere lo que sí parece adecuado es contar con figuras que faci liten información. Se opte o no por el gravar.




\subsubsection{¿Cuáles son las desviaciones contempladas? Gastos fiscales}



En segundo lugar, la relación entre el causante y el contribuyente es primordial. El cómo el sistema considere esta lo convertirá en un impuesto progresivo o no o un impuesto que tiene en cuenta las condiciones del contribuy ente o no. 

Por ejemplo, una exención generalizada a los contribuyentes que presentan minusvalía parece una decisión que incrementa la equidad del sistema, pues esta persona cuenta con mayores dificultades de generar remas. Sin embargo, una exención general a los hijos independientemente de su patrimonio propio o renta propia puede considerarse como regresiva para el sistema pues simplemente está teniendo en cuenta una circunstancia personal que nada tiene que ver con la capacidad económica. Es decir, tal como se defina el impuesto según esta relación entre causante y contri buyente tendremos un impuesto que incrementa la equidad del sistema y progresivo o no. Los puntos de decisión que se encontrará un legislador los podemos definir como: 
• Grado de parentesco. • Patrimonio o renta previa del contTibuyente. • Circunstancias personales del contribuyente (minusvalía, hij<Ís etc ... ). • Relación económica y social del contribuyente con el causante (dependencia). • Interés general que considere el legislador (conservación del patrimonio familiar).




\subsubsection{¿Cómo se determina la renta sujeta a gravamen? }
Según su delinición en el SEC 2010. Los impuestos sobre el capital son impuestos que gravan a intervalos irregulares y muy poco frecuentes el valor de los activos o el patrimonio neto de las unidades institucionales o el valor de los activos transferidos entre unidades institucionales como resultado de sucesiones, donaciones inter vivos u otras transferencias.

Así pues, en cuanto a la definición del patrimonio estaremos a lo indicado anteriormente. El patrimonio se definirá como el conjunto de bienes y derechos con contenido económico deducido de costes o deudas pendientes. En este ·caso el patrimonio a gravar será sólo el que reciba el contribuyente y no todo el patrimonio del causante. Así un heredero solo tributará por su porción de herencia y un donatario sólo tributará por el donativo recibido. Al definirse el patrimonio de manera no global, es importante establecer en este tipo de tributos normas claras de valoración y de desgravación. Es decir, que deudas serán desgravable y de ser conjuntas con otros elementos del patrimonio como lo serán.

Los impuestos sobre el capital comprenden: 
• los impuestos sobre las transferencias de capital: los derechos de sucesión y los derechos sobre las donaciones inter vivos que gravan el capital de los beneficiarios, excepto los impuestos sobre las ventas de activos; 
los gravámenes sobre el capital: los gravámenes ocasionales o excepcionales sobre el valor de los activos (o el patrimonio neto) de las unidades institucionales. 


La determinación de la base sujeta a gravamen, así como los tipos aplicados sobre la cuota que debe efectivamente liquidar el sujeto pasivo, es extremadamente heterogénea (el más, sin duda) de entre todos los tipos de activos en los que se deposita el ahorro.



Por ejemplo, pensemos en dos formas clásicas de riqueza características del Impuesto sobre Bienes Inmuebles. ¿Cuál debería ser el objeto sujeto a gravamen: el terreno o la edificación? El impuesto sobre la propiedad inmobiliaria diseñarse gravando el terreno o gravando las mejoras realizadas en ese terreno, como una casa, un edificio de oficinas o un centro comercial, o ambas.\footnote{%
    Esta distinción fue destacada por Henry George, un pensador del siglo XIX que creía que los rendimientos obtenidos por el trabajo y el capital reflejaban la actividad productiva de estos factores, pero que los rendimientos que correspondían a los propietarios de terrenos no lo hacían. Al propietario del terreno se le paga únicamente porque posee el terreno y no debido a ningún esfuerzo dedicado a la producción. George, por lo tanto, hizo una propuesta radical: eliminar los impuestos sobre el trabajo y el capital y sustituirlos por un único impuesto sobre el valor del terreno. George no propuso gravar las mejoras realizadas en el terreno porque este enfoque socavaría su premisa básica de que la propiedad, y no el esfuerzo, debería ser gravada.
} La forma en la que se diseña no es, en absoluto, baladí. Por ejemplo, las distorsiones generadas pueden provocar conductas deseadas que potencien la recaudación fiscal de una ciudad. Este ejemplo, sencillo, no es más que una forma particular de ilustrar uno de los grandes desafíos de esta figura impositiva: la obligación tributaria de un individuo es el producto del tipo impositivo y la base sujeta a gravamen, el valor que la jurisdicción asigna a esa forma de riqueza; aunque en la mayoría de los casos las jurisdicciones intentan hacer que los valores tasados correspondan a los valores de mercado, esto resulta muy difícil para el recaudador de impuestos, que a menudo desconoce el valor de mercado y, por tanto, debe hacer una estimación, quizá basada en los valores de mercado de bienes comparables que se hayan vendido recientemente.

\paragraph{Heterogeneidad en las formas de riqueza}
Los valores de mercado y los valores tasados divergen en una medida que depende de la precisión del procedimiento de estimación de la jurisdicción. La proporción entre el valor tasado y el valor de mercado se denomina ratio de tasación. Si todas las propiedades tienen el mismo tipo legal y el mismo ratio de tasación, sus tipos impositivos efectivos son los mismos. Volviendo al ejemplo de la vivienda, supóngase que los ratios de tasación difieren entre propiedades. Nacho y Ana poseen ambos propiedades que valen $100.000€$; la primera tasada en $100.000€$ y la segunda en $80.000€$. Claramente, incluso si se enfrentan al mismo tipo legal (digamos, el $2\%$), el tipo efectivo de Nacho del $2\%$ es superior al $1,6\%$ de Ana ($=1.600/100.000$). 

Para analizar el impuesto sobre la riqueza, desde el principio hay que darse cuenta de que literalmente existen miles de jurisdicciones que gestionan sus sistemas de impuestos sobre la riqueza de manera más o menos independiente (en España: local, autonómico y estatal). Ninguna jurisdicción incluye una medida integral de la riqueza en su base imponible. En muchos países, por ejemplo en España o Estados Unidos, diferentes áreas eximen del impuesto la riqueza personal distinta de las viviendas. Normalmente, las estructuras y el terreno sobre el que están construidas están sujetos al impuesto, pero los tipos efectivos difieren sustancialmente entre jurisdicciones.

Así, aunque continuamos describiendo el objeto de este apartado como el impuesto sobre la riqueza, ahora debería quedar claro que no existe tal cosa. La variedad de impuestos sobre la riqueza, así como sus diversas formas, son cruciales para evaluar los efectos económicos del sistema en su conjunto.






Los impuestos sobre la propiedad se aplican a una tasa determinada localmente sobre el valor tasado de la propiedad residencial y comercial, que es el valor asignado por la jurisdicción, normalmente el pueblo o la ciudad. Las localidades intentan hacer coincidir los valores tasados con los valores de mercado de las propiedades en la mayoría de los casos, pero esto es difícil cuando una propiedad no ha sido vendida recientemente. Como resultado, la tasa efectiva del impuesto sobre la propiedad (la proporción entre los pagos del impuesto sobre la propiedad y los valores de mercado de las propiedades) puede desviarse significativamente de la tasa legal. Las tasas efectivas se calculan multiplicando la tasa impositiva nominal por la proporción de tasación, el porcentaje del valor de la propiedad que está sujeto a tributación. Por ejemplo, una tasa impositiva nominal del 2% y una proporción de tasación del 50% darían lugar a una tasa impositiva efectiva del 1%.


Hay una serie de cuestiones económicas interesantes planteadas por la fiscalidad sobre la propiedad. Algunas de estas cuestiones fueron tratadas en otros temas, como la capitalización de los impuestos sobre la propiedad en los precios de las viviendas. Ahora, consideramos otras cuestiones importantes de la política del impuesto sobre la propiedad, como la incidencia del impuesto sobre la propiedad y los diferentes tipos de fiscalidad sobre la propiedad.














El tipo de gravamen suele aplicarse a l a base imponible reducida en un m ínimo exento, este principio se basa en que los patrimonios que deben soportar la carga son los mayores. En este sentido. el m ínimo exento es ya un primer paso para l a progresividad del tributo, que ya en la definición de l a base tenía en cuenta las circunstancias personales de cada contribuyente. El tipo a aplicar será único o progresivo según la opción del legislador. pero en cualquier caso debernos recordar que a l gravar un stock estaremos ante tipos m enores que los de los i mpuestos sobre l a rl'nta.




Los impuestos sobre las ganancias de capital dependiendo de la legislación se registran como impuestos sobre el capital o como impuestos corrientes sobre la renta, el patrimonio. Es decir, son impuestos que giran sobre un stock, pero solo de tipo temporal. Es decir, el patrimonio solo se grava cuando cambia de esfera jurídica. sea por el motivo que sea. 



Este tipo de tributos tienen dos elementos impor1antes de definición. La definición del patrimonio que se intercambia o traslada y la relación entre los contribuyentes.








\subsubsection{Impuesto sobre la propiedad}

\textcolor{red}{\textbf{Origen:} GRUBER}
Fiscalidad sobre la propiedad



¿Quién soporta el impuesto sobre la propiedad?

El impuesto sobre la propiedad es una fuente de mucho debate a nivel estatal y local, con una serie de estados imponiendo límites a la capacidad de las localidades para aumentar sus impuestos sobre la propiedad. Este debate refleja la opinión de que los impuestos sobre la propiedad son cargas costosas para los propietarios de viviendas de ingresos medios. Sin embargo, como sabemos [\authorfont{Ver Tema 3.B.8}], la incidencia de un impuesto no está determinada por quién paga la factura del impuesto. Entonces, ¿quién soporta realmente el impuesto sobre la propiedad?










Efectos sobre la incidencia y la eficiencia

La incidencia del impuesto sobre la propiedad ha sido una fuente de debate de larga duración entre los economistas de las finanzas públicas, como se resume en ZODROW (2001). Hay tres escuelas de pensamiento sobre la incidencia del impuesto sobre la propiedad. 

La cuestión de quién soporta en última instancia la carga del impuesto sobre la propiedad es controvertida. Analizamos tres perspectivas diferentes y después tratamos de conciliarlas.

Perspectiva tradicional: el impuesto sobre la propiedad como un impuesto especial sobre el consumo 
La perspectiva tradicional sostiene que el impuesto sobre la propiedad es un impuesto especial sobre el consumo que recae sobre el terreno y las estructuras. Por tanto, la incidencia del impuesto está determinada por las formas de las correspondientes curvas de oferta y demanda.

Según la perspectiva tradicional, por tanto, el impuesto sobre el terreno recae sobre los propietarios del terreno (o, más precisamente, los propietarios del terreno en el momento en que se aplica el impuesto); y el impuesto sobre las estructuras se traslada a los inquilinos.\footnote{%
    \textbf{Terreno}. Mientras la cantidad de terreno sea fija, su curva de oferta es perfectamente vertical, y los propietarios del terreno soportan toda la carga de un impuesto aplicado sobre él. Además, como ya analizamos ([\authorfont{Ver Tema 4.B.8}]), bajo ciertas circunstancias el impuesto se capitaliza en el valor del terreno. Los posibles compradores de terreno tienen en cuenta el hecho de que, si compran el terreno, también compran una corriente futura de obligaciones tributarias. En la medida en que el terreno no tenga una oferta fija, el análisis requiere una modificación. Por ejemplo, la oferta de terreno urbano puede ampliarse en los límites de las áreas urbanas que son adyacentes a tierras agrícolas. De manera similar, la oferta puede incrementarse si son factibles los rellenos de terrenos o la recuperación de terrenos baldíos. En tales casos, el impuesto sobre el terreno es soportado tanto por los propietarios como por los usuarios del terreno, en proporciones que dependen de las elasticidades de la demanda y de la oferta. Pero una curva de oferta vertical para el terreno suele ser una buena aproximación a la realidad.

\textbf{Estructuras.} Para comprender la perspectiva tradicional del impuesto sobre las estructuras, comenzamos considerando el mercado nacional de capital. El capital puede utilizarse para muchos fines: construcción de estructuras, equipos para la fabricación, proyectos del sector público como presas, etcétera. En cualquier momento dado, el capital tiene algún precio que lo raciona entre usos alternativos. Según la perspectiva tradicional, a largo plazo, la industria de la construcción puede obtener todo el capital que demande al precio de mercado. Así, la curva de oferta de estructuras es perfectamente horizontal. Bajo estas condiciones, antes del impuesto la oferta de estructuras es horizontal al precio vigente. De esta forma, al imponerse el impuesto el resultado es totalmente diferente al anterior. El precio recibido por los oferentes de estructuras será el mismo que el precio antes de que se impusiera el impuesto. Sin embargo, los demandantes de estructuras pagarán un precio que excede el precio original precisamente en la cantidad del impuesto. Por tanto, la carga se traslada enteramente a los inquilinos. Este resultado, por supuesto, se deriva del supuesto de una curva de oferta horizontal.
} En este sentido, la parte del impuesto sobre la propiedad correspondiente al terreno es soportada por las personas en proporción a la cantidad de renta de alquiler que reciben, y la parte del impuesto correspondiente a las estructuras es soportada por las personas en proporción a la cantidad de vivienda que consumen. De ello se deduce que el impacto de la parte del impuesto correspondiente al terreno sobre la progresividad depende de si la proporción de renta procedente de la propiedad del terreno tiende o no a aumentar con la renta. Existe un acuerdo bastante generalizado en que sí lo hace, por lo que esta parte del impuesto es progresiva. Del mismo modo, la progresividad del impuesto sobre las estructuras depende de manera crucial de si la proporción de renta dedicada a la vivienda aumenta o disminuye a medida que aumenta la renta. Si disminuye, entonces la parte del impuesto correspondiente a las estructuras es regresiva, y viceversa.

Una enorme cantidad de trabajo econométrico se ha centrado en cómo responden realmente los gastos en vivienda a los cambios en la renta. La capacidad de alcanzar un consenso se ha visto obstaculizada por el desacuerdo sobre qué concepto de renta utilizar. Algunos investigadores que utilizan la renta anual tienden a encontrar que la proporción de renta dedicada a la vivienda disminuye a medida que aumenta la renta, lo que sugiere que el impuesto es regresivo. Otros investigadores creen que alguna medida de la renta permanente es más relevante para comprender las decisiones sobre vivienda. Según esta perspectiva, el hecho de que la renta anual de una familia en un año determinado resulte ser mayor o menor que su renta permanente tiene poco impacto sobre el consumo de vivienda de ese año. Las decisiones sobre vivienda se toman en el contexto de las perspectivas a largo plazo de la familia, no de las variaciones anuales.

Por supuesto, si la renta permanente es la variable apropiada, entonces debe encontrarse alguna manera de estimarla. Un enfoque consiste en definir la renta permanente como la media de las rentas anuales de varios años. Resulta que los gastos en vivienda responden más a los cambios en la renta permanente que a los cambios en la renta anual. De hecho, aunque la evidencia es mixta, una conclusión razonable es que el consumo de vivienda es aproximadamente proporcional a la renta permanente. Por tanto, la parte del impuesto correspondiente a las estructuras probablemente no sea ni regresiva ni progresiva. Desafortunadamente, los análisis basados en la renta anual, que sugieren que el impuesto es regresivo, generalmente dominan las discusiones públicas sobre el impuesto.

\textbf{La nueva perspectiva: el impuesto sobre la propiedad como un impuesto sobre el capital} 
La perspectiva tradicional utiliza un marco estándar de equilibrio parcial. Sin embargo, la visión del \textbf{impuesto sobre el capital} de la incidencia del impuesto sobre la propiedad reconoce la naturaleza de equilibrio general de la incidencia impositiva. Para los académicos, como ya vimos, aunque el análisis de equilibrio parcial es a menudo útil, puede producir resultados engañosos para impuestos que son grandes en relación con la economía. Según la nueva perspectiva, es mejor pensar en el impuesto sobre la propiedad como un impuesto general sobre la riqueza en el que algunos activos son gravados por debajo del tipo medio y otros son gravados por encima. 

Supóngase por el momento que el impuesto sobre la propiedad puede aproximarse como un impuesto uniforme sobre todo el capital. Entonces, el impuesto sobre la propiedad es simplemente un impuesto general sobre el factor capital. Supóngase además que la oferta de capital a la economía es fija: cuando un factor tiene una oferta fija, soporta toda la carga de un impuesto general aplicado sobre él. Por tanto, el impuesto sobre la propiedad recae enteramente sobre los propietarios del capital. Y dado que la proporción de renta procedente del capital tiende a aumentar con la renta, un impuesto sobre el capital tiende a ser progresivo. Así, el impuesto sobre la propiedad es progresivo, ¡una conclusión que da completamente la vuelta a la perspectiva tradicional!

No obstante, en la mayoría de países, el impuesto sobre la propiedad decididamente no es un impuesto uniforme. Los tipos varían según el tipo de propiedad y la jurisdicción en la que se encuentra. Por tanto, el impuesto sobre la propiedad es un conjunto de impuestos especiales sobre el capital. Según la nueva perspectiva, el capital tiende a migrar desde las áreas en las que se enfrenta a un tipo impositivo alto hacia aquellas en las que el tipo es bajo. En un proceso que recuerda al modelo de HARBERGER ([\authorfont{Ver Tema 4.B.8}]), a medida que el capital migra hacia áreas con tipos impositivos bajos, su tasa de rendimiento antes de impuestos allí se reduce. Al mismo tiempo, la tasa de rendimiento antes de impuestos en las áreas con impuestos altos aumenta a medida que el capital se marcha. El proceso continúa hasta que las tasas de rendimiento después de impuestos son iguales en toda la economía. En general, a medida que el capital se desplaza, los rendimientos de otros factores de producción también cambian. El impacto sobre los otros factores depende en parte de su movilidad. El terreno, que es perfectamente inmóvil, no puede trasladar el impuesto. (Al menos en esta conclusión, las perspectivas nueva y antigua están de acuerdo). Del mismo modo, los tipos de capital menos móviles son los que tienen mayor probabilidad de soportar el impuesto. La incidencia final depende de cómo esté organizada la producción, de la estructura de la demanda de los consumidores y del grado en que los distintos factores sean móviles.\footnote{%
    Además, cabe estacar que hemos supuesto que la cantidad de capital disponible para la economía es fija. Sin embargo, a largo plazo, la oferta de capital puede depender del tipo impositivo. Si el impuesto sobre la propiedad disminuye la oferta de capital, la productividad del trabajo y, por tanto, el salario real disminuyen. Si el impuesto aumenta la acumulación de capital, ocurre exactamente lo contrario.
}

En definitiva, la nueva perspectiva muestra como el impuesto sobre la propiedad es un impuesto general sobre el capital en el que algunos tipos de capital son gravados a tipos superiores a la media y otros a tipos inferiores. El efecto general del impuesto es reducir el rendimiento del capital, lo que tiende a ser progresivo en su impacto sobre la distribución de la renta. Las diferencias en los tipos impositivos crean efectos de impuesto especial, que tienden a perjudicar a los factores inmóviles en las jurisdicciones con impuestos elevados. El proceso de ajuste puesto en marcha por estos efectos de impuesto especial es muy complicado, y no se sabe mucho sobre sus efectos sobre la progresividad. Tampoco puede decirse mucho acerca de la importancia de los efectos a largo plazo creados por los cambios en el tamaño del stock de capital. Si los efectos de impuesto especial y a largo plazo no contrarrestan con demasiada fuerza el efecto general, el impacto global del impuesto sobre la propiedad sería progresivo.


\textbf{El impuesto sobre la propiedad como una tarifa de usuario}

La discusión hasta ahora ha ignorado el hecho de que las comunidades utilizan los impuestos sobre la propiedad para comprar servicios públicos como educación y protección policial. En el modelo de TIEBOUT, el impuesto sobre la propiedad es simplemente el coste de comprar servicios públicos, y cada individuo compra exactamente la cantidad que desea. Así, el impuesto sobre la propiedad en realidad no es en absoluto un impuesto; se parece más a una tarifa de usuario por los servicios públicos. Esta perspectiva tiene dos implicaciones importantes:
\begin{la}
    \item \textbf{No es un gravamen}. La noción de la incidencia del impuesto sobre la propiedad carece de sentido porque el gravamen no es un impuesto en el sentido normal de la palabra; y,
    \item \textbf{Tarifa de usuario}. El impuesto sobre la propiedad no crea ninguna carga excesiva. Debido a que es simplemente la tarifa por los servicios públicos, no distorsiona el mercado de la vivienda más que el precio de cualquier otro bien.
\end{la}

El vínculo entre los impuestos sobre la propiedad y los servicios recibidos es a menudo tenue, por lo que no deberíamos tomar demasiado literalmente la noción del impuesto sobre la propiedad como una tarifa de usuario. No obstante, esta línea de razonamiento tiene implicaciones interesantes. Por ejemplo, si a las personas les importan los servicios públicos que reciben, esperamos que los efectos depresores de unos impuestos elevados sobre la propiedad sobre los valores de las viviendas sean contrarrestados por los servicios públicos financiados mediante estos impuestos. En un artículo clásico, OATES (1969) construyó un modelo econométrico de determinación del valor de la propiedad.\footnote{%
    Por supuesto, entre comunidades, los factores que influyen en los precios de las viviendas sí difieren, que incluyen tanto características físicas de las viviendas, como el número de habitaciones, y características de las propias comunidades, como la distancia a un centro urbano. Estos factores deben ser considerados al intentar separar los efectos de los impuestos sobre la propiedad y de los bienes públicos locales sobre los valores de la propiedad. OATES utilizó un análisis de regresión múltiple para hacerlo.
} Los resultados de su regresión sugieren que los aumentos del tipo del impuesto sobre la propiedad disminuyen los valores de las viviendas, mientras que los aumentos del gasto por alumno aumentan los valores de las viviendas, permaneciendo iguales el resto de cosas.  Estos resultados deben interpretarse con cautela. Por una parte, el gasto por alumno puede no ser una medida adecuada de los servicios públicos locales. Las localidades proporcionan muchos servicios públicos distintos de la educación, como protección policial, parques y bibliotecas. Además, incluso si la educación fuera el único bien público local, el gasto por alumno podría no ser una buena medida de la calidad educativa. Es posible, por ejemplo, que los gastos en una comunidad determinada sean elevados porque la comunidad tiene que pagar mucho a sus profesores, sus escuelas no son administradas eficientemente o sus estudiantes son particularmente difíciles de educar.

Con posterioridad al estudio de OATES, muchos otros investigadores han examinado las relaciones entre los valores de las propiedades, los impuestos sobre la propiedad y los bienes públicos locales utilizando datos de diferentes áreas geográficas y empleando diferentes conjuntos de variables explicativas. Aunque los resultados son algo mixtos, la conclusión general de OATES parece ser válida: los impuestos sobre la propiedad y el valor de los servicios públicos locales se capitalizan en los precios de las viviendas. Véase, por ejemplo, WEIMER y WOLKOFF (2001). Así, si dos comunidades tienen el mismo nivel de servicios públicos, pero la primera tiene impuestos más altos que la segunda (quizá porque su coste de proporcionar los servicios es mayor), esperamos que la primera tenga valores de las propiedades más bajos, permaneciendo iguales las demás cosas. 

De manera más general, estos resultados implican que, para comprender lo acomodados que están los miembros de una comunidad, no podemos considerar los tipos del impuesto sobre la propiedad de manera aislada. También deben considerarse los servicios gubernamentales y los valores de las propiedades.


Conciliación de las tres perspectivas 
Las tres perspectivas del impuesto sobre la propiedad no son mutuamente excluyentes. Cada una puede ser válida en diferentes contextos. Si, por ejemplo, queremos averiguar las consecuencias de eliminar todos los impuestos sobre la propiedad y sustituirlos por un impuesto nacional sobre las ventas, la nueva perspectiva es apropiada porque un cambio que afecta a todas las comunidades requiere un marco de equilibrio general. Por otro lado, si una comunidad determinada está considerando reducir su tipo del impuesto sobre la propiedad y compensar la pérdida de ingresos mediante un impuesto local sobre las ventas, la perspectiva tradicional ofrece la mayor comprensión. Esto se debe a que una única comunidad es tan pequeña en relación con la economía que su oferta de capital es esencialmente perfectamente horizontal. Finalmente, cuando los impuestos y los beneficios se modifican conjuntamente y las personas son suficientemente móviles como para poder escoger entre comunidades, la perspectiva de la tarifa de usuario es útil.

No obstante, la incidencia del impuesto sobre la propiedad depende en gran medida de cuál sea el modelo apropiado para analizar el impuesto. Bajo el modelo tradicional, el impuesto sobre la propiedad es soportado íntegramente por los residentes locales en proporción a sus gastos en vivienda y, por lo tanto, es proporcional o algo regresivo. Bajo la visión del impuesto sobre el capital, la parte común del impuesto sobre la propiedad es soportada en gran medida por los propietarios del capital, por lo que es mucho más progresivo. Bajo la visión del impuesto sobre los beneficios, no hay incidencia, ya que el impuesto sobre la propiedad es simplemente un precio pagado por un servicio recibido.

La evidencia existente no puede distinguir entre estos modelos, pero la teoría económica sugiere que la visión del impuesto sobre el capital tiene un mérito considerable, de modo que el impuesto puede ser más progresivo de lo que normalmente se percibe (ZODROW, 2001).
















¿Por qué odia tanto la gente el impuesto sobre la propiedad?

El 7 de junio de 1978, los votantes de California aprobaron una iniciativa estatal de limitación del impuesto sobre la propiedad conocida como Proposición 13. Sus disposiciones clave eran (1) establecer un límite del 1 por ciento al tipo del impuesto sobre la propiedad que cualquier localidad podía imponer, (2) limitar el valor tasado de la propiedad a su valor de 1975,¹³ y (3) prohibir a los gobiernos estatales y locales imponer cualesquiera impuestos adicionales sobre la propiedad sin la aprobación de una mayoría de dos tercios en una votación local. La Proposición 13 inició un movimiento para limitar el impuesto sobre la propiedad que sigue siendo fuerte hoy en día. Las encuestas de opinión pública indican regularmente que a la gente le desagrada el impuesto sobre la propiedad incluso más que el impuesto federal sobre la renta.

¿Por qué es tan impopular el impuesto sobre la propiedad? Se han propuesto varias explicaciones:

Debido a que las transacciones del mercado de la vivienda normalmente se producen con poca frecuencia, el impuesto sobre la propiedad debe aplicarse sobre un valor estimado. En la medida en que esta valoración se realice de manera incompetente (o corrupta), el impuesto se percibe como injusto.

El impuesto sobre la propiedad es altamente visible. Bajo los impuestos federales sobre la renta y sobre las nóminas, los pagos se retienen de las nóminas de los trabajadores, y el empleador envía los ingresos al gobierno. En contraste, el impuesto sobre la propiedad suele ser pagado directamente por el contribuyente. Además, los pagos vencen trimestral o anualmente, de modo que cada pago llega como un gran impacto. Es difícil saber hasta qué punto tomar en serio este argumento. Incluso los contribuyentes que de alguna manera no se dan cuenta del hecho de que los impuestos federales sobre la renta y sobre las nóminas son retenidos durante el año reciben cada abril un recordatorio claro de cuánto han pagado. Puede haber suficiente ira en ese único mes para durar todo un año.

El impuesto sobre la propiedad se percibe como regresivo. 
Esta percepción se debe en parte al continuo predominio de la perspectiva tradicional del impuesto sobre la propiedad en el debate público. Se ve reforzada por el hecho de que algunos propietarios, particularmente los ancianos, no tienen suficiente efectivo para realizar los pagos del impuesto sobre la propiedad y, por tanto, pueden verse obligados a vender sus viviendas. Algunos estados han respondido a este fenómeno introduciendo cortacircuitos que proporcionan beneficios a los contribuyentes (normalmente en forma de un reembolso en los impuestos estatales sobre la renta) que dependen del exceso de los pagos del impuesto sobre la propiedad residencial sobre alguna proporción especificada de la renta. Una solución mejor sería aplazar los pagos del impuesto hasta el momento en que se transfiera la propiedad.

A los contribuyentes pueden desagradarles otros impuestos tanto como el impuesto sobre la propiedad, pero se sienten impotentes para hacer algo respecto a los otros. Es relativamente fácil dirigir los ataques contra el impuesto sobre la propiedad, que se aplica localmente. Los residentes de Canaan, Nueva York, demostraron este hecho cuando demandaron al tasador local después de que sus facturas del impuesto sobre la propiedad se duplicaran con creces en un plazo de 10 años [Smith, 2005]. En contraste, organizar una campaña contra el impuesto federal sobre la renta es muy difícil, aunque solo sea porque sería necesaria una campaña nacional y, por tanto, implicaría grandes costes de coordinación.




A la luz de la hostilidad generalizada hacia el impuesto, es natural preguntarse si puede mejorarse. Una propuesta muy modesta consiste en mejorar los procedimientos de tasación. El uso de ordenadores y técnicas modernas de valoración puede hacer que las tasaciones sean más uniformes. En comparación con el sistema actual de tipos impositivos efectivos diferentes dentro de una jurisdicción, unos tipos impositivos uniformes probablemente mejorarían la eficiencia. Las cuestiones de equidad son más complicadas. Superficialmente, parece una violación de la equidad horizontal que dos personas con propiedades idénticas paguen impuestos diferentes sobre ellas. Sin embargo, el fenómeno de la capitalización exige que distingamos cuidadosamente entre los propietarios en el momento en que se aplica el impuesto y los propietarios actuales. Una propiedad con un tipo impositivo indebidamente alto se venderá por un precio más bajo, permaneciendo iguales las demás cosas. Así, un tipo impositivo alto no necesariamente hace que un individuo que compra la propiedad después de que se imponga el impuesto esté en peor situación. De hecho, igualar los ratios de tasación podría generar todo un nuevo conjunto de desigualdades horizontales.

Una reforma más ambiciosa del impuesto sobre la propiedad consistiría en convertirlo en un impuesto personal sobre el patrimonio neto, cuya base es la diferencia entre el valor de mercado de todos los activos y pasivos del contribuyente. Una ventaja de tal sistema frente a un impuesto sobre la propiedad es que, al permitir la deducción de los pasivos, proporciona un mejor índice de la capacidad de pago. Además, debido a que es un impuesto personal, pueden incorporarse exenciones al sistema y los tipos pueden variarse para alcanzar el grado deseado de progresividad.

Un impuesto personal sobre el patrimonio neto es una especie de impuesto general sobre la riqueza, y analizamos las cuestiones administrativas y económicas asociadas con la tributación de la riqueza en el Capítulo 21. En el contexto de la reforma del impuesto sobre la propiedad, es particularmente importante señalar que, debido a que los individuos pueden tener activos y pasivos en diferentes jurisdicciones, un impuesto sobre el patrimonio neto tendría indudablemente que ser administrado por el gobierno federal. Esto nos lleva a lo que muchas personas consideran que es la principal justificación del sistema actual de tributación de la propiedad. Cualesquiera que sean sus defectos, el impuesto sobre la propiedad puede ser administrado localmente sin ninguna ayuda de los gobiernos federal o estatales. Por tanto, proporciona al gobierno local una considerable autonomía fiscal. Según esta perspectiva, la eliminación del impuesto sobre la propiedad acabaría destruyendo la independencia económica de las unidades locales de gobierno.

La experiencia de California después de la Proposición 13 es coherente con esta noción. Debido a que la Proposición 13 limitó la capacidad de las comunidades para recaudar dinero mediante impuestos sobre la propiedad, esa medida aumentó la importancia de los ingresos estatales, y esto parece haber trasladado el poder sobre la política educativa desde las localidades al gobierno estatal. De manera similar, Cheung [2008] encontró que la Proposición 13 también condujo a un desplazamiento hacia las asociaciones de propietarios de viviendas, que son instituciones privadas que tienen autoridad para imponer gravámenes, proporcionar servicios públicos y hacer cumplir regulaciones a sus miembros. Así, el papel político del impuesto sobre la propiedad debe tomarse en serio en cualquier discusión sobre su reforma.




\begin{specialblock}{\textbf{Extra:} Otros impuestos directos}
    La definición de un impuesto como directo ya dijimos que era contenciosa. pero quedémonos con lo que hemos venido e>..plicando. Serán im puestos directos aquellos que gravan capacidades de renta o riqueza. periódicamente o puntualmente. En este sentido existen más impuestos directos.
    
    VI I I . J . Impuestos sobre el valor de los bienes inmuebles (IBI).
    Estos impuestos gravan los bienes inmuebles asociados a su valor. que puede ser de mercado o registral. Este tipo de tributos son muy comunes en muchos países (España. EE.UU., Grecia etc . . . ) y sirven para financiar generalmente necesidades locales, como los colegios en EE.UU. Dependiendo de la naturaleza jurídica (de origen napoleónico, sajón etc . . .) de la jurisdicción donde este el impuesto. puede basarse en una valoración de mercado de los inmuebles o una registra! basada en un Catastro.
    
    V I I I . U . I m puestos sobre e l valor de los bienes m uebles (IBM).
    La valoración de bienes muebles, suele tornar la forma de impuesto sobre la matriculación de carácter periódico y no debe confundirse con tasas que se pueden pagar asociadas a los bienes muebles. Este tipo de tributos, existen por su facil idad de gestión y fácil recaudación. en lo que parece ser una manifestación de capacidad económica. que cada día es más común. Este tipo de tributos pueden existir sobre multitud de bienes muebles o incluso sobre cuentas bancarias o títulos valores.
    
    V J I I . U I . Impuestos sobre el incremento del valor de los bienes (Plusvalía).
    Cuando se produce u n incremento del valor de los bienes anteriores, puede girarse un impuesto sobre dicho fenómeno. Que estaría muy en conflicto con la idea de ganancia patrimonial que expresamos en los impuestos sobre la renta. De hecho, la existencia de tributos así medidos puede considerarse doble imposición. En cualquier caso, este tipo de tributos consisten básicamente en extraer este concepto de ganancia patrimonial y no considerarlo una renta. sino una manifestación económica propia y que es digna de tributación. de hecho, algunos sistemas tributarios sólo gravan las ganancias patrimoniales de esta forma.
\end{specialblock}




\clearpage
\section{Imposición directa: Instantáneos}
\puntosclave{
    
}





\subsection{Imposición sobre el ahorro: Transmisiones patrimoniales} (Los dos: Hecho)
Un impuesto importante sobre el capital es el impuesto sobre las transmisiones patrimoniales. Los impuestos sobre las transmisiones patrimoniales pueden adoptar una extensa variedad de formas y definiciones legales (ITP, ISD, etc.) pero, en esencia, son todos el mismo tipo de impuesto en términos económicos: se grava la transferencia de un bien determinado, propiedad del individuo $x$, al individuo $y$. Como vemos, esta es la razón por la que debe ser tratada en el ámbito de la tributación instantánea: por definición, la transferencia unívoca de un bien entre individuos no puede ser realizada más de en una ocasión seguidamente.

El ejemplo por antonomasia es el Impuesto sobre Sucesiones y Donaciones. En la mayoría de países, entre ellos España, los impuestos sobre donaciones y sucesiones están unificados bajo una misma norma legal, de modo que el patrimonio total está constituido por los activos que se dejan al fallecimiento más el importe total de las donaciones realizadas a lo largo de la vida (por encima de determinadas exenciones).

Veintisiete de los treinta países de la OCDE disponían de alguna forma de tributación sobre las transmisiones patrimoniales en 2012. Estados Unidos obtiene una proporción relativamente elevada de sus ingresos públicos mediante estos impuestos: los impuestos sobre transmisiones representan, en promedio, únicamente el $0,4\%$ de los ingresos públicos de la OCDE, pero representan el $0,6\%$ de los ingresos públicos en Estados Unidos. Sin embargo, trece países de la OCDE aplican además un impuesto sobre el patrimonio (wealth tax) anual sobre el valor de los activos poseídos por sus ciudadanos. Este impuesto no existe en Estados Unidos. El país medio de la OCDE obtiene el $0,9\%$ de sus ingresos públicos mediante impuestos sobre el patrimonio y sobre las transmisiones patrimoniales, una proporción algo superior al $0,6\%$ correspondiente a Estados Unidos.

\subsubsection{¿Por qué gravar la riqueza? Argumentos a favor del impuesto sobre sucesiones}
Al analizar el impuesto sobre sucesiones, primero debemos plantearnos la siguiente cuestión: ¿por qué deberían los gobiernos gravar el stock de riqueza (ya sea anualmente o en el momento del fallecimiento), en lugar de gravar únicamente el flujo anual de renta (o el consumo)? 

Existen al menos tres argumentos a favor de gravar la riqueza.



Los opositores al impuesto sobre sucesiones, sin embargo, también mantienen opiniones sólidas y bien fundamentadas sobre esta cuestión. Señalan que gran parte de la renta gravada por el impuesto sobre sucesiones ya fue gravada cuando los titulares del patrimonio estaban vivos, por lo que supone una doble imposición sobre su renta. Esta doble imposición puede llevar a las personas más ricas a obtener menos ingresos y ahorrar menos. Los opositores al impuesto sobre sucesiones también lo consideran un impuesto sobre la muerte que impone una penalización gubernamental por el propio hecho de morir, un insulto añadido al daño definitivo. Como afirma Grover Norquist, presidente de Americans for Tax Reform: Los ricos han pagado todos los malditos impuestos que se hayan inventado jamás. ¿Por qué deberían pagar impuestos solo porque fallecen?

Los argumentos de los opositores parecen estar imponiéndose en los últimos años: la ley fiscal de 2001 aumentó sustancialmente los umbrales a partir de los cuales los patrimonios están sujetos al impuesto, elevando el umbral mínimo hasta patrimonios de 5,43 millones de dólares en 2015.


\paragraph{Desincentivos: ¿Sobre el receptor o el dondante?}
Un tercer argumento que ha sido planteado con frecuencia por los defensores del impuesto sobre sucesiones es la afirmación de que permitir que los hijos de familias ricas hereden toda la riqueza de sus padres les priva de toda motivación para trabajar con esfuerzo y alcanzar el éxito por sí mismos. Como afirmó Andrew Carnegie en 1891: \textit{El padre que deja a su hijo una enorme fortuna generalmente adormece los talentos y las energías de su hijo, y lo tienta a llevar una vida menos útil y menos digna de la que, de otro modo, habría llevado} (HOLTZ-EAKIN et al., 1993).\footnote{%
    De hecho, HOLTZ-EAKIN et al. encontraron que, cuanto mayor es la herencia que recibe un individuo, mayor es la probabilidad de que abandone la población activa.
}

No obstante, el problema de los incentivos es más complicado de lo que sugiere Carnegie, porque debemos tener en cuenta el comportamiento del donante, no solo el del receptor. Considérese a Samuel, un individuo que está motivado para trabajar duro durante su vida con el fin de dejar una gran sucesión a sus hijas. La presencia de un impuesto sobre sucesiones podría desalentar el esfuerzo laboral de Samuel. Por otro lado, con un impuesto sobre sucesiones, tiene que acumularse una mayor cantidad de riqueza para dejar un legado determinado después de impuestos. Así, la presencia de un impuesto sobre sucesiones podría inducir a Samuel a trabajar más duro para mantener el valor neto de su sucesión. En consecuencia, el impacto de un impuesto sobre sucesiones sobre el esfuerzo laboral de un donante es lógicamente indeterminado.\footnote{%
    Del mismo modo, no podemos predecir cómo afecta un impuesto sobre sucesiones al comportamiento de ahorro del donante. Es fácil describir escenarios en los que ahorra menos y escenarios en los que ahorra más. Otro argumento que se suele plantear es que el impuesto sobre sucesiones también puede afectar a la forma en la que se producen las transferencias, no solo a la cantidad tranferida. Por ejemplo, un impuesto sobre los legados de capital físico crea incentivos para transmitir riqueza en forma de capital humano. Así, en lugar de dar a cada hija $80.000€$ en acciones y bonos, Samuel podría gastar $80.000€$ en la educación universitaria de cada una de ellas. Por tanto, un impuesto sobre sucesiones podría conducir a una sobreinversión en capital humano.
} En este sentido, KOPCZUK y SLEMROD (2001) examinaron datos estadounidenses de declaraciones del impuesto sobre sucesiones presentadas entre 1916 y 1996 para evaluar el efecto de los tipos del impuesto sobre sucesiones sobre las sucesiones declaradas. Encontraron una relación negativa entre la magnitud del tipo impositivo que prevalecía diez años antes de la muerte y el tamaño de la sucesión, lo que sugiere que los aumentos del tipo impositivo reducen la acumulación de riqueza. 

Como vemos, incluso si Carnegie tuviera razón en que el impuesto sobre sucesiones lleva a los herederos potenciales a trabajar más, también podría generar incentivos para que los donantes trabajaran menos. La teoría por sí sola no nos dice qué tendencia domina. 

\paragraph{Beneficios inmateriales: ¿compensa esto el impuesto?}
Algunos sostienen que el gobierno protege los derechos de propiedad y supervisa la transferencia de la propiedad del causante a sus herederos. Como compensación por proporcionar estos servicios, el Estado tiene derecho a una parte de la sucesión. Quienes se oponen al impuesto sobre sucesiones creen que la prestación de tales servicios es un derecho fundamental por el que no se tiene que pagar. Como lo expresó la actriz Whoopi Goldberg: \textit{No quiero que me cobren impuestos solo porque morí. Simplemente no creo que esté bien}. Además, parece arbitrario seleccionar las transferencias de propiedad como objetos especiales de tributación. Si Ana gasta $10.000$ euros en un viaje a Estados Unidos, Carmen gasta $10.000€$ en la educación universitaria de su hija y Samuel deja $10.000€$ a su hijo, ¿por qué debería Samuel enfrentarse a un impuesto especial?

En este sentido, los defensores sostienen que, en última instancia, toda propiedad pertenece a la sociedad en su conjunto. Durante la vida de un individuo, la sociedad le permite disponer como desee de la propiedad que ha conseguido acumular. Pero, al morir, la propiedad revierte a la sociedad, que puede disponer de ella a voluntad. Según esta perspectiva, aunque las personas puedan tener derecho a lo que ganan, sus descendientes no tienen ningún derecho ético convincente sobre ello. No obstante, detrás de tales afirmaciones se encuentran muchos juicios de valor controvertidos. Los oponentes creen que es fundamentalmente incorrecto argumentar que una persona posee riqueza únicamente por voluntad de la sociedad, o que la sociedad tenga alguna vez algún derecho válido sobre la riqueza personal.


Otro problema del impuesto sobre sucesiones es similar al planteado anteriormente respecto a la tributación de las ganancias de capital sobre la base del devengo: para poder pagar el impuesto, puede ser necesario vender el activo. Esto podría constituir un problema importante, por ejemplo, en el caso de las explotaciones agrícolas familiares, donde los hijos desean continuar con la actividad agrícola tras la muerte de sus padres, pero la única forma de hacer frente al impuesto sobre sucesiones consiste en vender la explotación.












Con diferencia, el impuesto sobre la riqueza más importante en Estados Unidos es el impuesto sobre la propiedad, que es particularmente crucial para las operaciones de los gobiernos locales. En consecuencia, posponemos nuestra discusión del impuesto sobre la propiedad hasta el Capítulo 22, en el cual discutimos las unidades subnacionales de gobierno, y en su lugar nos centramos aquí en los impuestos sobre regalos y legados.

Impuestos sobre sucesiones y donaciones

Estos impuestos se aplican a intervalos irregulares cuando tienen lugar determinados acontecimientos: el impuesto sobre sucesiones a la muerte del poseedor de la riqueza (causante), y el impuesto sobre donaciones cuando la propiedad se transfiere entre personas vivas (inter vivos). Los impuestos sobre sucesiones y donaciones representan solo alrededor del $1\%$ por ciento de los ingresos fiscales del gobierno federal de Estados Unidos (Economic Report of the President, 2009). La incidencia legal del impuesto federal no afecta a las vidas de la mayoría de los ciudadanos. Menos del $1\%$ por ciento de todos los causantes tienen sucesiones que están sujetas al impuesto. Algunos han sugerido que el papel de los impuestos sobre sucesiones y donaciones debería ampliarse. Sin embargo, la voluntad política parece ser más bien la contraria. Analicemos los argumentos a favor y en contra.



Disposiciones

La tributación de las donaciones y la tributación de las sucesiones están inextricablemente vinculadas. Supóngase que las sucesiones están gravadas y las donaciones no. Si Samuel desea transmitir su riqueza a sus hijas y sabe que será gravada a su muerte, entonces puede evitar el impuesto realizando la transferencia como una donación inter vivos. Surgirían oportunidades similares si existiera un impuesto sobre donaciones pero no un impuesto sobre sucesiones. 

El impuesto unificado sobre transferencias es similar en su estructura básica al impuesto sobre la renta personal. Después de calcular la sucesión bruta, se restan diversas deducciones y exenciones, quedando la sucesión imponible. La deuda tributaria se determina aplicando una escala progresiva de tipos a la sucesión imponible.

Cálculo de la base imponible 
La sucesión bruta consiste en toda la propiedad poseída por el causante en el momento de la muerte, incluidos bienes inmuebles, acciones, bonos y pólizas de seguro. También incluye las donaciones realizadas durante la vida del causante. Para determinar la sucesión imponible, se permiten deducciones por gastos funerarios, costes de liquidación de la sucesión (honorarios de abogados) y cualesquiera deudas pendientes de la sucesión. Las donaciones a organizaciones benéficas son deducibles sin límite.

Están disponibles las siguientes deducciones:

• En 2009, a cada sucesión se le permitía una exención vitalicia de 3,5 millones de dólares. No se aplica ningún impuesto federal sobre sucesiones a las sucesiones que sean inferiores a la exención vitalicia. La exención no está indexada a la inflación.

• Todas las transferencias cualificadas a los cónyuges —mediante donación o legado— son deducibles al determinar la base imponible. Así, la sucesión de una multimillonaria que deja 2 millones de dólares a sus hijos y el resto a su marido no soporta ninguna deuda tributaria. Debido a la deducción conyugal, la mayoría de las parejas casadas no pagan ningún impuesto sobre sucesiones hasta que ambos cónyuges han muerto.

• Cada individuo tenía derecho a una exclusión anual de donaciones de 13.000 dólares por receptor. (El receptor no necesita ser un familiar). Considérese una familia con tres hijos. Cada año, Mamá puede dar 13.000 dólares a cada hijo, al igual que Papá. Juntos, por tanto, la pareja puede dar anualmente a sus tres hijos 78.000 dólares libres de impuestos. Curiosamente, existen algunas pruebas de que las personas ricas no aprovechan plenamente las ventajas fiscales de distribuir la riqueza antes de la muerte. ¿Por qué? Hay una historia sobre un hombre rico que dio a cada uno de sus hijos 1 millón de dólares cuando alcanzaron la edad de 21 años. Cuando se le preguntó por qué lo hizo, el millonario explicó que quería que sus hijos pudieran decirle que «se fuera al infierno» —que tuvieran una independencia financiera total—. Parece que la mayoría de las personas preferirían que sus hijos no pudieran decirles que se fueran al infierno. Por tanto, estas personas mantienen el control de su riqueza durante el mayor tiempo posible, incluso a costa de una deuda tributaria mayor de lo necesario.

Estructura de tipos 
La base imponible está sujeta a tipos impositivos marginales crecientes. Para 2009, el tipo máximo era del 45 por ciento. Es difícil decir si este tipo es eficiente o no. Como de costumbre, la respuesta depende de la capacidad de respuesta del comportamiento a los cambios en el tipo impositivo. Pero, como se indicó anteriormente, se sabe poco sobre cómo las decisiones económicas se ven afectadas por los impuestos sobre sucesiones y donaciones.



Problemas especiales 

Surgen varias dificultades en la administración de un impuesto sobre sucesiones y donaciones.

Propiedad poseída conjuntamente 
Supóngase que un marido y una mujer poseen conjuntamente una propiedad. A efectos de la tributación de las sucesiones, ¿debería considerarse esto una sucesión o dos? Como ya analizamos antes en relación con la unidad de tributación, no es necesario volver a plantear los problemas aquí. Según la legislación federal actual, la mitad del valor de la propiedad poseída conjuntamente se incluye ahora en la sucesión bruta del primer cónyuge que muere, independientemente de la medida relativa en que los cónyuges contribuyeron a la acumulación de la propiedad.¹¹

Empresas de propiedad restringida 
Supóngase que Samuel quiere legar su empresa, que es el único activo que posee, a sus hijas. Debido a que no hay efectivo en la sucesión de Samuel, las hijas pueden tener que vender la empresa para pagar el impuesto sobre sucesiones adeudado. Para reducir la probabilidad de que se produzca tal acontecimiento, la ley puede permitir que los impuestos sobre sucesiones correspondientes a empresas de propiedad restringida sean prorrateados durante un periodo temporal definido, a tipos de interés favorables. Además, al calcular la sucesión bruta, las explotaciones agrícolas familiares y las empresas familiares que cumplen los requisitos se valoran por debajo de su valor justo de mercado. Tales disposiciones reflejarían un juicio de valor según el cual es socialmente deseable per se que la misma familia controle una determinada empresa durante varias generaciones y, también, el poder político de los propietarios de pequeñas empresas.




Reforma de los impuestos sobre sucesiones y donaciones

Para quienes desean ampliar el papel de los impuestos sobre sucesiones y donaciones, el enfoque más directo sería reducir la exención vitalicia. Sin embargo, si el impuesto sobre sucesiones ha de desempeñar alguna vez un papel importante en el sistema de ingresos, deben idearse métodos para hacer frente a la elusión mediante fideicomisos y otros instrumentos similares.

Algunos teóricos de la tributación proponen integrar el sistema de impuestos sobre sucesiones y donaciones en el impuesto sobre la renta personal. Las donaciones y las herencias serían gravadas como renta de los receptores. Como se señaló anteriormente, tales ingresos son renta y, de acuerdo con la definición de renta de HAIG-SIMONS, deberían por tanto incluirse en la renta bruta ajustada. Para tener en cuenta el hecho de que la renta en esta forma tiende a ser irregular, tendría que idearse alguna forma de promediación.

Existe, sin embargo, resistencia popular a gravar las donaciones y las herencias como renta ordinaria. Un método diferente de trasladar el foco de la tributación de sucesiones y donaciones del donante al receptor es un impuesto sobre adquisiciones, bajo el cual cada individuo es gravado sobre el total de adquisiciones a lo largo de su vida procedentes de herencias y donaciones. La escala de tipos podría hacerse progresiva e incluir una exención, si así se deseara. El atractivo de tal sistema es que relaciona las obligaciones tributarias con la capacidad de pago del receptor en lugar de con la sucesión. Surgirían dificultades administrativas por la necesidad de que los contribuyentes mantuvieran registros de todas las donaciones y sucesiones de tamaño considerable. Pero si alguna vez se decide gravar las transferencias de riqueza de manera más agresiva, un impuesto sobre adquisiciones merece una consideración seria. Por otro lado, para quienes se oponen a la tributación de las transferencias de riqueza por razones filosóficas o económicas, la mejor reforma del impuesto es abolirlo.

































\clearpage
\section{Conclusión} 


\subsection{Recapitulación}


\subsection{Extensiones y relación con otras partes del Temario}



\subsection{Opinión}

\subsubsection{Idea final (Salida o cierra)}

\clearpage
\pagenumbering{gobble}
\MyBibliography

- SCHANZ, G. von (1896): "Der Einkommensbegriff und die Einkommensteuergesetze", FinanzArchiv, vol. 13.
- HAIG, R. M. (1921): "The Concept of Income: Economic and Legal Aspects", en The Federal Income Tax, Columbia University Press.
- SIMONS, H. C. (1938): Personal Income Taxation: The Definition of Income as a Problem of Fiscal Policy, University of Chicago Press.
- WAGNER, A.: Finanzwissenschaft, tradición hacendística alemana de la renta-producto (referencia genérica de escuela, ya usada en el Tema 16).
- KALDOR, N. (1955): An Expenditure Tax, George Allen & Unwin, Londres.
- MEADE, J. E. (dir.) (1978): The Structure and Reform of Direct Taxation (Informe Meade), Institute for Fiscal Studies / George Allen & Unwin.
- ANDREWS, W. D. (1974): "A Consumption-Type or Cash Flow Personal Income Tax", Harvard Law Review, vol. 87, núm. 6.
- VICKREY, W. S. (1947): Agenda for Progressive Taxation, The Ronald Press Company, Nueva York (tesis doctoral, Columbia University).
- BOSKIN, M. (2000): "From Edgeworth to Vickrey to Mirrlees: The Vickrey Lecture", Atlantic Economic Journal, vol. 28, núm. 1.
- BUCHANAN, N. H. (2006): "Should We Adopt William Vickrey's Cumulative Averaging Income Tax System? Progressivity and Simplicity in Tax Reform".
- Tokyo Foundation for Policy Research (2021): "A Proposal for Lifetime Income Taxation: The Vickrey Mechanism Revisited".
- SAMUELSON, P. A. (1956): "Social Indifference Curves", Quarterly Journal of Economics, vol. 70, núm. 1.
- BECKER, G. S. (1981): A Treatise on the Family, Harvard University Press.
- BERGSTROM, T. C. (1989): "A Fresh Look at the Rotten Kid Theorem — and Other Household Mysteries", Journal of Political Economy, vol. 97, núm. 5.
- CHIAPPORI, P.-A. (1988): "Rational Household Labor Supply", Econometrica, vol. 56, núm. 1; y (1992): "Collective Labor Supply and Welfare", Journal of Political Economy, vol. 100, núm. 3.
- McELROY, M. B. y HORNEY, M. J. (1981): "Nash-Bargained Household Decisions: Toward a Generalization of the Theory of Demand", International Economic Review, vol. 22, núm. 2.
- MANSER, M. y BROWN, M. (1980): "Marriage and Household Decision-Making: A Bargaining Analysis", International Economic Review, vol. 21, núm. 1.
- BROWNING, M. y CHIAPPORI, P.-A. (1998): "Efficient Intra-Household Allocations: A General Characterization and Empirical Tests", Econometrica, vol. 66, núm. 6.
- BICK, A. y FUCHS-SCHÜNDELN, N. (2018): "Taxation and Labour Supply of Married Couples across Countries: A Macroeconomic Analysis", Review of Economic Studies, vol. 85, núm. 3.
- MUSGRAVE, P. B. (1963): Taxation of Foreign Investment Income: An Economic Analysis, Johns Hopkins University Press
- DESAI, M. A. y HINES, J. R. (2003): "Evaluating International Tax Reform", National Tax Journal, vol. 56, núm. 3.
- WEISBACH, D. A. (2014): "The Use of Neutralities in International Tax Policy", Oxford University Centre for Business Taxation, Working Paper 14/14.
- TILLY, C. (1990): Coercion, Capital, and European States, AD 990-1990, Basil Blackwell, Oxford.
- BESLEY, T. y PERSSON, T. (2009): "The Origins of State Capacity: Property Rights, Taxation, and Politics", American Economic Review, vol. 99, núm. 4.
- Pollock v. Farmers' Loan \& Trust Co., 157 U.S. 429 (1895).
- Decimosexta Enmienda a la Constitución de los Estados Unidos (1913).
- Revenue Act de 1913 (Estados Unidos).
- Income Tax Act 1842 (Reino Unido, Sir Robert Peel).
- SMITH, A. (1776): An Inquiry into the Nature and Causes of the Wealth of Nations, Libro V (los cuatro cánones de la imposición).
- WICKSELL, K. (1896): Finanztheoretische Untersuchungen, Jena.
- LINDAHL, E. (1919): Die Gerechtigkeit der Besteuerung, Lund.
- MILL, J. S. (1848): Principles of Political Economy, Libro V, cap. II (el sacrificio igual).
- EDGEWORTH, F. Y. (1897): "The Pure Theory of Taxation", The Economic Journal, vol. 7, núm. 25.
- ROBBINS, L. (1932): An Essay on the Nature and Significance of Economic Science, Macmillan, Londres.
- BRENNAN, G. y BUCHANAN, J. M. (1980): The Power to Tax: Analytical Foundations of a Fiscal Constitution, Cambridge University Press.
- SURREY, S. S. (1973): Pathways to Tax Reform: The Concept of Tax Expenditures, Harvard University Press.
- BITTKER, B. I. (1969): "The Tax Expenditure Budget: A Reply to Professors Surrey and Hellmuth", National Tax Journal, vol. XXII, núm. 4.
- TITMUSS, R. M. (años 50): concepto de fiscal welfare, tradición de política social británica.
- MEHROTRA, A. K. y ZELENAK, L. (2022): reseña de la trayectoria intelectual de Stanley Surrey sobre gastos fiscales y equidad horizontal, Tax Notes Federal.
- HALL, R. E. y RABUSHKA, A. (1985): The Flat Tax, Hoover Institution Press.
- RAMSEY, F. P. (1927): "A Contribution to the Theory of Taxation", The Economic Journal, vol. 37, núm. 145.
- CHAMLEY, C. (1986): "Optimal Taxation of Capital Income in General Equilibrium with Infinite Lives", Econometrica, vol. 54, núm. 3.
- JUDD, K. L. (1985): "Redistributive Taxation in a Simple Perfect Foresight Model", Journal of Public Economics, vol. 28, núm. 1.
- STRAUB, L. y WERNING, I. (2020): "Positive Long-Run Capital Taxation: Chamley-Judd Revisited", American Economic Review, vol. 110, núm. 1.
- SAEZ, E. (2010): "Do Taxpayers Bunch at Kink Points?", American Economic Journal: Economic Policy, vol. 2, núm. 3.
- CHETTY, R., LOONEY, A. y KROFT, K. (2009): "Salience and Taxation: Theory and Evidence", American Economic Review, vol. 99, núm. 4.
- COASE, R. H. (1937): "The Nature of the Firm", Economica, vol. 4, núm. 16.
- BERLE, A. A. y MEANS, G. C. (1932): The Modern Corporation and Private Property, Macmillan, Nueva York.
- OCDE (2015): Neutralising the Effects of Hybrid Mismatch Arrangements, Action 2 - 2015 Final Report, Proyecto BEPS OCDE/G20.
- OECD Tax Database, Table I.7, Top Statutory Personal Income Tax Rates (serie histórica 1975-2024).
- Tax Policy Center / Tax Foundation: series históricas de tipos marginales OCDE, 1975-2013.
- Tax Cuts and Jobs Act de 2017 (Estados Unidos), régimen GILTI y exención parcial de dividendos de fuente extranjera.
- Reforma sueca de imposición individual de 1971 (referencia comparada de transición agregación-desagregación familiar).

% ANEXOS
\clearpage
\anexo{Problema previo: Debate doctrinal sobre el concepto de renta}\label{anx:teorias_renta}
Antes de presentar SHS como solución, hay que presentar el problema que resuelve. La doctrina hacendística clásica distingue dos grandes tradiciones: la teoría de la fuente o renta-producto (de origen alemán decimonónico, Wagner/Hermann), que solo considera renta el flujo periódico y regular derivado de una fuente productiva estable (excluye, por ejemplo, las plusvalías esporádicas o las herencias, por no ser "fruto" recurrente de una fuente), y la teoría del incremento patrimonial neto o renta-entrada (accretion theory), que considera renta cualquier incremento del poder económico del sujeto, provenga de donde provenga y tenga o no carácter periódico. [Probable] Esta contraposición no es un excurso erudito prescindible: es la que explica por qué el IRPF español de 1978 pudo llamarse "sintético" —adopta decididamente la segunda tradición— y por qué, pese a ello, conserva rasgos de la primera (ej. distinción entre rendimientos y ganancias patrimoniales, que es un resto conceptual de la vieja distinción fuente/incremento). Arrancar aquí te permite justificar por qué SHS, y no otra definición, es el benchmark adecuado: es la culminación técnica de la teoría de la renta-entrada.









\end{document}